\documentclass[11pt]{article}
\usepackage[a4paper,margin=1.1in]{geometry}
\usepackage[T1]{fontenc}
\usepackage[utf8]{inputenc}
\usepackage[ngerman]{babel}

\usepackage{lmodern}
\usepackage{amsmath,amssymb,amsthm,mathrsfs}
\usepackage{microtype}
\usepackage{fancyhdr}
\usepackage{enumitem}
\usepackage{longtable}
\usepackage{array}
\setlength{\parindent}{1.5em}
\setlength{\parskip}{0.18em}
\pagestyle{fancy}
\fancyhf{}
\cfoot{\thepage}
\newcommand{\Om}{\Omega}
\newcommand{\Omm}{\Omega_m}
\newcommand{\Hm}{H_m}
\newcommand{\Norm}{\mathrm{N}}
\newcommand{\Normone}{\mathrm{N}_1}
\newcommand{\frp}{\mathfrak p}
\newcommand{\frP}{\mathfrak P}
\newcommand{\frG}{\mathfrak G}
\newcommand{\frA}{\mathfrak A}
\newcommand{\frB}{\mathfrak B}
\newcommand{\Gg}{\mathfrak G}
\newcommand{\Aa}{\mathfrak A}
\newcommand{\Bb}{\mathfrak B}
\newcommand{\pp}{\mathfrak p}
\newcommand{\eps}{\varepsilon}
\newcommand{\Na}{N_a}
\newcommand{\srcpage}[1]{\par\medskip\noindent\textsf{[source p.~#1]}\par\smallskip}

\lhead{Heinrich Weber, Lehrbuch der Algebra, Bd. II}
\rhead{literal repair cumulative}
\begin{document}
\section*{\S\,176. Der in $\Omega_m$ enthaltene reelle Körper $H_m$.}

Jede Zahl des Körpers $\Omega_m$ geht durch die Substitution $(r,r^{-1})$ in die conjugirt imaginäre Zahl über, die auch durch die Vertauschung $(i,-i)$ erhalten wird. Das Product zweier solcher conjugirt imaginärer Zahlen ist das Quadrat des absoluten Werthes einer jeden von ihnen.

Eine Zahl in $\Omega_m$ ist reell, wenn sie durch die Vertauschung $(r,r^{-1})$ ungeändert bleibt, und nur unter dieser Voraussetzung. Jede solche Zahl lässt sich rational durch die zweigliedrige Periode $r+r^{-1}$ ausdrücken, und gehört also einem reellen Körper vom Grade $\frac12\varphi(m)$ an, der ein Theiler von $\Omega_m$ ist. Diesen wollen wir den in $\Omega_m$ enthaltenen reellen Körper nennen und mit $H_m$ bezeichnen (obwohl auch noch andere reelle Körper, nämlich alle Theiler von $H_m$, in $\Omega_m$ enthalten sind).

Die $\frac12\varphi(m)=\frac12\mu$ Zahlen
\begin{equation}
1,\quad r+r^{-1},\quad r^2+r^{-2},\ldots,\quad r^{\frac12\mu-1}+r^{-\frac12\mu+1}
\tag{1}
\end{equation}
bilden eine Basis der ganzen Zahlen des Körpers $H_m$. Denn nach \S\,170, (7) ist
\begin{equation}
\begin{aligned}
\omega={}&x_0+x_1r+x_2r^2+\cdots+x_{\frac12\mu-1}r^{\frac12\mu-1}+x_{\frac12\mu}r^{\frac12\mu} \\
&+x_{-1}r^{-1}+x_{-2}r^{-2}+\cdots+x_{-\frac12\mu+1}r^{-\frac12\mu+1}
\end{aligned}
\tag{2}
\end{equation}
dann und nur dann eine ganze Zahl des Körpers $\Omega_m$, wenn die Coefficienten $x_0,x_1,\ldots$ ganze rationale Zahlen sind.

Die Bedingung dafür, dass diese Zahl dem Körper $H_m$ angehört, erhält man, wenn man die Substitution $(r,r^{-1})$ ausführt und die Differenz $=0$ setzt. Dividirt man noch durch $r-r^{-1}$, so ergiebt sich so:
\begin{align*}
0={}&(x_1-x_{-1})+(x_2-x_{-2})(r+r^{-1})+\cdots \\
&+(x_{\frac12\mu-1}-x_{-\frac12\mu+1})\frac{r^{\frac12\mu-1}-r^{-\frac12\mu+1}}{r-r^{-1}}
+x_{\frac12\mu}\frac{r^{\frac12\mu}-r^{-\frac12\mu}}{r-r^{-1}}.
\end{align*}
Die Divisionen lassen sich hier ausführen, und wenn man dann mit $r^{\frac12\mu-1}$ multiplicirt, so ergiebt sich für $r$ eine rationale Gleichung von niedrigerem als $\mu$tem Grade, die nur dann befriedigt sein kann, wenn alle ihre Coefficienten verschwinden. Dies führt zu den Gleichungen
\[
x_{\frac12\mu}=0,\qquad x_{\frac12\mu-1}=x_{-\frac12\mu+1},\ldots,\qquad x_1=x_{-1},
\]
und es ist also jede ganze Zahl $\omega$ des Körpers $H_m$ in der Form enthalten:
\[
\omega=x_0+x_1(r+r^{-1})+x_2(r^2+r^{-2})+\cdots+x_{\frac12\mu-1}(r^{\frac12\mu-1}+r^{-\frac12\mu+1}).
\]
Statt der Basis (1) kann man auch, wenn man
\[
\rho=r+r^{-1}
\]
setzt, die Potenzen von $\rho$:
\begin{equation}
1,\quad \rho,\quad \rho^2,\ldots,\quad \rho^{\frac12\mu-1}
\tag{3}
\end{equation}
als Basis der ganzen Zahlen von $H_m$ wählen. Denn die Grössen (3) lassen sich ganzzahlig durch die (1) ausdrücken und umgekehrt.

Die Zahl $\rho$ ist die Wurzel einer irreduciblen Gleichung vom Grade $\frac12\mu$, die wir mit $\psi(\rho)=0$ bezeichnen wollen. Ist $f(r)=0$ die Gleichung für $r$ (\S\,169), so besteht die identische Relation
\begin{equation}
f(x)=x^{\frac12\mu}\psi\left(x+\frac1x\right),
\tag{4}
\end{equation}
woraus durch Bildung der Ableitung
\begin{equation}
f'(r)=r^{\frac12\mu}\psi'(\rho)(1-r^{-2}).
\tag{5}
\end{equation}
Hieraus können wir die Grundzahl $\Delta_1$ des Körpers $H_m$ bilden, die, da $H_m$ nur reelle Zahlen enthält, positiv sein muss. Diese Grundzahl ist, wenn wir die Norm in Bezug auf $H_m$ mit $\Normone$ bezeichnen:
\[
\Delta_1=\pm\Normone\psi'(\rho).
\]
Ist aber $\Norm$ die in Bezug auf $\Omega_m$ gebildete Norm, so ist
\[
\Norm\psi'(\rho)=[\Normone\psi'(\rho)]^2=\Delta_1^2,
\]
und demnach ergiebt die Formel (5) mit Rücksicht auf \S\,169:
\[
\Delta=\Norm(1-r^{-2})\Delta_1^2.
\]
Nach \S\,169, (7) und (10) ist:
\begin{equation}
\begin{aligned}
\Norm(1-r^{-2})&=q &&\text{bei ungeradem }q,\\
&=4 &&\text{für }q=2,\\
\pm\Delta&=q\Delta_1^2 &&\text{bei ungeradem }q,\\
&=4\Delta_1^2 &&\text{für }q=2.
\end{aligned}
\tag{6}
\end{equation}
Setzt man darin
\[
\Delta=\pm q^{q^{\kappa-1}[\kappa(q-1)-1]},
\]
so folgt
\begin{equation}
\begin{aligned}
\Delta_1&=q^{\,q^{\kappa-1}\frac{q-1}{2}-\frac{q^{\kappa-1}+1}{2}} &&\text{bei ungeradem }q,\\
&=2^{(\kappa-1)2^{\kappa-2}-1} &&\text{für }q=2,
\end{aligned}
\tag{7}
\end{equation}
so dass $\Delta_1$ immer eine Potenz von $q$ ist.

\section*{\S\,177. Die Primideale im Körper $H_m$.}

Es ist jetzt die Frage zu untersuchen, in welcher Beziehung die Primideale des Körpers $H_m$ zu denen des Körpers $\Omega_m$ stehen.

Wenn $\frP$ irgend ein Primideal in $H_m$ vom Grade $f_1$ bedeutet, so muss $\frP$ Theiler einer natürlichen Primzahl $p$ sein, und $\frP$ kann also (nach \S\,171) nur dann durch die zweite oder eine höhere Potenz eines Primfactors $\frp$ in $\Omega_m$ theilbar sein, wenn $p=q$ ist.

Ist aber $\frP$ ein Theiler von $q$, so muss es eine Potenz von $\sigma=1-r$ sein. Wir setzen etwa
\[
\frP=\sigma^\lambda.
\]
Nehmen wir hiervon die Norm in Bezug auf $\Omega_m$, die das Quadrat der Norm in Bezug auf $H_m$ ist, so folgt (nach \S\,169)
\[
q^{2f_1}=q^\lambda,
\]
also $f_1=\frac12\lambda$. Da $f_1$ eine ganze Zahl ist, so muss $\lambda$ mindestens $=2$ sein. Es kann aber $\lambda$ auch nicht grösser als $2$ sein, da $\sigma^2$ mit der in $H_m$ enthaltenen ganzen Zahl $(1-r)(1-r^{-1})$ associirt ist, also $\sigma^2$ durch $\frP$ theilbar sein muss. Demnach ist $f_1=1$, und wir erhalten den Satz:

\begin{enumerate}[label=\arabic*.,leftmargin=2.2em]
\item Die Primzahl $q$ ist die $\frac12\mu$te Potenz einer in $H_m$ existirenden Primzahl ersten Grades.
\end{enumerate}

Es sei nun $p$ von $q$ verschieden, und $\frp$ ein Primfactor von $\frP$ im Körper $\Omega_m$ vom Grade $f$, der durch die Substitution $(r,r^{-1})$ in $\frp'$ übergeht. Dann ist $\frP$ sowohl durch $\frp$ als durch $\frp'$ theilbar, und andererseits ist $\frp\frp'$ in $H_m$ enthalten und also durch $\frP$ theilbar, und wir haben zu unterscheiden, ob $\frp,\frp'$ verschieden sind oder nicht.

Ist $\frp$ von $\frp'$ verschieden, so ist $\frP$ durch $\frp\frp'$ theilbar, und es folgt $\frP=\frp\frp'$, und durch Normbildung $f_1=f$.

Ist aber $\frp=\frp'$, so ist $\frP=\frp$, und die Normbildung ergiebt
\[
2f_1=f.
\]
Welcher dieser beiden Fälle eintritt, das hängt also davon ab, ob das Primideal $\frp$ die Substitution $(r,r^{-1})$ gestattet oder nicht, d. h. ob $(r,r^{-1})$ in der Gruppe, zu der das Primideal gehört, enthalten ist oder nicht. Nach \S\,172 ist dieser Unterschied dadurch bedingt, ob es irgend einen Exponenten $h$ giebt, für den die Congruenz
\[
p^h\equiv -1\pmod m
\]
erfüllt ist, oder ob es keinen solchen Exponenten giebt.

Wir fassen das Ergebniss folgendermaassen zusammen:

\begin{enumerate}[label=\arabic*.,leftmargin=2.2em,start=2]
\item Die von $q$ verschiedenen natürlichen Primzahlen $p$ zerfallen in zwei Arten $p_1,p_2$.

Die erste Art $p_1$ ist dadurch charakterisirt, dass für irgend einen Exponenten $h$
\[
p_1^h\equiv -1\pmod m,
\]
und diese Primzahlen zerfallen, wenn sie für den Modul $m$ zum Exponenten $f$ gehören und $\mu=ef$ gesetzt wird, im Körper $H_m$ in $e$ Primfactoren $\frac12 f$ten Grades. Ihre Primfactoren sind auch in $\Omega_m$ nicht weiter zerlegbar.

Die Primzahlen der zweiten Art $p_2$ sind dadurch bestimmt, dass keine ihrer Potenzen nach dem Modul $m$ mit $-1$ congruent wird, und sie zerfallen im Körper $H_m$ in $\frac12e$ Primfactoren $f$ten Grades. Ihre Primfactoren sind in $\Omega_m$ noch in je zwei Factoren zerlegbar.
\end{enumerate}

Bei den Primzahlen der ersten Art muss $f$ sicher eine gerade Zahl sein. Wenn $m$ ungerade ist, so genügt es auch, damit $p$ eine Primzahl von der ersten Art sei, dass sie zu einem geraden Exponenten gehöre; denn dann ist
\[
(p^{\frac12 f}-1)(p^{\frac12 f}+1)\equiv0\pmod m.
\]
Der erste dieser Factoren $p^{\frac12 f}-1$ ist aber nicht durch $m$ theilbar, weil sonst $p$ zum Exponenten $\frac12 f$ gehören würde. Folglich muss $p^{\frac12 f}+1$ durch $q$ theilbar sein. Hier können aber nicht beide Factoren $p^{\frac12 f}-1$ und $p^{\frac12 f}+1$ durch $q$ theilbar sein, weil sonst auch ihre Differenz, die $=2$ ist, durch $q$ theilbar wäre, und folglich ist $p^{\frac12 f}+1$ durch $m$ theilbar.

Ist aber $m$ eine Potenz von $2$, so ist die Congruenz
\[
p^h\equiv -1\pmod m
\]
nur dann möglich, wenn $p\equiv -1\pmod m$ ist. Denn jede gerade Potenz einer ungeraden Zahl ist $\equiv+1\pmod 8$ und kann also nicht $\equiv-1\pmod m$ sein. Ist aber $h$ ungerade, so ist
\[
\frac{p^h+1}{p+1}=p^{h-1}-p^{h-2}+\cdots+1
\]
selbst ungerade, und folglich ist $p^h+1$ durch keine höhere Potenz von $2$ theilbar, als $p+1$. Demnach können wir dem Satze 2. noch folgende nähere Bestimmung hinzufügen:

\begin{enumerate}[label=\arabic*.,leftmargin=2.2em,start=3]
\item Wenn $m$ ungerade ist, so gehören die Primzahlen $p_1$ zu einem geraden, die Primzahlen $p_2$ zu einem ungeraden Exponenten.

Ist $m$ gerade, so gehören die Primzahlen der Form $km-1$ zur ersten und alle anderen ungeraden Primzahlen zur zweiten Art.
\end{enumerate}

\section*{\S\,178. Die Einheiten des Körpers $H_m$.}

Die dem Körper $H_m$ angehörigen Einheiten sind von besonderer Wichtigkeit für die Anwendungen. Auf sie lassen sich auch die Einheiten des Körpers $\Omega_m$ zurückführen nach folgendem Satze:

\begin{enumerate}[label=\arabic*.,leftmargin=2.2em]
\item Jede Einheit des Körpers $\Omega_m$ ist das Product einer Einheitswurzel und einer reellen Einheit des Körpers $H_m$.
\end{enumerate}

Bezeichnen wir irgend eine Einheit des Körpers $\Omega_m$ mit $\mathscr E(r)$, so ist $\mathscr E(r^{-1})$ der conjugirt imaginäre Werth, der gleichfalls eine Einheit darstellt, und der Quotient $\mathscr E(r):\mathscr E(r^{-1})$ ist wieder eine Einheit, deren absoluter Werth
\begin{equation}
\sqrt{\frac{\mathscr E(r)}{\mathscr E(r^{-1})}\frac{\mathscr E(r^{-1})}{\mathscr E(r)}}
\tag{1}
\end{equation}
gleich $1$ ist, und folglich ist dieser Quotient eine Einheitswurzel und mithin $=\pm r^h$, worin $h$ irgend ein ganzzahliger Exponent ist (\S\,175, Satz 1., 2.). Wir haben also
\begin{equation}
\mathscr E(r)=\pm r^h\mathscr E(r^{-1}).
\tag{2}
\end{equation}

Es ist nun im weiteren Beweise ein kleiner Unterschied zu machen, je nachdem $m$ ungerade oder gerade ist.

Ist $m$ zunächst ungerade, so können wir in (2) den Exponenten $h$ gerade annehmen, da wir ihn nöthigenfalls durch $h+m$ ersetzen können. Ferner ist in der Formel (2) nur das obere Zeichen zulässig. Denn wenn das untere Zeichen stände, so würde folgen:
\[
\mathscr E(r)\equiv\mathscr E(1)\equiv-\mathscr E(1),\qquad 2\mathscr E(1)\equiv0\,[\mathrm{mod}\,(1-r)].
\]
Nach \S\,169 ist aber $1-r$ ein Theiler von $q$, also relativ prim zu $2$, und daher ist $\mathscr E(1)$ und folglich auch $\mathscr E(r)$ durch $(1-r)$ theilbar. Dies aber widerspricht der Voraussetzung, dass $\mathscr E(r)$ eine Einheit sein soll.

Demnach ergiebt sich aus (2):
\[
r^{-\frac12 h}\mathscr E(r)=r^{\frac12 h}\mathscr E(r^{-1}),
\]
woraus folgt, dass $r^{-\frac12 h}\mathscr E(r)$ eine reelle Einheit ist. Bezeichnen wir sie mit $e(r)$, so ist also
\begin{equation}
\mathscr E(r)=r^{\frac12 h}e(r),
\tag{3}
\end{equation}
woraus für diesen Fall das Theorem 1. bewiesen ist.

Wenn aber zweitens $m$ eine Potenz von $2$ ist, so ist $\mu=\frac12m$ und $r^\mu=-1$, und zugleich mit $r$ ist $-r$ eine $m$te Einheitswurzel. Wenn wir nöthigenfalls $h$ durch $h+\frac12m$ ersetzen, so können wir in (2) das obere Zeichen annehmen und setzen:
\begin{equation}
\mathscr E(r)=r^h\mathscr E(r^{-1}).
\tag{4}
\end{equation}
Es kommt jetzt darauf an, nachzuweisen, dass $h$ gerade sein muss.

Wenn wir $\mathscr E(r)$ durch die Basis $1,r,r^2,\ldots,r^{\mu-1}$ darstellen, so ergiebt sich
\begin{equation}
\mathscr E(r)=\sum_{s=0}^{\mu-1}x_sr^s,
\tag{5}
\end{equation}
worin die $x_s$ ganze rationale Zahlen sind. Nach (4) ist aber dann
\[
\sum_{s=0}^{\mu-1}x_sr^s=\sum_{s=1}^{\mu-1}x_sr^{h-s},
\]
und daraus folgt, dass $x_s=\pm x_{s'}$ ist, wenn $s+s'\equiv h\pmod\mu$. Ist nun $h$ ungerade, so ist aus zwei so verbundenen Zahlen die eine gerade, die andere ungerade (da $\mu$ eine Potenz von $2$ ist). Der Ausdruck (5) ergiebt also für $\mathscr E(r)$ ein Aggregat von Gliedern der Form
\[
x_s(r^s\pm r^{s'}).
\]
Es ist aber
\[
r^s\pm r^{s'}=r^s(1\pm r^{s'-s})=r^s(1\mp r^{\mu+s'-s})
\]
immer durch $1-r$ theilbar, und daraus würde folgen, dass $\mathscr E(r)$ durch $1-r$, was keine Einheit ist, theilbar wäre, was dem Begriffe der Einheit widerspricht. Also ist die Annahme unzulässig, dass in der Formel (4) der Exponent $h$ ungerade sei, und wir können wie oben
\[
\mathscr E(r)=r^{\frac12h}e(r)
\]
setzen, worin $e(r)$ eine reelle Einheit ist. Damit ist der Satz 1. bewiesen.

Wir wollen ein System von reellen Einheiten hervorheben: Nach \S\,169 sind die $\mu$ Grössen $1-r^n$ alle unter einander associirt. Demnach ist der Quotient
\[
\frac{1-r^n}{1-r^{n'}},
\]
worin $n,n'$ zwei relative Primzahlen zu $m$ bedeuten, eine Einheit. Daraus leitet man aber, wenn
\[
r=e^{\frac{2\pi i}{m}}
\]
gesetzt wird, die reelle Einheit
\begin{equation}
\frac{\sin\frac{n\pi}{m}}{\sin\frac{n'\pi}{m}}
=\frac{r^{\frac n2}-r^{-\frac n2}}{r^{\frac{n'}2}-r^{-\frac{n'}2}}
=r^{\frac{n'-n}{2}}\frac{1-r^n}{1-r^{n'}}
\tag{6}
\end{equation}
her. Ist $m$ gerade, so kann man $n'=n+\frac m2$ setzen, und erhält für diesen Fall die reellen Einheiten
\begin{equation}
\tau_n=\operatorname{tang}\frac{n\pi}{m},
\tag{7}
\end{equation}
worin $n$ jede ungerade Zahl bedeuten kann.

Ersetzt man $n$ durch $\frac12m-n$, so geht $\tau_n$ in den reciproken Werth über. Lässt man $n$ die Reihe der Zahlen
\[
1,3,5,\ldots,\frac m4-1
\]
durchlaufen, so erhält $\tau_n$ lauter positive echt gebrochene Werthe.

\clearpage

\srcpage{648}
\section*{Zweiundzwanzigster Abschnitt. Abel'sche Körper und Kreistheilungskörper.}

\section*{\S~179. Einfache Abel'sche Körper.}

Wir wenden uns nun zu einer der interessantesten Anwendungen der Theorie der algebraischen Zahlen.

Es handelt sich um den Beweis des allgemeinen Satzes:\footnote{Kronecker hat diesen Satz zuerst ausgesprochen, aber keinen vollständigen Beweis dafür veröffentlicht. Den ersten Beweis hat der Verfasser des vorliegenden Werkes in Bd.~8 der \textit{Acta Mathematica} bekannt gemacht (1886). In neuester Zeit ist ein zweiter Beweis von Hilbert veröffentlicht, der sich auf die im neunzehnten Abschnitte entwickelten Sätze stützt (Nachrichten d. Gesellschaft d. Wissenschaften zu Göttingen, 1896, Heft~1).}

\textbf{I. Alle im absoluten Rationalitätsbereich Abel'schen Zahlkörper sind Kreistheilungskörper.}

Haben wir diesen Satz bewiesen, so gewinnen die Untersuchungen des dritten Abschnittes über Kreistheilungskörper ein erhöhtes Interesse, weil damit gezeigt ist, dass auf dem dort angegebenen Wege nicht nur alle Kreistheilungskörper, sondern alle Abel'schen Körper überhaupt gefunden werden.

Es sei $\Om=R(x)$ ein Abel'scher Körper $m$-ten Grades, d.~h. ein Körper, der aus allen rationalen Functionen einer Wurzel $x$ einer irreduciblen Abel'schen Gleichung $m$-ten Grades besteht, wenn als Rationalitätsbereich der Körper $R$ der rationalen Zahlen betrachtet wird. Die Galois'sche Gruppe $\Gg$ dieser Gleichung, die also eine Abel'sche Gruppe ist und denselben Grad $m$ hat, wie der Körper $\Om$, ist auch die Gruppe des Körpers $\Om$.

Ist $x_1$ eine nicht primitive Zahl aus $\Om$, so ist $\Om_1=R(x_1)$ gleichfalls ein Abel'scher Körper, den wir einen echten Theiler von $\Om$ nennen und dessen Grad ein echter Theiler des Grades von $\Om$ ist. Denn ist $\Aa$ die Gruppe, zu der die Zahl $x_1$ gehört, so ist $\Gg\mid\Aa$ (\S~4 dieses Bandes) die Gruppe des Körpers $\Om_1$ und dies ist zugleich mit $\Gg$ eine Abel'sche Gruppe. Ein Theiler von $\Om$, der mit $\Om$ von gleichem Grade ist, ist mit $\Om$ identisch (vgl. Bd.~I, \S\S~144, 149, 156).

\srcpage{649}
Giebt es nun zwei nicht primitive Zahlen $x_1,x_2$ in $\Om$ von der Art, dass
\[
  \Om=R(x_1,x_2)
\]
gesetzt werden kann, so sind die Körper $\Om_1=R(x_1)$, $\Om_2=R(x_2)$ echte Theiler von $\Om$, und $\Om$ heisst aus den beiden Körpern $\Om_1,\Om_2$ zusammengesetzt.

Gestattet der Körper $\Om$ keine solche Darstellung, so heisst er einfach.

Wenn also als bewiesen vorausgesetzt wird, dass $\Om_1,\Om_2$ Kreistheilungskörper sind, d.~h. dass alle ihre Zahlen rational durch Einheitswurzeln darstellbar sind, so folgt dasselbe für $\Om$.

Wenn einer der Körper $\Om_1,\Om_2$ zerlegbar ist, so können wir die Zerlegung wiederholen, und müssen, da die Grade abnehmende ganze positive Zahlen sind, endlich zu einfachen Körpern gelangen.

Das Theorem I. braucht also nur für einfache Abel'sche Körper bewiesen zu werden.

Es ist aber daran zu erinnern, dass die Zerlegbarkeit eines Körpers $\Om$ keineswegs schon aus der Existenz eines Theilers folgt, und dass bei einem zerlegbaren Körper die Zerlegung in einfache Körper auf ganz verschiedene Arten geschehen kann, so dass bei den Körpern nicht die Gesetze der Theilbarkeit gelten, wie im Gebiete der Zahlen.

Ein Kennzeichen für die Einfachheit eines Körpers können wir aus seiner Gruppe ableiten.

Wenn die Gruppe $\Gg$ zwei echte Theiler $\Aa$ und $\Bb$ von der Art hat, dass jedes Element ein- und nur einmal aus einem Element von $\Aa$ und einem Element von $\Bb$ zusammengesetzt werden kann, so nennen wir die Gruppe $\Gg$ zerlegbar in die beiden Componenten $\Aa,\Bb$; $\Gg$ ist also zerlegbar, wenn in dem nach der Composition der Theile (\S~4) gebildeten Producte
\begin{equation}
  \Gg=\Aa\Bb
\tag{1}
\end{equation}
jedes Element von $\Gg$ ein- und nur einmal erscheint. Der Grad von $\Gg$ ist dann gleich dem Producte der Grade von $\Aa$ und $\Bb$.

Giebt es solche Theiler nicht, so heisst die Gruppe $\Gg$ unzerlegbar.

\srcpage{650}
Dann können wir den Satz beweisen:

\textbf{1. Ist die Gruppe eines Abel'schen Körpers zerlegbar, so ist auch der Körper zusammengesetzt.}

Nehmen wir, um dies zu beweisen, an, die Gruppe $\Gg$ des Körpers $\Om$ sei nach der Formel (1) zerlegbar und $m,a,b$ seien die Grade von $\Gg,\Aa,\Bb$. Wir nehmen eine zur Gruppe $\Bb$ gehörige Zahl $\xi$ des Körpers $\Om$, die also durch die Substitutionen von $\Gg$ in $a$ verschiedene conjugirte Werthe übergeht und ebenso eine zu $\Aa$ gehörige $b$-werthige Zahl $\eta$. Wir können dann eine Zahl
\[
  x=\alpha\xi+\beta\eta
\]
mit rationalen Zahlencoefficienten $\alpha,\beta$ bilden, die $m$ verschiedene conjugirte Werthe hat, und dann ist $\Om=R(x)$. Also ist $\Om$ aus $R(\xi)$ und $R(\eta)$ zusammengesetzt (Bd.~I, \S~143).

Erinnern wir uns nun an die in \S~9, 10 dieses Bandes abgeleitete Darstellung einer Abel'schen Gruppe durch eine Basis, so ergiebt sich durch wiederholte Anwendung des eben Bewiesenen:

\textbf{2. Jeder Abel'sche Körper $\Om$ lässt sich aus solchen zusammensetzen, deren Grad eine Primzahlpotenz und deren Gruppe cyklisch ist. Die Grade dieser zusammensetzenden Körper sind die Invarianten der Gruppe von $\Om$.}

Diese Sätze werden durch den folgenden ergänzt:

\textbf{3. Ein Abel'scher Körper von Primzahlpotenzgrad mit cyklischer Gruppe ist einfach.}

Wenn die Gruppe $\Gg$ cyklisch ist, so lassen sich ihre Elemente durch die Potenzen einer Basis, $G^\gamma$, darstellen, wenn $\gamma$ ein volles Restsystem nach dem Modul $m$ durchläuft. Wir erhalten alle Theiler $\Aa$ von $\Gg$ dargestellt durch
\[
  \Aa=G^{b\gamma},
\]
wenn $b$ ein Theiler von $m$ ist und $\gamma$ ein volles Restsystem nach dem Modul $a=m:b$ durchläuft.

Ist
\[
  \Aa'=G^{b'\gamma}
\]
ein zweiter Theiler von $\Gg$ vom Grade $a'$ und ist
\[
  b'\le b,
  \qquad a'\ge a,
\]
so ist, wenn wir jetzt annehmen, dass $m$ und mithin auch $a,b,a',b'$ Potenzen einer Primzahl $q$ sind, $b'$ ein Theiler von $b$ und folglich $\Aa$ ein Theiler von $\Aa'$.

\srcpage{651}
Wenn also $\xi$ eine Zahl des Körpers $\Om$ ist, die zu der Gruppe $\Aa$ gehört, so ist $\xi$ eine primitive Zahl des Körpers $b$-ten Grades $R(\xi)$, und in diesem Körper ist auch jede Zahl $\xi'$ enthalten, die die Substitutionen der Gruppe $\Aa'$ gestattet. Sind also $\Aa,\Bb$ und $\Aa',\Bb'$ echte Theiler von $\Gg$, so ist $R(\xi)$ von $\Om$ verschieden, und $R(\xi)$ wird durch Adjunction von $\xi'$ nicht erweitert. Es kann also auch nicht $\Om$ aus $R(\xi)$ und $R(\xi')$ zusammengesetzt werden, wodurch 3. bewiesen ist.

\section*{\S~180. Die Resolventen.}

Nach dem, was jetzt bewiesen ist, können wir uns in den weiteren Betrachtungen, die den Beweis des Satzes I. zum Ziele haben, auf solche Abel'sche Körper beschränken, deren Grad eine Primzahlpotenz und deren Gruppe cyklisch ist. Aus diesem Grunde genügte es auch für die beabsichtigte Anwendung, dass wir im vorigen Abschnitte uns auf die Betrachtung der Kreistheilungskörper $\Omega_m$ beschränkt haben, in denen $m$ eine Primzahlpotenz war. Wir behalten so viel als möglich die dort gebrauchte Bezeichnung bei und setzen
\begin{equation}
  m=q^\kappa,
\tag{1}
\end{equation}
worin $q$ eine Primzahl, $\kappa$ ein positiver Exponent ist, und in dem Falle, dass $q=2$ ist, nehmen wir $\kappa>1$ an. Mit $n$ bezeichnen wir jede durch $q$ nicht theilbare Zahl. Es ist $r$ eine primitive $m$-te Einheitswurzel und $\Omega_m$ der Körper der rationalen Functionen von $r$.

Es sei nun $\Om=R(x)$ ein einfacher Abel'scher Körper $m$-ten Grades, so dass $x$ Wurzel einer rationalen irreduciblen cyklischen Gleichung $m$-ten Grades ist. Bezeichnen wir die Wurzeln dieser Gleichung in geeigneter Reihenfolge mit
\begin{equation}
  x_0,x_1,x_2,\ldots,x_{m-1},
\tag{2}
\end{equation}
und nehmen $x_m=x_0$, $x_{m+1}=x_1,\ldots$ an, so ist
\[
  x_1=\Phi(x_0),\quad x_2=\Phi(x_1),\ldots,\quad x_0=\Phi(x_{m-1}),
\]
worin $\Phi$ eine ganze Function mit rationalen Coefficienten bedeutet, und die Gruppe des Körpers $\Om$ besteht aus den Potenzen der Substitution $(x_0,x_1)$ oder aus den cyklischen Permutationen

\srcpage{652}
der Wurzeln (2). Jede cyklische Function dieser Grössen ist eine rationale Zahl. Was wir zu beweisen haben, ist, dass sich die $x$ rational durch Einheitswurzeln ausdrücken lassen. Adjungiren wir dem Körper $\Om$ noch die $m$-te Einheitswurzel $r$, so entsteht ein Körper $R(x,r)$, der die beiden Körper $\Om=R(x)$ und $R(r)=\Omega_m$ als Theiler enthält.

Wenn eine Zahl $\omega=\Phi(x,r)$ des Körpers $R(x,r)$ die sämmtlichen Substitutionen $(x_0,x_1)$, $(x_0,x_2)$, \ldots, $(x_0,x_{m-1})$ gestattet, so ist sie im Körper $\Omega_m$ enthalten; denn dann ist
\[
  m\omega=\Phi(x_0,r)+\Phi(x_1,r)+\cdots+\Phi(x_{m-1},r),
\]
und darin sind die Coefficienten der Potenzen von $r$ nicht nur cyklische, sondern sogar symmetrische Functionen der Wurzeln (2), also rationale Zahlen. Dies gilt selbst noch in dem Falle, dass die Gleichung für $x$ durch Adjunction von $r$ reducirt wird.

Dies wenden wir an auf die Resolventen
\begin{equation}
  (r,x_0)=x_0+rx_1+r^2x_2+\cdots+r^{m-1}x_{m-1},
\tag{3}
\end{equation}
und allgemeiner, wenn $s$ einen beliebigen ganzzahligen Exponenten bedeutet:
\begin{equation}
  (r^s,x_0)=x_0+r^s x_1+r^{2s}x_2+\cdots+r^{(m-1)s}x_{m-1}.
\tag{4}
\end{equation}
Durch diese Resolventen lässt sich $x_0$ selbst wieder rational ausdrücken, nämlich:
\begin{equation}
  mx_0=\sum_{s=0}^{m-1}(r^s,x_0),
\tag{5}
\end{equation}
und wenn wir also nachweisen können, dass alle diese Resolventen $(r^s,x_0)$ Kreistheilungszahlen sind, so haben wir unser Ziel erreicht. Wenn wir
\begin{equation}
  (r^s,x_\lambda)=x_\lambda+r^s x_{\lambda+1}+r^{2s}x_{\lambda+2}+\cdots+r^{(m-1)s}x_{\lambda+m-1}
\tag{6}
\end{equation}
setzen, so geht $(r^s,x_\lambda)$ aus $(r^s,x_0)$ durch die Substitution $(x_0,x_\lambda)$ hervor, und es besteht die fundamentale Relation
\begin{equation}
  r^{\lambda s}(r^s,x_\lambda)=(r^s,x_0).
\tag{7}
\end{equation}
Alle diese Resolventen sind Zahlen des Körpers $R(x,r)$; wir betrachten vorzugsweise die beiden folgenden Verbindungen:
\begin{align}
  (r^s,x_0)^m &= \omega_s, \tag{8}\\
  (r^s,x_0)^{-1}(r,x_0)^s &= \alpha. \tag{9}
\end{align}

\srcpage{653}
Diese beiden Zahlen gestatten, wie aus (7) unmittelbar hervorgeht, die Substitutionen $(x_0,x_\lambda)$, und sind also Zahlen des Körpers $\Omega_m$.

Ist $n$ durch $q$ nicht theilbar, so kann wegen der Irreducibilität der Kreistheilungsgleichung keine von den Resolventen $(r^n,x_0)$ verschwinden, wenn sie nicht alle verschwinden; denn die Gleichung
\[
  (r,x_0)^m=0
\]
wäre eine Gleichung für $r$ mit rationalen Coefficienten, die alle Substitutionen $(r,r^n)$ gestattet.

Ebenso kann keine der Resolventen $(r^s,x_0)$ verschwinden, wenn nicht die ganze Schaar verschwindet, in der $s$ einen bestimmten grössten gemeinschaftlichen Theiler mit $m$ hat.

Wenn alle $(r,x_0)$ verschwinden, so ist die Formel (9) nicht anwendbar. Dann können wir aber die Resolventen $(r^q,x_0)$ direct betrachten, die zu einem Körper vom Grade $mq^{-1}$ gehören, der aus
\[
  y_0=x_0+x_q+x_{2q}+\cdots+x_{m-q}
\]
hervorgeht, und wenn auch diese Grössen verschwinden, so gehen wir zu den Resolventen $(r^{q^2},x_0)$ über u.~s.~f.

Wenn aber $(r,x_0)$ nicht verschwindet, so lassen sich durch die Formel (9) alle $(r^s,x_0)$ rational durch $(r,x_0)$ und durch Kreistheilungszahlen darstellen.

Demnach ist es für unseren Zweck ausreichend, wenn wir beweisen können, dass die nicht verschwindenden Resolventen $(r,x_0)$ Kreistheilungszahlen sind.

Lassen wir $n$ ein volles System incongruenter, durch $q$ nicht theilbarer Zahlen durchlaufen, so ist
\[
  \sum n\equiv0\pmod m,
\]
denn die Zahlen $n$ zerfallen in Paare einander zu $m$ ergänzender Zahlen. Daraus folgt nach (7), dass das Product
\begin{equation}
  \prod_n (r^n,x_0)=a
\tag{10}
\end{equation}
eine Zahl in $\Omega_m$ ist. Da aber diese Zahl sich überdies nicht ändert, wenn irgend eine der Substitutionen $(r,r^n)$ darin ausgeführt wird, so ist $a$ eine rationale Zahl.

\srcpage{654}
\section*{\S~181. Vorbereitung zum Beweis.}

Nach den Ausführungen des vorigen Paragraphen ist zum Beweis des grossen Theorems I nur noch der Nachweis erforderlich, dass die in $\Omega_m$ enthaltene Zahl
\[
  \omega=(r,x_0)^m
\]
die $m$-te Potenz einer Kreistheilungszahl ist.

\textbf{1. Die wesentliche Eigenschaft der zu $\omega$ conjugirten Zahlen $\omega_n$, auf die sich der Beweis stützt, ist die, dass}
\begin{equation}
  \omega_n^{-1}\omega_1^n=\alpha^m
\tag{1}
\end{equation}
\textbf{die $m$-te Potenz einer Zahl in $\Omega_m$ ist.}

Hierin kann $n$ jede durch $q$ nicht theilbare Zahl sein. Bilden wir von der rechten und linken Seite von (1) die Normen im Körper $\Omega_m$, so ergiebt sich, wenn wir die Norm der Zahl $\alpha$, die eine rationale Zahl ist, mit $a$ bezeichnen, da jedes $\omega_n$ dieselbe Norm wie $\omega$ hat,
\[
  N(\omega)^{n-1}=a^m.
\]
Wenn nun $m$ ungerade ist, so können wir hierin $n=2$ annehmen, und erhalten als Folgerung aus 1.:

\textbf{2. Die Norm von $\omega$ ist die $m$-te Potenz einer rationalen Zahl}
\begin{equation}
  N(\omega)=a^m.
\tag{2}
\end{equation}
Ist $m$ eine Potenz von 2, so ist diese Folgerung nicht mehr zulässig. Es ergiebt sich dann aus 1. nur, dass $N(\omega)^2$ die $m$-te Potenz einer rationalen Zahl ist. Gleichwohl müssen die Zahlen $\omega$, wie aus (10) des vorigen Paragraphen zu sehen ist, auch in diesem Falle noch der Bedingung 2. genügen, die dann aber nicht mehr von selbst schon in 1. enthalten ist, sondern als neue Eigenschaft hinzukommt.

Was wir zu beweisen haben, ist, dass aus 1. und 2. folgt, dass auch $\omega$ selbst die $m$-te Potenz einer Kreistheilungszahl, wenn auch in einem höheren Körper, ist.

Dazu ist aber erforderlich, die Zerlegung der Zahlen $\omega$ in ihre Primfactoren im Körper $\Omega_m$ zu untersuchen.

Die Primzahl $q$ ist, wie im \S~169 nachgewiesen, mit der $\varphi(m)$-ten Potenz einer in $\Omega_m$ existirenden Primzahl $\sigma=1-r$

\srcpage{655}
associirt. Ist also $\sigma^k$ die in $\omega$ enthaltene Potenz von $\sigma$, so setzen wir
\[
  \omega=\sigma^k\omega',
\]
worin $\omega'$ mit der Primzahl $q$ theilerfremd ist. Die Bildung der Norm ergiebt nach 2.
\[
  a^m=q^k b,
\]
worin $b$ eine rationale Zahl ist, die in einfachster Gestalt die Primzahl $q$ weder im Zähler noch im Nenner enthält. Daraus ergiebt sich aber, dass $k$ durch $m$ theilbar sein muss, und daraus das erste Resultat:

\textbf{3. Der Exponent der in $\omega$ enthaltenen Primzahl $\sigma$ ist immer durch $m$ theilbar.}

Es sei nun ferner $p$ eine von $q$ verschiedene Primzahl und $\pp_\xi$ ein Primfactor von $p$ im Körper $\Omega_m$. Wir lassen $\xi$ die im \S~172 in den verschiedenen Fällen näher bestimmte Zahlenreihe $\xi_1,\xi_2,\ldots,\xi_e$ durchlaufen, so dass
\begin{equation}
  p=\pp_{\xi_1}\pp_{\xi_2}\cdots\pp_{\xi_e}
\tag{3}
\end{equation}
wird, und nehmen an, es sei $\pp_\xi^{k_\xi}$ die in $\omega$ enthaltene Potenz von $\pp_\xi$. Die in $\omega_n$ enthaltene Potenz von $\pp_{n\xi}$ ist dann $\pp_{n\xi}^{k_\xi}$, und wenn wir $n'$ aus der Congruenz
\begin{equation}
  nn'\equiv1\pmod m
\tag{4}
\end{equation}
bestimmen, so ist $\pp_\xi^{k_{n'\xi}}$ die in $\omega_n$ enthaltene Potenz von $\pp_\xi$. Demnach ist in
\begin{equation}
  \omega_n^{-1}\omega_1^n
\tag{5}
\end{equation}
die Potenz $\pp_\xi^{n k_\xi-k_{n'\xi}}$ enthalten, und wegen 1. muss der Exponent dieser Potenz durch $m$ theilbar sein.

Wir erhalten also für jedes durch $q$ nicht theilbare $n$ die Congruenz
\begin{equation}
  n k_\xi\equiv k_{n'\xi}\pmod m,
\tag{6}
\end{equation}
in der wir nun auch $\xi$ durch $n\xi$ ersetzen können, wodurch sich
\begin{equation}
  n k_{n\xi}\equiv k_\xi\pmod m
\tag{7}
\end{equation}
ergiebt. Nehmen wir hierin $\xi=1$ und setzen dann $\xi,\xi'$ für $n,n'$ und $k$ für $k_1$, so folgt aus (7)
\begin{equation}
  k_\xi\equiv\xi'k\pmod m,
\tag{8}
\end{equation}
worin nun $k$ für alle $\xi$ denselben Werth hat und $\xi,\xi'$ durch die Congruenz
\[
  \xi\xi'\equiv1\pmod m
\]
zusammenhängen.

\srcpage{656}
Nehmen wir $n=p$ an, so wird $\pp_{p\xi}=\pp_\xi$ (\S~172), also ist auch $k_{p\xi}=k_\xi$, und aus (7) und (8) folgt:
\begin{equation}
  k(p-1)\equiv0\pmod m.
\tag{9}
\end{equation}
Hieraus ergeben sich nun wichtige Folgerungen: Wenn zunächst $p-1$ nicht durch $q$ theilbar ist, so folgt, dass $k$ durch $m$ theilbar sein muss, und wir haben also den Satz:

\textbf{4. Die Primfactoren einer Primzahl $p$, die nach dem Modul $q$ nicht mit $1$ congruent ist, sind in $\omega$ nur in solchen Potenzen enthalten, deren Exponenten durch $m$ theilbar sind.}

Wenn in dem Falle, wo $m$ eine Potenz von 2 ist, $p-1$ durch 2, aber nicht durch 4 theilbar ist, wenn also $p\equiv-1\pmod4$ ist, so folgt aus (9) nur, dass $k$ durch $\tfrac12m$, nicht aber, dass $k$ durch $m$ theilbar ist. Der Exponent der in $\omega$ enthaltenen Potenz von $\pp_\xi$ ist, von Vielfachen von $m$ abgesehen, gleich $k$, und wenn $k$ durch $m$ theilbar ist, so ist alles wie im Falle 4. Ist aber $k$ nur durch $\tfrac12m$, nicht durch $m$ theilbar, also
\[
  k\equiv\tfrac12m\pmod m,
\]
so ist $k-\tfrac12m$ durch $m$ theilbar, und wir können mit Rücksicht auf die Zerlegung (3) den Satz so formuliren:

\textbf{5. Ist $m$ eine Potenz von 2 und $p$ eine Primzahl congruent mit $-1$ nach dem Modul 4, so lässt sich der Exponent $h$ so bestimmen, dass alle Primfactoren von $p$ in}
\[
  \omega\sqrt p^{\,-hm}
\]
\textbf{nur zu solchen Potenzen enthalten sind, deren Exponent durch $m$ theilbar ist.}

Es sei endlich $p-1$ durch $q$, oder, falls $m$ gerade ist, durch 4 theilbar. Ist $m_1$ der grösste gemeinschaftliche Theiler von $m$ und $p-1$, und ist
\[
  m=m_1m_2,
\]
so muss $k$ nach (9) durch $m_2$ theilbar sein, und wenn wir
\[
  k=hm_2
\]
setzen, so ist in $\omega$ nach (8) das Product enthalten:
\begin{equation}
  \left(\prod_\xi\pp_\xi^{\xi'}\right)^{hm_2}.
\tag{10}
\end{equation}
Ausserdem kommen in $\omega$ die Primfactoren $\pp_\xi$ nur noch in solchen Potenzen vor, deren Exponenten durch $m$ theilbar sind.

\srcpage{657}
In diesem Falle durchläuft nun $\xi$ nach \S~172, III ein volles System durch $q$ nicht theilbarer Reste nach dem Modul $m_1$, und wir können in (10) unter $\xi'$ die kleinste positive Lösung der Congruenz
\[
  \xi\xi'\equiv1\pmod{m_1}
\]
verstehen, weil, wenn wir die Exponenten $\xi'$ in (10) nach dem Modul $m_1$ verändern, nur $m$-te Potenzen der Primfactoren $p$ hinzutreten.

Dann können wir auf das Product (10) das Kummer'sche Theorem [\S~174, (16)] anwenden, indem wir dort $m$ durch $m_1$ ersetzen, und erhalten, wenn wir unter $r_1=r^{m_2}$ eine $m_1$-te Einheitswurzel, unter den $\eta$ gewisse Perioden der $p$-ten Einheitswurzeln verstehen:
\[
  \prod_\xi\pp_\xi^{\xi'}=(r_1,\eta)^{m_1},
\]
und folglich
\begin{equation}
  \left(\prod_\xi\pp_\xi^{\xi'}\right)^{hm_2}=(r_1,\eta)^{hm}.
\tag{11}
\end{equation}
Demnach haben wir das Theorem:

\textbf{6. Ist $p$ eine Primzahl und $p-1$ durch $q$, oder, wenn $q=2$ ist, durch 4 theilbar, so lässt sich der positive Exponent $h$ so bestimmen, dass die in $\Omega_m$ enthaltene Zahl}
\[
  \omega(r_1,\eta)^{-hm}
\]
\textbf{die Primfactoren von $p$ nur zu solchen Potenzen enthält, deren Exponent durch $m$ theilbar ist.}

Es ist jetzt auch nicht mehr nöthig, den Satz 5. von dem Satze 6. zu unterscheiden; denn für $m_1=2$ geht 6. in 5. über, weil dann $r_1=-1$ wird, und nach Bd.~I, \S~171
\[
  (r_1,\eta)^m=(-1,\eta)^m=\sqrt{p^m}
\]
ist.\footnote{In des Verfassers Abhandlung in Bd.~8 der \textit{Acta mathematica} ist es versäumt, den Fall des Satzes 5. besonders hervorzuheben.}

Fassen wir das Ergebniss der Sätze 3. bis 6. in eine Formel zusammen, so ergiebt sich, wenn $\eps$ ein Einheitsfunctional, $\varphi$ irgend ein Functional des Körpers $\Omega_m$ bedeutet, und $p$ eine Reihe von Primzahlen durchläuft, die congruent mit 1 nach dem Modul $q$ sind, worunter dieselbe Primzahl $p$ auch mehrmals vorkommen kann,
\begin{equation}
  \omega=(r,x_0)^m=\eps\varphi^m\prod_p(r^{m_2},\eta)^m.
\tag{12}
\end{equation}

\srcpage{658}
Die in diesen Formeln vorkommenden Resolventen der Kreistheilung $(r^{m_2},\eta)$ sind Kreistheilungszahlen in einem höheren Körper. Ihre $m$-ten Potenzen sind aber im Körper $\Omega_m$ enthaltene Zahlen, die durch die Substitution $(r,r^n)$ in
\[
  (r^{nm_2},\eta)^m
\]
übergehen. Ebenso sind die Verbindungen
\begin{equation}
  (r^{nm_2},\eta)^{-1}(r^{m_2},\eta)^n
\tag{13}
\end{equation}
im Körper $\Omega_m$ enthalten [\S~174, (9), (11)]. Diese Resolventen sind specielle Fälle der allgemeinen Resolvente $(r,x_0)$. Aus (12) erhalten wir, wenn wir mit $\varphi_n,\eps_n$ die conjugirten Functionale zu $\varphi$ und $\eps$ bezeichnen
\begin{equation}
  \omega_n=\eps_n\varphi_n^m\prod_p(r^{m_2 n},\eta)^m.
\tag{14}
\end{equation}

Diese Formel lehrt uns die wichtigsten Eigenschaften der Functionale $\varphi_n$ und $\eps_n$ kennen, zunächst:

\textbf{7. Die Functionale $\varphi_n^m$ sind mit Zahlen des Körpers $\Omega_m$ associirt, gehören also der Hauptclasse an (\S~152).}

Nach 1. ist $\omega_n^{-1}\omega_1^n=\alpha^m$ die $m$-te Potenz einer Zahl $\Omega_m$. Da auch die Verbindungen (13) Zahlen in $\Omega_m$ sind, so können wir
\[
  \prod_p(r^{nm_2},\eta)^{-1}(r^{m_2},\eta)^n=\frac{\alpha}{\beta}
\]
setzen, und erhalten aus (14)
\[
  \eps_n^{-1}\eps_1^n(\varphi_n^{-1}\varphi_1^n)^m=\beta^m,
\]
worin $\beta$ eine Zahl in $\Omega_m$ ist.

Daraus geht hervor, dass das Product
\[
  \beta\varphi_n\varphi_1^{-n}=\eps
\]
ein Einheitsfunctional des Körpers $\Omega_m$ ist, und daraus ergeben sich für die Functionale $\varphi_n$ und die Einheitsfunctionale $\eps_n$ die folgenden beiden Sätze:

\textbf{8. Die conjugirten Functionale $\varphi_n$ haben die Eigenschaft, dass alle Producte}
\[
  \varphi_n\varphi_1^{-n}
\]
\textbf{der Hauptclasse angehören.}

\textbf{9. Die Einheitsfunctionale $\eps_n$ haben die Eigenschaft, dass die Producte}
\[
  \eps_n\eps_1^{-n}=\eps^m
\]
\textbf{$m$-te Potenzen von Einheitsfunctionalen sind.}


\srcpage{659}
Es kommt nun für den Beweis des Haupttheorems I. alles auf den Beweis der beiden Sätze an, die aus den Eigenschaften der Functionale $\varphi$ und $\eps$ zu folgern sind:

\textbf{A. Die Functionale $\varphi_n$ gehören der Hauptclasse an.}

Können wir dieses Lemma voraussetzen, so ist $\varphi_n$ mit einer Zahl $\alpha_n$ associirt, und wir können die Formel (14) mit etwas veränderter Bedeutung von $\eps_n$ so darstellen:
\begin{equation}
  \omega_n=\eps_n\alpha_n^m\prod_p(r^{m_2 n},\eta)^m.
\tag{15}
\end{equation}
In dieser Formel ist aber $\eps_n$ eine numerische Einheit des Körpers $\Omega_m$, die dem Satze 9. genügt.

Können wir also aus dem Satze 9. noch das zweite Lemma beweisen:

\textbf{B. Eine numerische Einheit $\eps_n$ des Körpers $\Omega_m$, die dem Satze 9. genügt, ist die $m$-te Potenz einer Einheit in $\Omega_m$, multiplicirt mit einer Potenz von $r$;}

so können wir in (15) die $m$-te Wurzel ziehen, und erhalten, wenn mit $\rho$ eine Einheitswurzel von möglicherweise höherem Grade als $m$ und mit $\alpha$ eine Zahl des Körpers $\Omega_m$ bezeichnet wird:
\begin{equation}
  (r,x_0)=\rho\alpha\prod_p(r^{m_2},\eta),
\tag{16}
\end{equation}
wo nun rechts lauter Kreistheilungszahlen stehen und wodurch daher das Theorem I. erwiesen ist.

\srcpage{659}
\section*{\S\,182. Beweis des ersten Hülfssatzes für ein ungerades $m$.}

Wir haben im vorigen Paragraphen den Beweis des Theorems I. von zwei Hülfssätzen abhängig gemacht, deren Beweis nun ferner zu suchen ist.

Wir formuliren den ersten dieser Hülfssätze so:

\begin{enumerate}[label=\alph*),leftmargin=2.2em]
\item Ist $\varphi_n$ ein System von Null verschiedener conjugirter Functionale des Körpers $\Omega_m$, von denen bekannt ist, dass
\[
\varphi_n^m\quad\text{und}\quad\varphi_n^{-1}\varphi_1^n
\]
der Hauptclasse angehören, so gehören die Functionale $\varphi$ selbst der Hauptclasse an.
\end{enumerate}

Der Satz wird bewiesen sein, wenn von irgend einer Potenz $\varphi^k$ von $\varphi$, deren Exponent nicht durch $q$ theilbar ist, bewiesen werden kann, dass sie der Hauptclasse angehört. Denn sind $\alpha,\beta$ Zahlen in $\Omega_m$ und ist
\[
\varphi^m=\varepsilon'\alpha,
\qquad
\varphi^k=\varepsilon''\beta,
\]
so können wir die ganzen rationalen Zahlen $x,y$ so bestimmen, dass
\[
mx+ky=1
\]
wird, und dann wird
\[
\varphi=\varphi^{mx}\varphi^{ky}=\varepsilon^{\prime\prime\prime}\alpha^x\beta^y,
\]
also ist auch, wenn $\varepsilon',\varepsilon'',\varepsilon^{\prime\prime\prime}$ Einheiten sind, $\varphi$ selbst mit einer Zahl associirt.

Nach dem heutigen Stande der Sache ist der Beweis für ein gerades $m$ auf einem ganz anderen Wege zu führen, als für ein ungerades, und wir beschäftigen uns also zunächst mit dem leichteren Falle des ungeraden $m$\footnote{Vgl. übrigens den neuen Beweis von Hilbert, bei dem der Unterschied zwischen geradem und ungeradem $m$ weit mehr zurücktritt (S. 642).}.

Die in a) über $\varphi_n$ gemachten Voraussetzungen können wir wenn $\alpha,\beta$ wieder irgend welche Zahlen des Körpers $\Omega_m$ und $\varepsilon$ Einheitsfunctionale bedeuten, durch die beiden Formeln ausdrücken:
\begin{align}
\varphi_n^m&=\varepsilon\beta, \tag{1}\\
\varphi_n&=\varepsilon\alpha\varphi_1^n. \tag{2}
\end{align}
Wir werden nun die Aufgabe schrittweise lösen.

Ist zunächst, wie früher, $\sigma$ der in $q$ enthaltene Primfactor, der eine in $\Omega_m$ existirende Zahl ist, und $\sigma^k$ die in $\varphi$ enthaltene Potenz von $\sigma$, so setzen wir
\begin{equation}
\varphi=\sigma^k\psi,
\tag{3}
\end{equation}
worin $\psi$ ein zu $q$ theilerfremdes Functional bedeutet, das denselben Bedingungen (1), (2) wie $\varphi$ genügt, und wenn eines von beiden in die Hauptclasse gehört, so gilt das auch vom anderen. Der Satz a) braucht also nur noch unter der Voraussetzung bewiesen zu werden, dass $\varphi$ theilerfremd zu $q$ ist.

Es sei nun $p$ eine von $q$ verschiedene Primzahl, die in Primfactoren zerlegt den Ausdruck hat:
\begin{equation}
p=\prod_\xi \frp_\xi,
\tag{4}
\end{equation}
worin $\xi$ die in \S\,172 näher bestimmte Zahlenreihe durchläuft. Es möge $\frp_\xi^{k_\xi}$ die in $\varphi$ enthaltene Potenz von $\frp_\xi$ sein, so dass, wenn wir
\[
\varphi=\chi\prod_\xi \frp_\xi^{k_\xi}
\]
setzen, $\chi$ theilerfremd zu $p$ und mit keinen anderen Primzahlen verwandt ist (s. den Schluss von \S\,173), als $\varphi$ selbst. Daraus ergiebt sich
\begin{equation}
\varphi_n=\chi_n\prod_\xi \frp_{n\xi}^{k_\xi},
\tag{5}
\end{equation}
und hierin lassen wir $n$ abermals die Reihe der Zahlen $\xi$ durchlaufen. Setzen wir noch
\[
\psi=\prod_\xi \chi_\xi,
\]
so ist auch $\psi$ theilerfremd zu $p$ und mit keinen anderen Primzahlen verwandt als $\varphi$, und wir erhalten aus (5) mit Benutzung von (4):
\begin{equation}
\prod_\xi\varphi_\xi=\psi p^{\sum k_\xi},
\tag{6}
\end{equation}
und daraus ergiebt sich, dass das Functional $\psi$ denselben beiden Bedingungen (1), (2) genügt, wie $\varphi$.

Es zeigt sich ferner, dass, wenn die $\varphi_\xi$ in der Hauptclasse liegen, auch $\psi$ in der Hauptclasse enthalten ist.

Wenn aber $p-1$ nicht durch $q$ theilbar ist, so können wir auch das Umgekehrte schliessen. Denn nach (2) ist $\varphi_\xi$ äquivalent mit $\varphi^\xi$ und folglich ist:
\[
\prod_\xi\varphi_\xi\ \text{äquivalent}\ \varphi^{\sum\xi},
\qquad
\text{äquivalent}\ \psi,
\]
und wenn $\psi$ mit einer Zahl associirt ist, so gilt dasselbe von $\varphi^{\sum\xi}$. Nach dem Satze \S\,172 ist aber, wenn $p-1$ nicht durch $q$ theilbar ist, $\sum\xi$ auch nicht durch $q$ theilbar, und folglich ist auch $\varphi$ selbst in der Hauptclasse enthalten.

Diese Betrachtung können wir nun, wenn $\psi$ noch mit weiteren Primzahlen, die nach dem Modul $q$ nicht mit 1 congruent sind, verwandt ist, wiederholen, und kommen so zu dem Satze:
\begin{center}
Der Satz a) ist allgemein bewiesen, wenn er unter der Voraussetzung bewiesen werden kann, dass $\varphi$ nur mit solchen Primzahlen verwandt ist, die nach dem Modul $q$ mit 1 congruent sind.
\end{center}

Machen wir diese Voraussetzung über die Functionale $\varphi_n$, so können wir das Kummer'sche Theorem in der Fassung des \S\,174, 3. anwenden.

Wir bezeichnen in der Folge mit $\varepsilon$ Einheitsfunctionale, auf deren genauere Kenntniss nichts ankommt, und setzen
\begin{equation}
\Phi_n=\prod_t \varphi_{nt'}^{t}=\varepsilon\vartheta_n^m,
\tag{7}
\end{equation}
wenn $t$ die durch $q$ nicht theilbaren Reste von $m$, die kleiner als $m$ sind, durchläuft, und $t'$ durch $tt'\equiv1\pmod m$ bestimmt ist. Dann ist $\vartheta_n$ eine Kreistheilungszahl in einem höheren Körper, deren $m$te Potenz in $\Omega_m$ enthalten ist und die der zweiten Bedingung genügt, dass
\begin{equation}
\vartheta_n^{-1}\vartheta_1^n=\alpha
\tag{8}
\end{equation}
eine Zahl des Körpers $\Omega_m$ ist.

Wir führen nun ein vielfach gebrauchtes Zeichen ein, indem wir unter $E(x)$ die grösste ganze Zahl verstehen, die in der reellen Zahl $x$ enthalten ist, und unter $(x)$ den Rest, der stets ein echter Bruch ist, und, wenn $x$ selbst eine ganze Zahl sein sollte, gleich Null zu setzen ist. Dann ist
\begin{equation}
x=E(x)+(x).
\tag{9}
\end{equation}
Ist $x$ eine ganze Zahl, so ist nach dieser Bezeichnungsweise
\[
m\left(\frac xm\right)
\]
auch eine ganze Zahl, und zwar der kleinste positive Rest der Theilung von $x$ durch $m$.

In der Formel (7) ist der Index der Functionale $\varphi$ nur nach dem Modul $m$ bestimmt. Der Exponent aber ist auf seinen kleinsten Rest nach dem Modul $m$ reducirt. Ersetzen wir also in (7) den Index $t'$ durch $n't'$, so muss $t$ durch den Rest der Division von $nt$ durch $m$, d. h. [nach (9)] durch
\[
m\left(\frac{nt}{m}\right)=nt-mE\left(\frac{nt}{m}\right)
\]
ersetzt werden, und wir können setzen
\[
\Phi_n=\prod_t \varphi_{t'}^{nt-mE(\frac{nt}{m})}=\varepsilon\vartheta_n^m,
\]
andererseits ist
\[
\Phi_1=\prod_t \varphi_{t'}^t=\varepsilon\vartheta_1^m,
\]
also nach (8)
\begin{equation}
\Phi_1^n\Phi_n^{-1}=\prod_t \varphi_{t'}^{mE(\frac{nt}{m})}=\varepsilon\alpha^m.
\tag{10}
\end{equation}
Hiernach ist der Quotient
\[
\left(\frac{\prod_t \varphi_{t'}^{E(\frac{nt}{m})}}{\alpha}\right)^m
\]
ein Einheitsfunctional, und da er zugleich die $m$te Potenz eines Functionals in $\Omega_m$ ist, so ist die Basis dieser Potenz auch eine Einheit, d. h. es ist
\[
\prod_t \varphi_{t'}^{E(\frac{nt}{m})}=\varepsilon\alpha,
\]
und es gehört dies Product der Hauptclasse an. Nun ist aber nach unserer Voraussetzung
\[
\varphi_{t'}\ \text{äquivalent mit}\ \varphi^{t'},
\]
und daher ist
\[
\prod_t \varphi_{t'}^{E(\frac{nt}{m})}\ \text{äquivalent mit}\ \varphi^{\sum t'E(\frac{nt}{m})},
\]
und dies gilt für jedes beliebige durch $q$ nicht theilbare $n$.

Können wir also nachweisen, dass sich $n$ so bestimmen lässt, dass auch die Summe
\[
S_n=\sum_t t'E\left(\frac{nt}{m}\right)
\]
durch $q$ nicht theilbar ist, so haben wir das Theorem a) bewiesen.

Wir wollen zunächst eine Umformung der Summe $S_n$ vornehmen, durch die die Frage auf den weit einfacheren Fall zurückgeführt wird, in dem $m$ eine Primzahl ist.

Wir setzen, wenn $m$ eine höhere als die erste Potenz von $q$ ist,
\[
m=qm',\qquad t=t_1+qt_2,\qquad
\begin{array}{l}
t_1=1,2,\ldots,q-1,\\
t_2=0,1,\ldots,m'-1.
\end{array}
\]
Darin ist $m'$ noch eine Potenz von $q$, der Summationsbuchstabe $t_1$ durchläuft die Zahlenreihe $1,2,\ldots,q-1$, und $t_2$ die Zahlenreihe $0,1,2,\ldots,m'-1$. Wir bestimmen $t_1'$ aus der Congruenz
\[
t_1t_1'\equiv1\pmod q
\]
und erhalten
\[
t'\equiv t_1'\pmod q.
\]
Es ist dann
\begin{equation}
S_n=\sum_t t'E\left(\frac{tn}{m}\right)
\equiv
\sum_{t_1}t_1'\sum_{t_2}E\left(\frac{nt_2}{m'}+\frac{nt_1}{m}\right)
\pmod q.
\tag{11}
\end{equation}
Wir nehmen jetzt $n<q$ an, also auch $nt_1<m$, so dass $nt_1:m$ ein echter Bruch ist. Dann ist
\[
\begin{array}{ll}
\text{a)}& E\left(\dfrac{nt_2}{m'}+\dfrac{nt_1}{m}\right)=E\left(\dfrac{nt_2}{m'}\right),\\[1.2em]
\text{b)}& \text{oder}=E\left(\dfrac{nt_2}{m'}\right)+1,
\end{array}
\]
und zwar tritt der Fall b) nur dann ein, wenn zwischen
\[
\frac{nt_2}{m'}\quad\text{und}\quad \frac{nt_2}{m'}+\frac{nt_1}{m}
\]
eine ganze Zahl liegt. Dies ist aber nur dann der Fall, wenn der Unterschied zwischen $nt_2:m'$ und der nächst grösseren ganzen Zahl kleiner als $nt_1:m$ ist, also wenn
\begin{equation}
E\left(\frac{nt_2}{m'}\right)+1-\frac{nt_2}{m'}<\frac{nt_1}{m}.
\tag{12}
\end{equation}
Lassen wir in dem Ausdrucke auf der linken Seite von (12)
\begin{equation}
E\left(\frac{nt_2}{m'}\right)+1-\frac{nt_2}{m'}
\tag{13}
\end{equation}
$t_2$ alle seine Werthe durchlaufen, so erhalten wir lauter positive, die Einheit nicht übersteigende rationale Brüche mit dem Nenner $m'$, von denen keine zwei einander gleich sind. Denn wären zwei der Werthe, die aus (13) durch Substitution von $t_2',t_2''$ für $t_2$ entstehen, einander gleich, so müsste $n(t_2'-t_2'')/m'$ eine ganze Zahl sein, was, da $n$ relativ prim zu $m'$ und $t_2'-t_2''$ kleiner als $m'$ ist, nicht sein kann. Die Ausdrücke (13) durchlaufen daher in irgend einer Reihenfolge die Zahlenwerthe
\begin{equation}
\frac1{m'},\ \frac2{m'},\ldots,\frac{m'-1}{m'},\ 1,
\tag{14}
\end{equation}
und wenn $a$ einen beliebigen positiven Bruch bedeutet, so ist $E(m'a)$ die Anzahl der Zahlen der Reihe (14), die nicht grösser als $a$ sind. Nach (12) ist daher die Anzahl der Werthe von $t_2$, für die der Fall b) eintritt, gleich $E(nt_1/q)$, und es ergiebt sich
\[
\sum_{t_2}E\left(\frac{nt_2}{m'}+\frac{nt_1}{m}\right)
=
\sum_{t_2}E\left(\frac{nt_2}{m'}\right)+E\left(\frac{nt_1}{q}\right).
\]
Um $S_n$ zu bilden, multipliciren wir diese Gleichung mit $t_1'$ und summiren. Dadurch ergiebt sich
\[
S_n\equiv \sum_{t_1}t_1'\sum_{t_2}E\left(\frac{nt_2}{m'}\right)+\sum_{t_1}t_1'E\left(\frac{nt_1}{q}\right)\pmod q.
\]
Darin ist nun $\sum t_1'=\sum t_1\equiv0\pmod q$, weil die $t_1$ paarweise die Summe $q$ ergeben, und es folgt:
\begin{equation}
S_n\equiv\sum_{t_1}t_1'E\left(\frac{nt_1}{q}\right)\pmod q.
\tag{15}
\end{equation}
Die Summe, die hier auf der rechten Seite steht, ist ebenso gebildet, wie $S_n$ selbst in dem Falle, wo $m$ eine Primzahl ist. Für diesen Fall ergiebt sich aber
\[
\frac{nt_1}{q}=E\left(\frac{nt_1}{q}\right)+\left(\frac{nt_1}{q}\right),
\]
\[
\frac{(q-n)t_1}{q}=E\left(\frac{(q-n)t_1}{q}\right)+\left(\frac{(q-n)t_1}{q}\right),
\]
und daraus durch Addition
\[
t_1=E\left(\frac{nt_1}{q}\right)+E\left(\frac{(q-n)t_1}{q}\right)+\left(\frac{nt_1}{q}\right)+\left(\frac{(q-n)t_1}{q}\right),
\]
und dabei kann die Summe der beiden positiven echten Brüche, die doch eine ganze Zahl sein muss, nur den Werth 1 haben. Es folgt also
\[
E\left(\frac{nt_1}{q}\right)+E\left(\frac{(q-n)t_1}{q}\right)=t_1-1,
\]
und daraus durch Multiplication mit $t_1'$ und Summation über alle Werthe von $t_1$ von 1 bis $q-1$
\[
\sum_{t_1}t_1'E\left(\frac{nt_1}{q}\right)+\sum_{t_1}t_1'E\left(\frac{(q-n)t_1}{q}\right)\equiv -1\pmod q.
\]
Folglich können sicher nicht beide Summen auf der linken Seite durch $q$ theilbar sein. Dann zeigt der Ausdruck (15), dass von den beiden Summen $S_n$ und $S_{q-n}$ gewiss eine durch $q$ untheilbar ist.

Dies aber war zu beweisen.

\srcpage{666}
\section*{\S\,183. Beweis des zweiten Hülfssatzes für ein ungerades $m$.}

Das zweite Lemma, was nach \S\,181 noch zu beweisen ist, ist folgendes:

\begin{enumerate}[label=\arabic*.,leftmargin=2.2em,start=2]
\item Ist $\varepsilon_n$ ein System conjugirter numerischer Einheiten des Körpers $\Omega_m$, das der Bedingung genügt, dass
\begin{equation}
\varepsilon_n\varepsilon_1^{-n}=\mathscr E^m
\tag{1}
\end{equation}
die $m$te Potenz einer Einheit in $\Omega_m$ ist, so sind die $\varepsilon_n$ selbst $m$te Potenzen von Einheiten in $\Omega_m$, multiplicirt mit einer Einheitswurzel der Form $r^{nk}$.
\end{enumerate}

Auch dieser Beweis ist für den Fall eines ungeraden $m$ ganz elementar. Wir betrachten daher hier zunächst diesen Fall.

Nach \S\,178, 1. können wir jede Einheit des Körpers $\Omega_m$ in der Form darstellen:
\begin{equation}
\varepsilon=\pm r^k\mathscr E(r),\qquad
\varepsilon_n=\pm r^{nk}\mathscr E(r^n),
\tag{2}
\end{equation}
worin $\mathscr E(r)$ eine reelle Einheit des Körpers $\Omega_m$ ist. Es genügt also $\mathscr E(r)$ der Bedingung
\begin{equation}
\mathscr E(r)=\mathscr E(r^{-1}),
\tag{3}
\end{equation}
woraus sich nach (2) ergiebt:
\begin{equation}
\varepsilon_1\varepsilon_{-1}=\mathscr E(r)^2;
\tag{4}
\end{equation}
aus (1) aber findet sich, wenn man $n=-1$ setzt,
\[
\varepsilon_1\varepsilon_{-1}=\mathscr E^m,
\]
d. h. die linke Seite von (4) ist die $m$te Potenz einer Einheit in $\Omega_m$, und es ist also auch
\begin{equation}
\mathscr E(r)^2=\mathscr E^m
\tag{5}
\end{equation}
die $m$te Potenz einer solchen Einheit.

Da nun $m$ ungerade ist, so lässt sich die ganze rationale Zahl $x$ so bestimmen, dass
\begin{equation}
2x+m=1
\tag{6}
\end{equation}
wird, und daraus folgt
\[
\mathscr E(r)=\mathscr E(r)^{2x}\mathscr E(r)^m.
\]
Wenn wir also nach (5)
\[
\mathscr E^x\mathscr E(r)=e(r)
\]
setzen, so folgt
\begin{equation}
\mathscr E(r)=e(r)^m.
\tag{7}
\end{equation}
Durch (7) ist aber der Satz 2. bewiesen und damit zugleich für ein ungerades $m$ das ganze Theorem I.

Auch dieser Schluss versagt für ein gerades $m$, weil dann die Gleichung (6) nicht mehr lösbar ist.

\srcpage{667}
\section*{\S\,184. Vorläufiges über den Fall eines geraden $m$.}

Für den Fall, dass $m$ eine Potenz von 2 ist, schlagen wir einen ganz anderen Weg ein, über den hier zunächst einige vorläufige Bemerkungen Platz finden mögen.

Wir haben schon im \S\,182 gesehen, dass das erste Lemma bewiesen ist, wenn wir nachweisen können, dass irgend eine Potenz von $\varphi$, deren Exponent nicht durch $q$ theilbar ist, zur Hauptclasse gehört. Dabei ist von den beiden Eigenschaften des Functionals $\varphi$ nur die erste benutzt, dass die $m$te Potenz von $\varphi$ in der Hauptclasse enthalten ist.

Nun kennen wir nach den allgemeinen Sätzen in \S\,154 in jedem Körper einen Exponenten $h$, nämlich die Anzahl der Idealclassen, für den $\varphi^h$ zur Hauptclasse gehört, wenn $\varphi$ ein beliebiges Functional in diesem Körper ist.

Wenn wir also beweisen können, dass die Classenzahl $h$ im Körper $\Omega_m$, wenn $m$ eine Potenz von 2 ist, eine ungerade Zahl ist, so ist damit ohne alles Weitere das Lemma 1. bewiesen.

Dies soll das Ziel unserer Betrachtungen in den nächsten Abschnitten sein.

In Bezug auf das zweite Lemma liegt die Sache ähnlich. Hier kann man aus den Voraussetzungen in \S\,183, 2. ganz wie oben die Formel \S\,183, (5) herleiten, aus der aber nur zu schliessen ist, dass $\mathscr E(r)$ die $\frac12m$te Potenz einer reellen Einheit $e(r)$ ist.

Nehmen wir also das erste Lemma als bewiesen an, so ergiebt die Formel \S\,181, (14):
\begin{equation}
(r^n,x_0)=\rho_n\sqrt{e(r^n)}\,\alpha_n\prod_p (r^{m_2n},\eta),
\tag{1}
\end{equation}
worin $\rho_n$ eine Einheitswurzel vom Grade $m^2$, und $\alpha_n$ eine Zahl in $\Omega_m$ ist, und es wäre noch zu beweisen, dass $e(r^n)$ das Quadrat einer Einheit in $\Omega_m$ ist.

Da $e(r)$ eine reelle Einheit ist, so ist
\begin{equation}
e(r^n)=e(r^{-n}),
\tag{2}
\end{equation}
und wir wollen noch festsetzen, dass das Vorzeichen der Wurzel so bestimmt sei (was wir bei passender Annahme über $\rho_n$ immer annehmen können), dass
\begin{equation}
\sqrt{e(r^n)}=\sqrt{e(r^{-n})}.
\tag{3}
\end{equation}
Das Product $(r^n,x_0)(r^{-n},x_0)$ ist nach der Voraussetzung eine Zahl des Körpers $\Omega_m$, die sich durch die Substitution $(r,r^{-1})$ nicht ändert, d. h. eine reelle Zahl. Ferner ist nach \S\,174, (5)
\[
(r^{m_2n},\eta)(r^{-m_2n},\eta)=\pm p,
\]
also gleichfalls reell. Ebenso ist $\alpha_n\alpha_{-n}$ reell, und daraus ergiebt sich nach (1), dass
\[
\rho_n\rho_{-n}=\frac{(r^n,x_0)(r^{-n},x_0)}{e(r^n)\alpha_n\alpha_{-n}\prod_p (r^{m_2n},\eta)(r^{-m_2n},\eta)}
\]
eine reelle Zahl ist, die, weil sie eine Einheitswurzel ist, nur $=\pm1$ sein kann.

Nun können wir in der hieraus für $n=1$ sich ergebenden Gleichung
\[
(r,x_0)(r^{-1},x_0)=\pm e(r)\alpha_1\alpha_{-1}\prod_p(\pm p)
\]
jede Substitution $(r,r^n)$ ausführen, woraus hervorgeht, dass das Zeichen von $\rho_n\rho_{-n}$ von $n$ unabhängig ist, dass also
\begin{equation}
\rho_n\rho_{-n}=\rho_1\rho_{-1}=\pm1
\tag{4}
\end{equation}
ist. Nun leiten wir aus (1) weiter her:
\begin{equation}
\rho_n\rho_1^{-1}\sqrt{e(r^n)}\sqrt{e(r)}^{-n}=\frac{(r^n,x_0)(r,x_0)^{-n}}{\alpha_n\alpha_1^{-1}\prod_p(r^{m_2n},\eta)(r^{m_2},\eta)^{-n}},
\tag{5}
\end{equation}
und die rechte Seite ist eine Zahl des Körpers $\Omega_m$. Die linke Seite zeigt aber, dass es eine Einheit ist, und wir können demnach, wenn $\mathscr E(r)$ eine reelle Einheit bedeutet,
\begin{equation}
\rho_n\rho_1^{-1}\sqrt{e(r^n)}\sqrt{e(r)}^{-n}=r^k\mathscr E(r)
\tag{6}
\end{equation}
setzen. Substituiren wir diesen Werth in die Formel (5) und machen die Substitution $(r,r^{-1})$, so ergiebt sich
\begin{equation}
\rho_{-n}\rho_{-1}^{-1}\sqrt{e(r^n)}\sqrt{e(r)}^{-n}=r^{-k}\mathscr E(r),
\tag{7}
\end{equation}
und durch Multiplication von (6) und (7) mit Rücksicht auf (4)
\begin{equation}
e(r^n)e(r)^{-n}=\mathscr E(r)^2.
\tag{8}
\end{equation}

In unseren Betrachtungen über die Classenzahl wird sich noch der Satz ergeben:
\begin{center}
Eine Einheit $e(r)$ des reellen Körpers $H_m$, die mit allen ihren Conjugirten positiv ist, ist das Quadrat einer Einheit im Körper $H_m$.
\end{center}
Die Formel (8) zeigt aber, dass der Einheit $\pm e(r)$ diese Eigenschaft zukommt, und dass daher $e(r)$ (da auch $-1=i^2$ das Quadrat einer Einheit ist) das Quadrat einer Einheit des Körpers $\Omega_m$ ist.

Damit ist die Frage auf den Beweis der beiden erwähnten Sätze zurückgeführt, der sich als Schlussstein einer langen, aber an interessanten Beziehungen äusserst reichen Kette von Betrachtungen ergiebt, denen die nächsten Abschnitte gewidmet sein sollen.

\srcpage{670}
\section*{Dreiundzwanzigster Abschnitt. Classenzahl.}

\section*{\S~185. Der Dirichlet'sche Satz über die Einheiten.}

Die Untersuchungen über die Anzahl der Idealclassen in einem algebraischen Körper, die uns nun die nothwendigen Ergänzungen der Beweise des vorigen Abschnittes geben sollen, bedienen sich der Methoden, die Dirichlet in die Zahlentheorie eingeführt hat.\footnote{Dirichlet, \textit{Recherches sur diverses applications de l'analyse infinitésimale à la théorie des nombres}. Crelle's Journal, Bd.~19, 21. Dirichlet's Werke, Bd.~I (1839, 1840). (Hierher gehören übrigens auch die nachgelassenen Abhandlungen von Gauss, ``De nexu inter multitudinem classium etc.'' Gauss' Werke, Bd.~II.) Die Anwendung dieser Principien auf die Kreistheilungszahlen hat zuerst Kummer gemacht [Crelle's Journal, Bd.~35, 40 (1847, 1850); Liouville's Journal, Bd.~16 (1851)]. Die allgemeine Theorie der Einheiten ist gleichfalls von Dirichlet begründet [Dirichlet's Werke, Bd.~I, S.~622, 633, 639 (1840, 1842, 1846)]. Die allgemeine Theorie der Einheiten und die Bestimmung der Classenzahlen hat Dedekind gegeben: Dirichlet-Dedekind, Vorlesungen über Zahlentheorie, 4.~Auflage, \S~183 f. Minkowski, Geometrie der Zahlen.}

Sie sind nicht mehr rein algebraisch, insofern dabei auch transcendente Functionen und Zahlen, wie der natürliche Logarithmus und die Zahl $\pi$ vorkommen. Auch wird von Grenzübergängen, wie in der Integralrechnung, Gebrauch gemacht, aus denen wir schon in \S~156 ein schönes Resultat gezogen haben.

Die Untersuchungen, von denen wir zunächst zu handeln haben, sind ganz allgemein, d.~h. sie gelten für jeden algebraischen Zahlkörper, und es bietet keinen Vortheil der Einfachheit, sie auf specielle Körper, wie etwa die Kreistheilungskörper, zu beschränken.

\srcpage{671}
Es sei also $\Omega$ ein Körper $n$ten Grades, und
\begin{equation}
\Omega_1,\Omega_2,\ldots,\Omega_n
\tag{1}
\end{equation}
die conjugirten Körper. Die irreducible Gleichung $n$ten Grades, deren Wurzeln uns diese conjugirten Körper bestimmen, wird im Allgemeinen sowohl reelle als imaginäre Wurzeln haben, von denen aber die letzteren immer paarweise vorkommen, so dass zu jedem imaginären Körper $\Omega^{(i)}$ ein anderer gehört, dessen Zahlen mit denen von $\Omega^{(i)}$ conjugirt imaginär sind, der also aus $\Omega^{(i)}$ durch die Substitution $(i,-i)$ hervorgeht.

Solche imaginäre Paare fassen wir zu einer Einheit zusammen und bezeichnen die Anzahl der reellen Körper und der Paare imaginärer Körper der Reihe (1) mit $\nu$. Sind alle conjugirten Körper reell, so ist $n=\nu$; sind sie alle imaginär, so ist $n=2\nu$.

Ausser im Falle des rationalen Körpers ist $\nu$ nur noch in dem Falle eines imaginären quadratischen Körpers $=1$. Sonst ist $\nu$ immer grösser als 1.

Die Anzahl der imaginären Paare ist $n-\nu$, und die Anzahl der reellen Körper $2\nu-n$.

In der Determinante, durch die in \S~145 die Quadratwurzel aus der Grundzahl, $\sqrt{\Delta}$, definirt ist, vertauschen sich durch die Substitution $(i,-i)$ je zwei conjugirt imaginäre Reihen, und $\sqrt{\Delta}$ wechselt also sein Zeichen dann, wenn diese Zahl ungerade ist, sonst nicht. Im ersten Falle ist $\sqrt{\Delta}$ imaginär, im zweiten reell. Daraus ergiebt sich also, dass die Grundzahl $\Delta$ des Körpers $\Omega$ positiv oder negativ ist, je nachdem $n-\nu$ gerade oder ungerade ist.

Behalten wir von zwei conjugirt imaginären Körpern nur den einen bei, so ergiebt sich die Reihe der conjugirten Körper
\[
\Omega_1,\Omega_2,\ldots,\Omega_\nu,
\]
denen wir ein Zeichensystem
\[
\delta_1,\delta_2,\ldots,\delta_\nu
\]
mit der Bestimmung zuordnen, dass $\delta_s=1$ sein soll, wenn $\Omega_s$ reell, und $=2$, wenn $\Omega_s$ imaginär ist, also $\sum \delta_s=n$.

Es möge $\eta$ eine von Null verschiedene Zahl des Körpers $\Omega$ bedeuten und
\[
\eta_1,\eta_2,\ldots,\eta_n
\]
die conjugirten Zahlen. Diesem Zahlensysteme ordnen wir ein anderes Zahlensystem zu
\[
\lambda_1,\lambda_2,\ldots,\lambda_\nu,
\]
\srcpage{672}
das wir durch die Gleichung
\begin{equation}
\lambda_s=\delta_s\log |\eta_s|
\tag{2}
\end{equation}
definiren, worin $|\eta_s|$ den absoluten Werth von $\eta_s$ und $\log|\eta_s|$ den reellen natürlichen Logarithmus der positiven Zahl $|\eta_s|$ bedeutet. Ist $\Omega_s$ reell und das Zeichen $\pm$ so gewählt, dass $\pm\eta_s$ positiv ist, so ist
\[
\lambda_s=\log(\pm\eta_s).
\]
Ist aber $\eta_s,\eta_s'$ ein imaginäres Paar, so ist
\[
\lambda_s=\log\eta_s\eta_s',
\]
und zu $\eta_s$ und $\eta_s'$ gehört dasselbe $\lambda_s$, so dass die Anzahl der verschiedenen $\lambda_s$ gleich $\nu$ ist. Aus dieser Bestimmung ergiebt sich noch
\begin{equation}
\lambda_1+\lambda_2+\cdots+\lambda_\nu=\log N_a(\eta),
\tag{3}
\end{equation}
wenn, wie früher, $N_a$ die absolute Norm bedeutet.

Die Zahlen $\lambda_1,\lambda_2,\ldots,\lambda_\nu$ wollen wir die conjugirten Logarithmen der Zahl $\eta$ nennen.

Wenn $\eta$ eine ganze Zahl des Körpers $\Omega$ ist, so ist die absolute Norm eine natürliche Zahl, die nur dann $=1$ ist, wenn $\eta$ eine Einheit ist, und daraus folgt nach (3):

\textbf{1. Die Summe der conjugirten Logarithmen einer ganzen Zahl $\eta$ ist positiv, und nur dann gleich Null, wenn $\eta$ eine Einheit ist.}

Bedeutet $\omega_1,\omega_2,\ldots,\omega_n$ eine Basis der ganzen Zahlen in $\Omega$ (eine Minimalbasis von $\Omega$, \S~145), und $\omega_{1,s},\omega_{2,s},\ldots,\omega_{n,s}$ für $s=1,2,\ldots,n$ die conjugirten Zahlen, so lassen sich alle ganzen Zahlen des Körpers $\Omega_s$ in der Form darstellen:
\begin{equation}
\eta_s=x_1\omega_{1,s}+x_2\omega_{2,s}+\cdots+x_n\omega_{n,s},
\tag{4}
\end{equation}
mit ganzen rationalen Coëfficienten $x_1,x_2,\ldots,x_n$. Die Determinante dieses Systems
\[
\sum \pm \omega_{1,1}\omega_{2,2}\cdots\omega_{n,n}=\sqrt{\Delta}
\]
ist die Quadratwurzel aus der Discriminante $\Delta$ des Körpers, und demnach immer von Null verschieden. Demnach lässt sich das System (4) nach den Unbekannten $x_1,x_2,\ldots,x_n$ auflösen, und wenn wir nun festsetzen, dass die absoluten Werthe der Zahlen $\eta_s$ nicht über eine gewisse Grenze hinausgehen sollen, so ergeben sich aus diesen Auflösungen Grenzen für die ganzen rationalen Zahlen $x_i$.

\srcpage{673}
Diese einfache Bemerkung liefert uns den folgenden wichtigen Satz:

\textbf{2. Es gibt in einem algebraischen Körper $\Omega$ nur eine endliche Anzahl ganzer Zahlen $\eta$ von der Beschaffenheit, dass die absoluten Werthe der conjugirten Zahlen $\eta_s$ unter einer gegebenen Grenze liegen.}

Wir machen jetzt von dem allgemeinen Satze \S~155, (25) Gebrauch, nach dem, wenn $f(x_1,x_2,\ldots,x_n)$ irgend eine positive quadratische Form von $n$ Variablen mit der Determinante $D$ ist, die Variablen $x_i$ sich als ganze rationale Zahlen, nicht alle verschwindend, so bestimmen lassen, dass
\[
f(x_1,x_2,\ldots,x_n)< n\sqrt[n]{0{,}404D},
\]
d.~h. kleiner als eine von der Determinante $D$ allein abhängige Zahl wird.

Diesen Satz wenden wir, wie im \S~156, auf die Form an:
\[
f(x_1,x_2,\ldots,x_n)=\frac{|\eta_1|^2}{c_1^2}+\frac{|\eta_2|^2}{c_2^2}+\cdots+\frac{|\eta_n|^2}{c_n^2},
\]
worin die $\eta_s$ die Linearformen (4) sind, und im Falle eines reellen $\eta_s$
\[
|\eta_s|^2=(\eta_s)^2,
\]
im Falle eines imaginären Paares $\eta_s,\eta_s'$
\[
|\eta_s|^2=|\eta_s'|^2=\eta_s\eta_s'.
\]
Im letzteren Falle ist $\eta_s\eta_s'$ die Summe zweier reeller Quadrate, und demnach ist, wenn die $c_s$ reell angenommen werden, $f$ eine positive Form. Wir wollen über die bis jetzt willkürlichen reellen Grössen $c_s$ noch festsetzen, dass, wenn $\eta_s,\eta_s'$ ein conjugirt imaginäres Paar bilden,
\begin{equation}
c_s=c_s'
\tag{5}
\end{equation}
sein soll; ausserdem wollen wir die $c_s$ noch der Bedingung unterwerfen:
\begin{equation}
c_1c_2\cdots c_n=1.
\tag{6}
\end{equation}
Die Determinante der Function $f$ bilden wir, wie im \S~156, dadurch, dass wir sie zunächst als Form der Variablen $\eta_1,\eta_2,\ldots,\eta_n$ auffassen, und die Determinante dieser Form ist wegen (6) gleich $\pm1$. Die Determinante von $f$ in Bezug auf die Variablen $x_i$ ist also (Bd.~I, \S~56), vom Vorzeichen abgesehen, gleich dem
\srcpage{674}
Quadrate der Determinante der Linearformen $\eta_s$, d.~h. der Grundzahl des Körpers $\Omega$.

Daraus folgt, dass man eine von den $c_s$ unabhängige, nur durch den Körper $\Omega$ bestimmte Zahl $A$ finden kann, von der Art, dass, was auch die $c_s$ sein mögen, $f(x_1,x_2,\ldots,x_n)$ für ganzzahlige nicht verschwindende $x$ unter $A$ herunterinkt. Da $f$ die Summe von positiven Summanden ist, so gilt dasselbe von jedem dieser Summanden, und damit ist der folgende Satz bewiesen:

\textbf{3. Ist $c_s$ ein beliebiges, den Bedingungen (5), (6) genügendes System reeller positiver Zahlen, so gibt es eine durch den Körper $\Omega$ allein bestimmte reelle Zahl $A$ von der Art, dass man immer eine ganze Zahl $\eta$ des Körpers $\Omega$ bestimmen kann, die mit ihren conjugirten Zahlen $\eta_s$ den Bedingungen genügt}
\begin{equation}
|\eta_s|<c_sA.
\tag{7}
\end{equation}

Beiläufig folgt hieraus, dass man in jedem algebraischen Körper $\Omega$, ausgenommen dem rationalen und dem imaginären quadratischen, von Null verschiedene ganze Zahlen finden kann, deren absoluter Werth unter jede gegebene Grenze herunterinkt.

Wir wollen den Ausdruck des Satzes 3 durch Einführung der conjugirten Logarithmen von $\eta$ etwas umformen. Wir setzen $\log A=k$, so dass $k$ eine durch $\Omega$ bestimmte reelle Zahl ist, ferner
\begin{equation}
c_s=e^{\frac{\gamma_su}{\delta_s}},
\tag{8}
\end{equation}
worin $u$ eine beliebige reelle Grösse sein kann, und die Zahlen $\gamma_s$ wegen (5) und (6) der Bedingung genügen:
\begin{equation}
\gamma_1+\gamma_2+\cdots+\gamma_\nu=\sum_{s=1}^{\nu}\gamma_s=0.
\tag{9}
\end{equation}
Da überdies
\[
|\eta_s|=e^{\frac{\lambda_s}{\delta_s}}
\]
ist, so folgt aus (7) und (8):
\begin{equation}
\lambda_s-\gamma_su<k\delta_s.
\tag{10}
\end{equation}
Die Summe der $\nu$ Ausdrücke auf der linken Seite dieser Ungleichung ist wegen (9) gleich $\sum\lambda_s$, also nach dem Satze 1 nicht negativ, und wir finden
\[
0\leq \sum_s(\lambda_s-\gamma_su)<kn.
\]

\srcpage{675}
Daraus ergiebt sich wegen (10), dass ein Glied dieser Summe, $\lambda_s-\gamma_su$, nicht kleiner sein kann, als $-(n-\delta_s)k$, weil sonst die Gesammtsumme negativ wäre, und wenn wir der Einfachheit halber die untere Grenze noch etwas kleiner nehmen
\begin{equation}
-nk\delta_s<\lambda_s-\gamma_su<k\delta_s.
\tag{11}
\end{equation}

Damit ist bewiesen, dass es für jedes gegebene System der Zahlen $u,\gamma_s$, wenn nur die Bedingung (9) erfüllt ist, eine ganze Zahl $\eta$ des Körpers $\Omega$ gibt, die mit ihren conjugirten Zahlen den Grenzbedingungen (11) genügt.

Es sei nun ferner $g_s$ ein System reeller Grössen, von dem wir nur verlangen, dass
\[
\gamma_1g_1+\gamma_2g_2+\cdots+\gamma_\nu g_\nu=\alpha
\]
von Null verschieden sei. Damit ist [nach (9)] ausgeschlossen, dass alle $g_s$ einander gleich sind; wenn sie aber das nicht sind, so werden sich zu jedem gegebenen System der $g_s$ unendlich viele Systeme $\gamma_s$ bestimmen lassen, die der Forderung (9) genügen.

Wird nun $k\sum\delta_sg_s=g$ gesetzt, so ergiebt sich aus (11) durch Multiplication mit $g_s$ und Summation
\begin{equation}
-ng+\alpha u<\sum g_s\lambda_s<g+\alpha u.
\tag{12}
\end{equation}

Nach dieser Grenzbedingung bestimmen wir nun, bei festgehaltenen $g_s$ und $\alpha$, eine Reihe von ganzen Zahlen
\begin{equation}
\eta,\eta',\eta'',\eta''',\ldots,
\tag{13}
\end{equation}
indem wir für $u$ eine Reihe von Annahmen machen:
\[
u,u',u'',u''',\ldots,
\]
die folgendermaassen näher bestimmt sind:

Wir wählen $u$ zunächst so, dass
\[
-ng+\alpha u=a,
\qquad
 g+\alpha u=a'
\]
positiv werden; dann setzen wir
\[
\delta=\frac{a'-a}{\alpha}=\frac{(n+1)g}{\alpha},
\]
und setzen
\[
u'=u+\delta,
\quad
u''=u'+\delta,
\quad
u'''=u''+\delta,
\quad\ldots,
\]
\[
a+\alpha\delta=a',
\quad a'+\alpha\delta=a'',
\quad a''+\alpha\delta=a''',
\quad\ldots,
\]
\srcpage{676}
und erhalten so aus (12), wenn $\lambda_s',\lambda_s'',\ldots$ die conjugirten Logarithmen von $\eta',\eta'',\ldots$ sind:
\begin{equation}
\begin{array}{rcl}
a&<&\sum g_s\lambda_s<a',\\
a'&<&\sum g_s\lambda_s'<a'',\\
a''&<&\sum g_s\lambda_s''<a''',\\
&&\cdots\cdots,
\end{array}
\tag{14}
\end{equation}
und die Reihe dieser Ungleichungen lässt sich unbegrenzt fortsetzen.

Alle diese Zahlen $\eta,\eta',\eta'',\ldots$ haben überdies nach (6) und (7) die Eigenschaft, dass ihre absoluten Normen $N,N',N'',\ldots$ unter einer bestimmten endlichen Grenze $e^{nk}$ bleiben.

Die Anzahl der möglichen Werthe von $N,N',N'',\ldots$ ist also endlich, nämlich gewiss nicht grösser, als die grösste in $e^{nk}$ enthaltene ganze Zahl. Ebenso ist die Anzahl aller möglichen Reste der Zahlen $x_1,x_2,\ldots,x_n$ [in (4)] nach allen Moduln $N,N',N'',\ldots$ nur eine endliche, und daraus folgt, dass, wenn wir die Reihe der Zahlen (13) nur weit genug fortsetzen, in der Reihe zwei verschiedene Zahlen $\eta^{(h)},\eta^{(k)}$ auftreten müssen, in denen $N^{(h)}=N^{(k)}$ und die Reste von $x_1,x_2,\ldots,x_n$ nach dem Modul $N^{(h)}$ genau dieselben sind.

Aus (14) aber folgt, wenn $h>k$ vorausgesetzt wird:
\begin{equation}
\sum g_s\bigl(\lambda_s^{(h)}-\lambda_s^{(k)}\bigr)>0.
\tag{15}
\end{equation}
Ist $N$ die absolute Norm dieser beiden Zahlen $\eta^{(h)},\eta^{(k)}$, so ist wegen der Uebereinstimmung der Reste der $x$ die Differenz $\eta^{(h)}-\eta^{(k)}$ durch $N$ theilbar. Andererseits ist nach \S~138 (13) die Zahl $N$ durch $\eta^{(h)}$ theilbar, und folglich ist auch $\eta^{(h)}-\eta^{(k)}$ durch $\eta^{(h)}$ theilbar. Daraus ergiebt sich, dass auch $\eta^{(k)}$ durch $\eta^{(h)}$ und ebenso $\eta^{(h)}$ durch $\eta^{(k)}$ theilbar ist. Beide Zahlen sind also associirt, und wenn wir
\[
\frac{\eta^{(h)}}{\eta^{(k)}}=\varepsilon
\]
setzen, so ist $\varepsilon$ eine Einheit des Körpers $\Omega$.

Setzen wir ausserdem, indem wir mit $|\varepsilon_s|$ den absoluten Werth von $\varepsilon_s$ bezeichnen,
\begin{equation}
\delta_s\log|\varepsilon_s|=l_s,
\tag{16}
\end{equation}
so ergiebt sich aus (15):
\begin{equation}
g_1l_1+g_2l_2+\cdots+g_\nu l_\nu>0.
\tag{17}
\end{equation}

\srcpage{677}
Dies beweist nun der Satz:

\textbf{4. Ist $g_1,g_2,\ldots,g_\nu$ ein beliebiges System reeller Zahlen, die nicht alle einander gleich sind, so lässt sich in jedem algebraischen Körper $\Omega$ eine Einheit $\varepsilon$ so bestimmen, dass die Summe}
\[
g_1l_1+g_2l_2+\cdots+g_\nu l_\nu
\]
\textbf{von Null verschieden ist, wenn $l_s$ das System der conjugirten Logarithmen von $\varepsilon$ ist.}

Dieser wichtige Satz, der das Fundament für das Folgende ist, rührt, wenn auch in etwas veränderter Fassung, von Dirichlet her. Diese Formulirung ist von Minkowski gegeben.

\section*{\S~186. Systeme unabhängiger Einheiten und Exponentensysteme der Einheiten.}

Wenn wir in dem zuletzt bewiesenen Satze $g_1=1$, $g_2=0,\ldots,g_\nu=0$ setzen, was, wenn wir von jetzt an $\nu>1$ voraussetzen, gestattet ist, so folgt, dass es in $\Omega$ eine Einheit $\varepsilon$ gibt, für die einer der conjugirten Logarithmen, $l_1$, von Null verschieden ist. Da jede Potenz von $\varepsilon$ dieselbe Eigenschaft hat, so gibt es auch unendlich viele solche Einheiten.

Dieser Satz gestattet nun eine sehr wichtige Verallgemeinerung:

\textbf{5. Ist $s\leq \nu-1$, so lässt sich im Körper $\Omega$ ein System von Einheiten mit den zugehörigen conjugirten Logarithmen}
\begin{equation}
\begin{array}{cccc}
\varepsilon_1: & l_{1,1},&l_{1,2},&\ldots,l_{1,\nu},\\
\varepsilon_2: & l_{2,1},&l_{2,2},&\ldots,l_{2,\nu},\\
&\multicolumn{3}{c}{\cdots\cdots\cdots}\,\\
\varepsilon_s: & l_{s,1},&l_{s,2},&\ldots,l_{s,\nu},
\end{array}
\tag{1}
\end{equation}
\textbf{derart bestimmen, dass eine beliebige der $s$-reihigen Determinanten der Matrix der $l_{h,k}$, etwa}
\[
L_s=\sum \pm l_{1,1}l_{2,2}\cdots l_{s,s}
\]
\textbf{von Null verschieden ist.}

Für $s=1$ ist der ausgesprochene Satz nach der am Eingange gemachten Bemerkung richtig. Wir nehmen also an, er
\srcpage{678}
sei bewiesen für $s-1$ und leiten ihn durch vollständige Induction allgemein her. Dazu ordnen wir die Determinante $L_s$ nach den Elementen der letzten Zeile, und schreiben sie so:
\begin{equation}
L_s=g_1l_{s,1}+g_2l_{s,2}+\cdots+g_sl_{s,s},
\tag{2}
\end{equation}
worin $g_s=L_{s-1}$ ist, und daher nach der gemachten Voraussetzung von Null verschieden angenommen werden kann. Substituiren wir diese Werthe von $g$ in die Formel des Satzes 4, \S~185, indem wir $g_{s+1}=0,\ldots,g_\nu=0$ annehmen, während $g_s$ nicht $=0$ ist, so sind, so lange $s<\nu$ ist, gewiss nicht alle $g_1,g_2,\ldots,g_\nu$ einander gleich, und es ergiebt sich, dass sich $\varepsilon_s$ so annehmen lässt, dass $L_s$ von Null verschieden ist, wie bewiesen werden sollte.

Auf $s=\nu$ ist diese Schlussweise nicht auszudehnen, und die Determinante $L_\nu$ muss auch in der That immer Null sein, weil für jede Einheit $\sum_{1}^{\nu} l_s$ verschwindet.

Wir wollen jetzt für den Fall, dass $s=\nu-1$ ist, die Determinante $L_s$ so bezeichnen:
\begin{equation}
L_{\nu-1}=L(\varepsilon_1,\varepsilon_2,\ldots,\varepsilon_{\nu-1}),
\tag{3}
\end{equation}
und ein System von Einheiten $\varepsilon_1,\varepsilon_2,\ldots,\varepsilon_{\nu-1}$, für welches diese Determinante von Null verschieden ist, ein System unabhängiger Einheiten nennen.

Der absolute Werth $L$ der Determinante $L(\varepsilon_1,\varepsilon_2,\ldots,\varepsilon_{\nu-1})$ heisst der Regulator des Systems $\varepsilon_1,\varepsilon_2,\ldots,\varepsilon_{\nu-1}$.\footnote{Dirichlet-Dedekind, Vorlesungen über Zahlentheorie, 4.~Auflage, \S~183.}

Wenn wir in der Determinante
\begin{equation}
\left|\begin{array}{cccc}
l_{1,1},&l_{1,2},&\ldots,&l_{1,\nu}\\
\cdots&\cdots&\cdots&\cdots\\
l_{\nu-1,1},&l_{\nu-1,2},&\ldots,&l_{\nu-1,\nu}\\
x_1,&x_2,&\ldots,&x_\nu
\end{array}\right|,
\tag{4}
\end{equation}
in der die $x_1,x_2,\ldots,x_\nu$ willkürliche Grössen sind, alle Colonnen zu der letzten addiren, so erhalten wir mit Rücksicht auf die Relationen $\sum_{1}^{\nu}l_{s,i}=0$ ihren Werth gleich
\[
\pm(x_1+x_2+\cdots+x_\nu)L.
\]
Wenn wir also alle $x$ mit Ausnahme von einem beliebigen
\srcpage{679}
$=0$, das eine $=1$ annehmen, so ergiebt sich aus (3) irgend eine der $(\nu-1)$-reihigen Determinanten der Matrix
\[
\begin{array}{ccc}
l_{1,1},&\ldots,&l_{1,\nu}\\
\cdots&\cdots&\cdots\\
l_{\nu-1,1},&\ldots,&l_{\nu-1,\nu}
\end{array}
\]
und alle diese Determinanten haben also (vom Vorzeichen abgesehen) denselben Werth. Es ist daher gleichgültig, welcher unter den $\nu$ conjugirten Werthen bei der Bildung des Regulators $L$ weggelassen wird.

Ist $\varepsilon_1,\varepsilon_2,\ldots,\varepsilon_{\nu-1}$ ein unabhängiges System von Einheiten, und sind $\lambda_1,\lambda_2,\ldots,\lambda_\nu$ die conjugirten Logarithmen irgend einer Einheit $\varepsilon$ in $\Omega$, so kann man die Zahlen $\xi_1,\xi_2,\ldots,\xi_{\nu-1}$ immer und nur auf eine Weise so bestimmen, dass
\begin{equation}
\begin{array}{rcl}
\xi_1l_{1,1}+\xi_2l_{2,1}+\cdots+\xi_{\nu-1}l_{\nu-1,1}&=&\lambda_1,\\
\xi_1l_{1,2}+\xi_2l_{2,2}+\cdots+\xi_{\nu-1}l_{\nu-1,2}&=&\lambda_2,\\
\multicolumn{3}{c}{\cdots\cdots\cdots}\,\\
\xi_1l_{1,\nu}+\xi_2l_{2,\nu}+\cdots+\xi_{\nu-1}l_{\nu-1,\nu}&=&\lambda_\nu
\end{array}
\tag{5}
\end{equation}
wird. Denn von diesen Gleichungen sind die ersten $\nu-1$ ein System linearer Gleichungen mit nicht verschwindender Determinante, und die $\nu$te Gleichung folgt aus den übrigen, weil die Summe der conjugirten Logarithmen einer jeden Einheit verschwindet.

Das System der Zahlen $\xi_1,\xi_2,\ldots,\xi_{\nu-1}$ nennen wir das Exponentensystem der Einheit $\varepsilon$ in Bezug auf das System $\varepsilon_1,\varepsilon_2,\ldots,\varepsilon_{\nu-1}$, und es ist nun zu beweisen:

\textbf{6. Dass die Exponenten jeder Einheit $\varepsilon$ rationale Zahlen sind, deren Nenner eine gewisse endliche Grenze nicht übersteigt, und die daher alle mit einem gemeinschaftlichen, nur von dem Systeme $\varepsilon_i$ abhängigen Nenner behaftet angenommen werden können.}

Um dies einzusehen, hat man Folgendes zu erwägen:

1) Wenn wir die Grössen $\xi_1,\xi_2,\ldots,\xi_{\nu-1}$ wie Variable betrachten, und jede von ihnen, von den anderen unabhängig, von 0 bis 1 gehen lassen, so bleiben die linken Seiten von (5) in endlichen, nur von den Einheiten $\varepsilon_1,\varepsilon_2,\ldots,\varepsilon_{\nu-1}$ abhängigen Grenzen eingeschlossen. In diesen Grenzen müssen also auch die conjugirten Logarithmen $\lambda$ der Einheit $\varepsilon$ bleiben, so lange
\srcpage{680}
die Exponenten $\xi_i$ auf nicht negative echte gebrochene Werthe beschränkt bleiben, und daraus ergiebt sich nach der Definition der conjugirten Logarithmen eine obere Grenze für $\varepsilon$ und seine Conjugirten.

Nach dem Satze 2, \S~185 gibt es aber in $\Omega$ nur eine endliche Anzahl ganzer Zahlen, also um so mehr nur eine endliche Anzahl von Einheiten, die dem absoluten Werthe nach mit allen ihren Conjugirten unter einer endlichen Grenze liegen.

Wenn wir nun eine Einheit, deren Exponenten in den Grenzen 0 und 1 liegen, mit Einschluss der unteren und mit Ausschluss der oberen Grenze eine in Bezug auf das System $\varepsilon_1,\varepsilon_2,\ldots,\varepsilon_{\nu-1}$ reducirte Einheit nennen, so haben wir den Satz:

\textbf{Es gibt nur eine endliche Anzahl von Einheiten in $\Omega$, die in Bezug auf ein unabhängiges System $\varepsilon_1,\varepsilon_2,\ldots,\varepsilon_{\nu-1}$ reducirt sind.}

2) Die Einheiten reproduciren sich durch Multiplication und Division. Sind
\[
\begin{array}{c}
\xi_1',\xi_2',\ldots,\xi_{\nu-1}',\\
\xi_1'',\xi_2'',\ldots,\xi_{\nu-1}''
\end{array}
\]
die Exponenten von zwei Einheiten $\varepsilon',\varepsilon''$, so sind die Summen und die Differenzen
\[
\xi_1'\pm\xi_1'',\ \xi_2'\pm\xi_2'',\ldots,\xi_{\nu-1}'\pm\xi_{\nu-1}''
\]
die Exponenten der Einheiten $\varepsilon'\varepsilon''$ und $\varepsilon':\varepsilon''$.

Dies ergibt sich unmittelbar aus dem Satze, dass der Logarithmus eines Productes oder eines Quotienten gleich der Summe oder der Differenz der Logarithmen der beiden Bestandtheile ist.

3) Die Einheiten $\varepsilon_1,\varepsilon_2,\ldots$ selbst haben die Exponenten
\[
1,0,\ldots,0;\quad 0,1,\ldots,0;\quad\ldots,
\]
und daraus folgt nach 2), dass jedes System von ganzen Zahlen, für $\xi_1,\xi_2,\ldots,\xi_{\nu-1}$ gesetzt, das Exponentensystem einer Einheit gibt, die sich durch Multiplication von Potenzen der $\varepsilon_1,\varepsilon_2,\ldots$ bilden lässt.

4) Aus 2) und 3) ergiebt sich, dass ein Exponentensystem $\xi_1,\xi_2,\ldots,\xi_{\nu-1}$ einer Einheit ein Exponentensystem bleibt, wenn alle seine Elemente mit einer und derselben ganzen Zahl multiplicirt werden, und dass es auch dann noch diesen Charakter behält, wenn seine Elemente $\xi_i$ um beliebige ganze Zahlen vermehrt oder vermindert werden. Gebrauchen wir also das Zeichen $(x)$ in dem Sinne, wie im \S~182, dass es den Ueber-
\srcpage{681}
schuss der Zahl $x$ über die grösste in $x$ enthaltene ganze Zahl bedeutet, so ergiebt sich Folgendes:

Ist $\xi_1,\xi_2,\ldots,\xi_{\nu-1}$ das Exponentensystem einer Einheit, und $m$ eine ganze Zahl, so ist auch
\begin{equation}
(m\xi_1),(m\xi_2),\ldots,(m\xi_{\nu-1})
\tag{6}
\end{equation}
das Exponentensystem einer Einheit.

Nun sind die Grössen $(m\xi_i)$, wenn sie nicht Null sind, positive echte Brüche, und nach 1) muss es sich also, wenn die Reihe der ganzen Zahlen hinlänglich weit fortgesetzt wird, ereignen, dass für zwei verschiedene Werthe $m',m''$ von $m$ die Zahlenreihe (6) übereinstimmt. Da hiernach $m'\xi_i,m''\xi_i$ denselben Ueberschuss über eine ganze Zahl haben, so ist $(m'-m'')\xi_i$ selbst eine ganze Zahl, und es gibt also eine ganze rationale Zahl $m$, die höchstens gleich der Anzahl der Exponentensysteme reducirter Einheiten ist, von der Eigenschaft, dass
\[
m\xi_1,m\xi_2,\ldots,m\xi_{\nu-1}
\]
ganze rationale Zahlen werden.

Diese Zahl $m$ kann sich ändern, wenn die Einheit $\varepsilon$ geändert wird. Da aber alle möglichen Werthe dieser Zahl unter einer endlichen Grenze liegen, so gibt es eine endliche Zahl, in der alle diese Werthe von $m$ enthalten sind und die selbst für $m$ genommen werden kann. Es gibt daher eine gewisse positive ganze Zahl $m$, die als Nenner aller der rationalen Zahlen genommen werden kann, die als Exponenten von Einheiten auftreten können. Damit ist der Satz 6. bewiesen.

\section*{\S~187. Fundamentalsysteme von Einheiten.}

Wir wählen jetzt an Stelle des Systemes unabhängiger Einheiten $\varepsilon_i$ ein anderes, dessen conjugirte Logarithmen
\begin{equation}
\lambda_{i,1},\lambda_{i,2},\ldots,\lambda_{i,\nu}
\qquad i=1,2,\ldots,\nu-1
\tag{1}
\end{equation}
sein mögen. Nach \S~186, (5) ist
\begin{equation}
\begin{array}{rl}
\lambda_{i,k}={}&\xi_{1,i}l_{1,k}+\xi_{2,i}l_{2,k}+\cdots+\xi_{\nu-1,i}l_{\nu-1,k},\\[2pt]
&i=1,2,\ldots,\nu-1;\quad k=1,2,\ldots,\nu,
\end{array}
\tag{2}
\end{equation}
wenn $\xi_{1,i},\xi_{2,i},\ldots,\xi_{\nu-1,i}$ die Exponenten einer der neuen Einheiten in Bezug auf das System $\varepsilon_1,\varepsilon_2,\ldots,\varepsilon_{\nu-1}$ bedeuten, die wir als rationale Zahlen mit dem gemeinsamen Nenner $m$ annehmen können.

\srcpage{682}
Es ist aber nach dem Multiplicationssatze der Determinanten
\begin{equation}
\left|\begin{array}{ccc}
\lambda_{1,1},&\ldots,&\lambda_{\nu-1,1}\\
\cdots&\cdots&\cdots\\
\lambda_{1,\nu-1},&\ldots,&\lambda_{\nu-1,\nu-1}
\end{array}\right|
=
\left|\begin{array}{ccc}
\xi_{1,1},&\ldots,&\xi_{1,\nu-1}\\
\cdots&\cdots&\cdots\\
\xi_{\nu-1,1},&\ldots,&\xi_{\nu-1,\nu-1}
\end{array}\right|
\left|\begin{array}{ccc}
l_{1,1},&\ldots,&l_{\nu-1,1}\\
\cdots&\cdots&\cdots\\
l_{1,\nu-1},&\ldots,&l_{\nu-1,\nu-1}
\end{array}\right|,
\tag{3}
\end{equation}
oder, wenn wir
\[
\Lambda_{\nu-1}=\sum\pm\lambda_{1,1}\lambda_{2,2}\cdots\lambda_{\nu-1,\nu-1}
\]
setzen, und beachten, dass die Determinante
\[
\sum\pm\xi_{1,1}\xi_{2,2}\cdots\xi_{\nu-1,\nu-1}
\]
eine rationale Zahl mit dem Nenner $m^{\nu-1}$ ist:
\begin{equation}
\Lambda_{\nu-1}=\frac{a}{m^{\nu-1}}L(\varepsilon_1,\varepsilon_2,\ldots,\varepsilon_{\nu-1}),
\tag{4}
\end{equation}
worin $a$ eine ganze rationale Zahl ist.

Daraus ergiebt sich, dass die neu eingeführten Einheiten immer und nur dann ein unabhängiges System bilden, wenn die Determinante der $\xi_{i,k}$, d.~h. die Zahl $a$, von Null verschieden ist.

Nun gibt es unter einer endlichen oder unendlichen Anzahl nicht verschwindender rationaler Brüche mit demselben Nenner immer einen dem absoluten Werthe nach kleinsten, und folglich können wir nach (4) das neue System unabhängiger Einheiten so wählen, dass sein Regulator $\pm\Lambda_{\nu-1}$ so klein als möglich wird. Ein solches System nennen wir ein Fundamentalsystem von Einheiten, und den Minimalwerth des Regulators selbst, also den Regulator eines Fundamentalsystems, nennen wir den Regulator des Körpers. Wir nehmen jetzt an, dass unser System $\varepsilon_1,\varepsilon_2,\ldots,\varepsilon_{\nu-1}$ selbst ein Fundamentalsystem sei.

Unter dieser Voraussetzung lässt sich beweisen, dass alle Exponenten von Einheiten ganze Zahlen sein müssen.

Wenn wir nämlich annehmen, es existire eine Einheit $\varepsilon$, deren nach den Formeln \S~186, (5) bestimmte Exponenten nicht alle ganze Zahlen sind, so gibt es nach \S~186, 4) auch eine Einheit $\varepsilon_0$, deren Exponenten echte Brüche sind, die nicht alle verschwinden.

Verstehen wir unter $\lambda_1,\lambda_2,\ldots,\lambda_{\nu-1}$ in den Formeln (5), \S~186 das System der conjugirten Logarithmen von $\varepsilon_0$, so ergiebt sich
\[
L(\varepsilon_0,\varepsilon_2,\ldots,\varepsilon_{\nu-1})=\xi_1L(\varepsilon_1,\varepsilon_2,\ldots,\varepsilon_{\nu-1}),
\]
\srcpage{683}
und wenn wir also annehmen, was offenbar keine Beschränkung der Allgemeinheit ist, dass $\xi_1$ ein nicht verschwindender echter Bruch sei, so widerspricht diese Formel der Annahme, dass
\[
\varepsilon_1,\varepsilon_2,\ldots,\varepsilon_{\nu-1}
\]
ein Fundamentalsystem sei, weil die Determinante $L_{\nu-1}$ verkleinert würde, wenn $\varepsilon_1$ durch $\varepsilon_0$ ersetzt wird.

Haben wir ein Fundamentalsystem, so können wir ein beliebiges System ganzer rationaler Zahlen $\xi_1,\xi_2,\ldots,\xi_{\nu-1}$ als Exponentensystem annehmen und eine Einheit $\varepsilon$ mit diesen Exponenten bilden:
\[
\varepsilon=\varepsilon_1^{\xi_1}\varepsilon_2^{\xi_2}\cdots\varepsilon_{\nu-1}^{\xi_{\nu-1}}.
\]

Es bleibt also noch die Frage zu beantworten, inwieweit eine Einheit durch das Exponentensystem bestimmt ist.

Nehmen wir an, es seien $\varepsilon',\varepsilon''$ zwei Einheiten mit demselben Exponentensysteme, dann besteht das Exponentensystem der Einheit $\varepsilon':\varepsilon''=\rho$ aus lauter Nullen, und folglich sind die conjugirten Logarithmen von $\rho$ alle gleich Null; $\rho$ ist daher eine ganze Zahl von der Eigenschaft, dass die absoluten Werthe aller mit $\rho$ conjugirten Zahlen gleich 1 sind, und daher, wie im \S~175, 2. bewiesen ist, eine Einheitswurzel. Damit ist das folgende Theorem bewiesen:

\textbf{1. Es gibt im Körper $\Omega$ ein System von $\nu-1$ fundamentalen Einheiten}
\begin{equation}
\varepsilon_1,\varepsilon_2,\ldots,\varepsilon_{\nu-1},
\tag{5}
\end{equation}
\textbf{welches die Eigenschaft hat, dass in der Form}
\begin{equation}
\varepsilon=\rho\varepsilon_1^{\xi_1}\varepsilon_2^{\xi_2}\cdots\varepsilon_{\nu-1}^{\xi_{\nu-1}}
\tag{6}
\end{equation}
\textbf{alle Einheiten des Körpers, jede nur einmal, enthalten sind, wenn $\xi_1,\xi_2,\ldots,\xi_{\nu-1}$ alle ganzen rationalen Zahlen und $\rho$ alle in $\Omega$ vorhandenen Einheitswurzeln durchläuft.}

Es ist nun leicht, aus einem Fundamentalsystem alle anderen abzuleiten, indem man in (6) für die Exponenten $\nu-1$ verschiedene Systeme ganzer Zahlen $\xi_{i,k}$ setzt, deren Determinante $=\pm1$ ist. Denn setzt man
\[
\varepsilon_i'=\rho_i\varepsilon_1^{\xi_{1,i}}\varepsilon_2^{\xi_{2,i}}\cdots\varepsilon_{\nu-1}^{\xi_{\nu-1,i}},
\]
so folgt aus (3):
\[
L(\varepsilon_1',\varepsilon_2',\ldots,\varepsilon_{\nu-1}')=\pm L(\varepsilon_1,\varepsilon_2,\ldots,\varepsilon_{\nu-1}).
\]

\srcpage{684}
Beide Determinanten haben also den Minimalwerth, und beide Systeme sind Fundamentalsysteme.

Der Fall $\nu=1$, den wir oben ausgeschlossen haben, tritt ausser im Falle des rationalen Körpers, in dem nur die beiden Einheiten $\pm1$ existiren, im Falle des imaginären quadratischen Körpers ein. In diesem Falle haben wir im zwanzigsten Abschnitte die ganzen Zahlen des Körpers in der Form
\[
\frac{x+y\sqrt{-m}}{2}
\]
dargestellt, worin $-m$ eine Stammdiscriminante bedeutet, so dass $m\equiv0$ oder $\equiv3\pmod 4$ ist, und $x,y$ beide gerade oder beide ungerade sind. Die Einheiten erhalten wir, wenn wir die Gleichung
\[
x^2+my^2=4
\]
auf alle mögliche Arten in ganzen rationalen Zahlen lösen. Diese Gleichung hat aber, wenn $m>4$ ist, nur die zwei Lösungen $x=\pm2,y=0$, und es gibt also in diesen Fällen, wie im Körper der rationalen Zahlen, nur die zwei Einheiten $\pm1$. Ist $m=4$, so findet man die vier Einheiten $\pm1,\pm i$, und ist endlich $m=3$, die sechs Einheiten
\[
\pm1,
\quad \pm\frac{-1+i\sqrt3}{2},
\quad \pm\frac{-1-i\sqrt3}{2}.
\]

In dem nächst einfachen Falle, $n=2$, $\nu=2$, d.~h. im reellen quadratischen Körper, fällt die Theorie der Einheiten zusammen mit der im \S~127 des ersten Bandes behandelten Theorie der Pell'schen Gleichung.


\section*{\S\,188. Reducirte Zahlen.}

\srcpage{684}
Ist $\alpha$ irgend eine von Null verschiedene ganze oder gebrochene Zahl des Körpers $\Omega$, so betrachten wir, wie bei den Einheiten, das System der conjugirten Logarithmen von $\alpha$, worunter wir, wie im \S~185, die reellen Zahlen
\begin{equation}
\lambda_1=\delta_1\log|\alpha_1|,\quad
\lambda_2=\delta_2\log|\alpha_2|,\ldots,\quad
\lambda_\nu=\delta_\nu\log|\alpha_\nu|
\tag{1}
\end{equation}
verstehen, die wir auch mit $\lambda_s(\alpha)$ bezeichnen. Ziehen wir das Fundamentalsystem $\varepsilon_1,\varepsilon_2,\ldots,\varepsilon_{\nu-1}$ von Einheiten zu Rathe, so

\srcpage{685}
können wir die linearen Gleichungen, analog wie im \S~186, (5) bilden:
\begin{equation}
\begin{aligned}
\xi_1l_{1,1}+\xi_2l_{2,1}+\cdots+\xi_{\nu-1}l_{\nu-1,1}+\delta_1\xi_\nu&=\lambda_1,\\
\xi_1l_{1,2}+\xi_2l_{2,2}+\cdots+\xi_{\nu-1}l_{\nu-1,2}+\delta_2\xi_\nu&=\lambda_2,\\
&\cdots\cdots\cdots\\
\xi_1l_{1,\nu}+\xi_2l_{2,\nu}+\cdots+\xi_{\nu-1}l_{\nu-1,\nu}+\delta_\nu\xi_\nu&=\lambda_\nu,
\end{aligned}
\tag{2}
\end{equation}
dessen Determinante [\S~186, (4)] den von Null verschiedenen Werth hat:
\begin{equation}
\left|
\begin{array}{ccccc}
l_{1,1}&l_{2,1}&\ldots&l_{\nu-1,1}&\delta_1\\
l_{1,2}&l_{2,2}&\ldots&l_{\nu-1,2}&\delta_2\\
\cdots&\cdots&\cdots&\cdots&\cdots\\
l_{1,\nu}&l_{2,\nu}&\ldots&l_{\nu-1,\nu}&\delta_\nu
\end{array}
\right|=nL(\varepsilon_1,\varepsilon_2,\ldots,\varepsilon_{\nu-1}).
\tag{3}
\end{equation}
Demnach sind die Unbekannten $\xi_1,\xi_2,\ldots,\xi_{\nu-1},\xi_\nu$ aus (2) eindeutig bestimmt.

Durch Addition der Gleichungen (2) findet man zunächst sehr einfach
\begin{equation}
n\xi_\nu=\log \Na(\alpha),
\tag{4}
\end{equation}
und $\xi_\nu$ bleibt also dasselbe für alle Zahlen $\alpha$, die unter einander associirt sind.

Die Zahlen $\xi_1,\xi_2,\ldots,\xi_{\nu-1}$ heissen die Exponenten der Zahl $\alpha$ [in Bezug auf das Fundamentalsystem $(\varepsilon_1,\varepsilon_2,\ldots,\varepsilon_{\nu-1})$].

Aus dieser Definition ergiebt sich ohne Weiteres, dass die Exponenten eines Productes oder eines Quotienten gleich der Summe oder der Differenz der entsprechenden Exponenten der einzelnen Bestandtheile sind. Die Exponenten einer Einheit sind ganze rationale Zahlen, und jedes System ganzer rationaler Zahlen $\xi_1,\xi_2,\ldots,\xi_{\nu-1}$ ist auch das Exponentensystem einer Einheit.

Die Exponenten associirter Zahlen unterscheiden sich also um ganze Zahlen von einander, und man kann zu jeder Zahl $\alpha$ eine associirte Zahl $\alpha_0$ finden, deren Exponenten zwischen 0 und 1 liegen. Solche Zahlen heissen \textbf{reducirte Zahlen} (in Bezug auf das Fundamentalsystem $\varepsilon_1,\ldots,\varepsilon_{\nu-1}$). Reducirte Einheiten sind alle und nur die in $\Omega$ vorhandenen Einheitswurzeln, deren Zahl wir mit $w$ bezeichnen wollen. Da die beiden Einheitswurzeln $\pm1$ in jedem Körper enthalten sind, so ist $w$ mindestens $=2$ und immer eine gerade Zahl.

Wenn zwei associirte Zahlen dasselbe Exponentensystem haben, so unterscheiden sie sich nur durch einen Factor, der

\srcpage{686}
eine Einheitswurzel ist. Das Exponentensystem der aus $\alpha$ abgeleiteten reducirten Zahl $\alpha_0$ ist durch $\alpha$ selbst völlig bestimmt. Hierdurch ist der Satz bewiesen:

\begin{center}
\textbf{II. Zu jeder (von Null verschiedenen) Zahl $\alpha$ des Körpers $\Omega$ giebt es $w$ und nicht mehr reducirte Zahlen $\rho\alpha_0$.}
\end{center}

Wenn statt des Fundamentalsystems $(\varepsilon_1,\varepsilon_2,\ldots,\varepsilon_{\nu-1})$ ein anderes angewendet wird, so erleiden die Exponenten $\xi_1,\xi_2,\ldots,\xi_{\nu-1}$ eine ganzzahlige lineare Substitution von der Determinante $\pm1$. Eine reducirte Zahl kann dann aufhören, für das neue Fundamentalsystem reducirt zu sein. Der Begriff der reducirten Zahl ist also von der Wahl des Fundamentalsystems abhängig. Die Anzahl der aus einer gegebenen Zahl abgeleiteten reducirten Zahlen ist aber immer dieselbe.

Durch diesen Satz haben wir den Zweck erreicht, aus dem ganzen unendlichen Systeme der unter einander associirten Zahlen eine bestimmte endliche Anzahl von Repräsentanten herausgehoben zu haben.

\section*{\S\,189. Grenzen der Anzahl der durch ein Ideal theilbaren ganzen Zahlen des Körpers $\Omega$.}

Unsere früheren Betrachtungen haben ergeben, dass jede ganze Zahl eines Körpers $\Omega$ nur eine endliche Anzahl von Idealen zu Theilern hat. Da jedes ganze Ideal ein Theiler seiner Norm ist, so giebt es also auch nur eine endliche Anzahl von ganzen Zahlen in $\Omega$, deren absolute Norm eine gegebene Grenze nicht übersteigt.

Aus diesen Zahlen greifen wir jetzt wieder einen Theil heraus, und fragen nach der Anzahl $T$ aller \textbf{ganzen Zahlen} des Körpers $\Omega$, deren absolute Norm eine positive Grösse $t$ nicht überschreitet, und die durch ein gegebenes Ideal $\mathfrak a$ theilbar sind.

Nehmen wir als vorläufiges Beispiel den rationalen Körper, so ist dort $T$ die Anzahl der natürlichen Zahlen, die kleiner als $t$ und durch eine gegebene ganze Zahl $m$ theilbar sind, also, in der früher (\S~182) gebrauchten Bezeichnung $T=E\left(\frac{t}{m}\right)$.

\srcpage{687}
Die Zahl $T$ ist eine im Allgemeinen nicht genauer zu bestimmende Function von $t$, die mit $t$ zugleich ins Unendliche wächst, die sich aber, da sie nur ganzzahlige Werthe hat, nur stufenweise, also unstetig ändert. Worauf es hier ankommt, ist der \textbf{Grenzwerth, dem sich das Verhältniss $T:t$ mit unendlich wachsendem $t$ nähert} (der in dem oben angeführten Beispiele $1:m$ ist).

Dieser Grenzwerth lässt sich durch Betrachtungen, wie wir sie schon im \S~155 benutzt haben, finden.

Wir bilden uns zunächst eine Basisform des gegebenen Ideals $\mathfrak a$:
\begin{equation}
\alpha=\alpha_1x_1+\alpha_2x_2+\cdots+\alpha_nx_n,
\tag{1}
\end{equation}
d.~h. eine lineare Function der $n$ Variablen $x_i$, von der Eigenschaft, dass durch Substitution aller ganzen rationalen Zahlen für $x_i$ aus $\alpha$ alle durch $\mathfrak a$ theilbare ganze Zahlen des Körpers $\Omega$ entstehen (\S~146). Mit $\alpha_{i,1},\alpha_{i,2},\ldots,\alpha_{i,n}$ bezeichnen wir die mit $\alpha_i$ conjugirten Werthe und bilden die lineare Substitution für $x_i$:
\begin{equation}
\begin{aligned}
y_1&=\alpha_{1,1}x_1+\alpha_{2,1}x_2+\cdots+\alpha_{n,1}x_n,\\
y_2&=\alpha_{1,2}x_1+\alpha_{2,2}x_2+\cdots+\alpha_{n,2}x_n,\\
&\cdots\cdots\cdots\\
y_n&=\alpha_{1,n}x_1+\alpha_{2,n}x_2+\cdots+\alpha_{n,n}x_n.
\end{aligned}
\tag{2}
\end{equation}
Ist $A$ die Determinante dieses Gleichungssystems
\[
A=\sum\pm\alpha_{1,1}\alpha_{2,2}\cdots\alpha_{n,n},
\]
so ist $A^2$ die Discriminante des Functionensystems $(\alpha_1,\alpha_2,\ldots,\alpha_n)$, und daher nach \S~147 gleich dem Producte aus dem Quadrate der absoluten Norm von $\alpha$ und der Körperdiscriminante $\Delta$. Da die absolute Norm von $\alpha$ zugleich die Norm des Ideals $\mathfrak a$ ist (\S~151), so setzen wir
\begin{equation}
A^2=N(\mathfrak a)^2\Delta.
\tag{3}
\end{equation}

Nun können wir, wie im \S~155, die $x_1,x_2,\ldots,x_n$ als Coordinaten eines Punktes $(x)$ in einem Raume von $n$ Dimensionen deuten, und wir lassen jedem Punkt $(x)$ einen Punkt $(x')$ entsprechen, dessen Coordinaten durch
\begin{equation}
x_1'=t^{\frac1n}x_1,
\quad x_2'=t^{\frac1n}x_2,
\ldots,
\quad x_n'=t^{\frac1n}x_n
\tag{4}
\end{equation}
bestimmt sind, worin $t$ eine beliebige positive Grösse bedeutet. Jedem Gebiete $S$ von Punkten $(x)$ entspricht ein Gebiet $S'$ von Punkten $(x')$. Alle Figuren in $S$ und entsprechenden Figuren in $S'$ ähnlich, und die Lineardimensionen in $S$ verhalten sich
zu denen in $S'$ wie $1:t^{\frac1n}$.

\srcpage{688}
Denken wir uns das Gebiet $S$ gegeben, so wird das entsprechende Gebiet $S'$ von $t$ abhängig sein und sich mit $t$ vergrössern. Die Punkte mit ganzzahligen Coordinaten $x_i$ nennen wir, wie früher, die Gitterpunkte. Die Anzahl der in einem endlichen Gebiete $S'$ liegenden Gitterpunkte ist immer endlich, wächst aber mit $t$ und soll mit $Z_t$ bezeichnet werden. Dann ist nach den Sätzen \S~155, (14) das Volumen des Gebietes $S$, d.~h. das über alle Punkte von $S$ erstreckte $n$-fache Integral
\[
\int\!\int\cdots\int dx_1\,dx_2\cdots dx_n
\]
gleich dem Grenzwerthe des Verhältnisses $Z_t:t$ für unendlich wachsendes $t$:
\begin{equation}
V=\lim_{t=\infty}\frac{Z_t}{t}.
\tag{5}
\end{equation}

Hierin liegt zugleich die genauere Präcisirung, die wir über das Gebiet $S$ machen müssen, wenn unsere Betrachtung anwendbar sein soll. Es muss $S$ so beschaffen sein, dass das $n$-fache Integral $V$ eine bestimmte Bedeutung hat. Jeder Gitterpunkt ist nun durch (1) das Bild einer ganzen durch $\mathfrak a$ theilbaren Zahl des Körpers $\Omega$, und wir können also auch sagen, dass $Z_t$ die Anzahl der durch $\mathfrak a$ theilbaren ganzen Zahlen in $\Omega$ ist, deren Bilder in dem Gebiete $S'$ liegen.

Um das Gebiet $S'$ in geeigneter Weise abzugrenzen, wählen wir ein System fundamentaler Einheiten des Körpers $\Omega$
\[
\varepsilon_1,\varepsilon_2,\ldots,\varepsilon_{\nu-1}
\]
mit dem System der conjugirten Logarithmen
\[
l_{i,1},l_{i,2},\ldots,l_{i,\nu}
\]
und dem Regulator
\begin{equation}
\pm L(\varepsilon_1,\varepsilon_2,\ldots,\varepsilon_{\nu-1})=\pm\sum\pm l_{1,1}l_{2,2}\cdots l_{\nu-1,\nu-1},
\tag{6}
\end{equation}
der mit $L$ bezeichnet werden mag.

Wir definiren jetzt ein System von Functionen $\xi_1,\xi_2,\ldots,\xi_\nu$ der Variablen $x_i$ durch folgende Bedingungen:

Es sei, wenn die $y_i$ die Bedeutung (2) haben,
\begin{equation}
z_i=\delta_i\log|y_i|,
\tag{7}
\end{equation}
so dass, wenn für die $x_i$ ein System ganzer rationaler Zahlen gesetzt wird, $z_i$ nach \S~185 in das System der conjugirten Logarithmen der Zahl $\alpha$ übergeht.

Es ist dann, wenn $y_k$ zu einem reellen Körper gehört,
\begin{equation}
z_k=\frac12\log y_k^2,
\tag{8}
\end{equation}
\srcpage{689}
und wenn $y_k,y_k'$ einem imaginären Paare angehören,
\begin{equation}
z_k=\log y_ky_k',
\tag{9}
\end{equation}
wodurch die Definition völlig eindeutig ist.

Nun bestimmen wir die Variablen $\xi_1,\xi_2,\ldots,\xi_{\nu-1},\xi_\nu$ aus den linearen Gleichungen
\begin{equation}
\begin{aligned}
\xi_1l_{1,1}+\cdots+\xi_{\nu-1}l_{\nu-1,1}+\delta_1\xi_\nu&=z_1,\\
\xi_1l_{1,2}+\cdots+\xi_{\nu-1}l_{\nu-1,2}+\delta_2\xi_\nu&=z_2,\\
&\cdots\cdots\cdots\\
\xi_1l_{1,\nu}+\cdots+\xi_{\nu-1}l_{\nu-1,\nu}+\delta_\nu\xi_\nu&=z_\nu,
\end{aligned}
\tag{10}
\end{equation}
deren Determinante nicht verschwindet.

Giebt man den Variablen $x_i$ ganzzahlige Werthe, so gehen nach \S~188, (2), (3) die $\xi_1,\xi_2,\ldots,\xi_{\nu-1}$ in das Exponentensystem der Zahl $\alpha$ über, und es ergiebt sich durch Addition der Gleichungen (10) [\S~185, (3)]:
\begin{equation}
\xi_\nu=\frac1n\log \Na(\alpha).
\tag{11}
\end{equation}

Wir gehen jetzt von dem Punkte $(x)$ zu dem Punkte $(x')$ über, indem wir $x_i'=t^{\frac1n}x_i$ setzen. Dadurch geht $y_i$ in
\[
y_i'=t^{\frac1n}y_i
\]
über, und $z_i$ in
\[
z_i'=z_i+\frac1n\delta_i\log t.
\]
Wenn wir also $z_i'$ an Stelle von $z_i$ setzen, und die dann aus (10) sich ergebenden Werthe der $\xi_1,\xi_2,\ldots,\xi_\nu$ mit $\xi_1',\xi_2',\ldots,\xi_\nu'$ bezeichnen, so ist
\begin{equation}
\xi_1'=\xi_1,
\quad \xi_2'=\xi_2,
\ldots,
\quad \xi_{\nu-1}'=\xi_{\nu-1},
\quad \xi_\nu'=\xi_\nu+\frac1n\log t.
\tag{12}
\end{equation}

Nunmehr wollen wir das Gebiet $S'$ dadurch abgrenzen, dass
\begin{equation}
0\leq\xi_1'<1,\ldots,\ 0\leq\xi_{\nu-1}'<1,
\qquad
\xi_\nu'<\frac1n\log t
\tag{13}
\end{equation}
sein soll. Dadurch erreichen wir, dass die in dem Gebiete $S'$ liegenden Gitterpunkte $(x')$ nur reducirte ganze Zahlen $\alpha$ darstellen, und alle und nur solche, deren absolute Norm kleiner als $t$ ist.

Unter den reducirten Zahlen sind aber immer je $w$ mit einander associirt, wenn $w$ die Anzahl der in $\Omega$ enthaltenen

\srcpage{690}
Einheitswurzeln bedeutet, und wenn wir also diese $w$ associirten Zahlen zu einem Complex zusammenfassen, so ist die Anzahl dieser Complexe gleich der Anzahl $T$ aller nicht associirten, durch $\mathfrak a$ theilbaren ganzen Zahlen in $\Omega$, deren absolute Norm kleiner als $t$ ist. Demnach ist die Anzahl der in $S'$ liegenden Gitterpunkte
\begin{equation}
Z_t=wT.
\tag{14}
\end{equation}

Für die Begrenzung des Gebiets $S$ erhält man aus (12) und (13):
\begin{equation}
0\leq\xi_1<1,\ldots,
0\leq\xi_{\nu-1}<1,
\quad \xi_\nu<0,
\tag{15}
\end{equation}
und nach (5) erhalten wir
\begin{equation}
\lim_{t=\infty}\frac{T}{t}
=\frac1w\int\!\int\cdots\int dx_1\,dx_2\cdots dx_n
=\frac1w V,
\tag{16}
\end{equation}
worin das $n$-fache Integral über das durch (15) bestimmte Gebiet $S$ auszudehnen ist.

\section*{\S\,190. Bestimmung des Volumens.}

Das Volumen $V$ lässt sich bestimmen nach der schon im \S~155 angeführten Transformationsformel für mehrfache Integrale, nach der, wenn $x_1,x_2,\ldots,x_n$ Functionen der Variablen $u_1,u_2,\ldots,u_n$ sind,
\begin{equation}
\begin{aligned}
V&=\int\!\int\cdots\int dx_1\,dx_2\cdots dx_n\\
&=\int\!\int\cdots\int\left(\sum\pm
\frac{\partial x_1}{\partial u_1}
\frac{\partial x_2}{\partial u_2}\cdots
\frac{\partial x_n}{\partial u_n}\right)du_1\,du_2\cdots du_n
\end{aligned}
\tag{1}
\end{equation}
ist, wobei die Functionaldeterminante in dem nach den Variablen $u$ genommenen Integrale mit ihrem absoluten Werthe zu nehmen ist.

Diese Functionaldeterminante ist auch gleich dem reciproken Werthe der Determinante
\[
\sum\pm
\frac{\partial u_1}{\partial x_1}
\frac{\partial u_2}{\partial x_2}\cdots
\frac{\partial u_n}{\partial x_n}.
\]

Wir führen zunächst die Variablen $u$ durch eine lineare Substitution ein, und bemerken, dass unter den mit $\Omega$ conjugirten Körpern nach unserer Annahme $n-\nu$ conjugirt imaginäre Paare

\srcpage{691}
und $2\nu-n$ reelle vorkommen. Demnach sind von den $n$ Functionen $y$ [\S~189, (2)]
\begin{equation}
y_k=\alpha_{1,k}x_1+\alpha_{2,k}x_2+\cdots+\alpha_{n,k}x_n
\tag{2}
\end{equation}
$2\nu-n$ reell, $2n-2\nu$ paarweise conjugirt. Wir wollen unter den $n$ Variablen $u$ reelle lineare Verbindungen der $x$ verstehen, nämlich die $2\nu-n$ reellen $y$ und die $2n-2\nu$ Bestandtheile der conjugirt imaginären Paare der $y$, so dass, wenn $y,y'$ ein imaginäres Paar ist,
\[
2u_1=y+y',\qquad 2iu_2=y-y',
\]
oder
\begin{equation}
y=u_1+iu_2,
\qquad y'=u_1-iu_2
\tag{3}
\end{equation}
gesetzt wird. Wenn man in der Functionaldeterminante der $y$ auf jedes Paar conjugirt imaginärer Zeilen zweimal den elementaren Determinantensatz von der Addition der Zeilen anwendet (Bd.~I, \S~22), so ergiebt sich
\[
\sum\pm
\frac{\partial u_1}{\partial x_1}
\frac{\partial u_2}{\partial x_2}\cdots
\frac{\partial u_n}{\partial x_n}
=
\frac1{(-2i)^{n-\nu}}
\sum\pm
\frac{\partial y_1}{\partial x_1}
\frac{\partial y_2}{\partial x_2}\cdots
\frac{\partial y_n}{\partial x_n}
=
\frac{A}{(-2i)^{n-\nu}},
\]
worin $A$ die in \S~189, (3) festgesetzte Bedeutung hat.

Die Körperdiscriminante ist positiv oder negativ, je nachdem die Anzahl der conjugirt imaginären Paare, also $n-\nu$, gerade oder ungerade ist, und demnach erhält man, wenn mit $\pm\Delta$ der absolute Werth der Körperdiscriminante bezeichnet wird, nach \S~189, (4), (16):
\begin{equation}
\lim\frac{T}{t}
=
\frac{2^{n-\nu}}{wN(\mathfrak a)\sqrt{\pm\Delta}}
\int\!\int\cdots\int du_1\,du_2\cdots du_n.
\tag{4}
\end{equation}
Es ist also weiter noch das Integral
\[
U=\int\!\int\cdots\int du_1\,du_2\cdots du_n
\]
zu behandeln.

Aus den Gleichungen \S~189, (10) erhält man zu jedem den Bedingungen (15) genügenden Werthsysteme der Variablen $\xi_1,\xi_2,\ldots,\xi_\nu$ ein bestimmtes endliches System der Werthe für $z_1,z_2,\ldots,z_\nu$, und daraus nach \S~189, (7) für jedes $y$ einen von Null verschiedenen absoluten Werth. Setzt man im Falle eines imaginären Paares (3)
\begin{align}
u_1&=e^{\frac12 z}\cos\vartheta,
&u_2&=e^{\frac12 z}\sin\vartheta,
\tag{5}\\
y_1y_1'&=u_1^2+u_2^2=e^z,
\tag{6}
\end{align}
\srcpage{692}
so ergiebt sich, wenn $\vartheta$ in den Grenzen
\[
0\leq\vartheta<2\pi
\]
genommen wird, bei feststehendem $z$ zu jedem Werthe von $\vartheta$ ein zugehöriges Werthepaar $u_1,u_2$. Ein solcher Winkel $\vartheta$ gehört zu jedem imaginären Paare.

Es gehört dann zu jedem Punkte des Gebietes $S$ ein und nur ein Werthsystem der Variablen $\xi_1,\xi_2,\ldots,\xi_\nu,\vartheta\ldots$ in der Begrenzung
\begin{equation}
0\leq\xi_1<1,\ldots,
0\leq\xi_{\nu-1}<1,
\qquad
\xi_\nu<0,
\qquad
0\leq\vartheta<2\pi,\ldots
\tag{7}
\end{equation}

Wenn umgekehrt irgend ein Werthsystem der Variablen $\xi,\vartheta$ in den Grenzen (7) gegeben ist, so erhält man hieraus die entsprechenden Bestandtheile $u_1,u_2$ eines imaginären $y$ eindeutig, von den reellen $y$ aber nur die absoluten Werthe, und es kann also jedem der reellen $y$ noch jedes der beiden Vorzeichen gegeben werden. Die Anzahl dieser möglichen Bestimmungen ist $2^{2\nu-n}$. Ist eine dieser Bestimmungen herausgegriffen, so sind die Variablen $x$ eindeutig bestimmt und sind für jedes endliche Werthsystem der $\xi,\vartheta$ endlich. Für ein gegen $-\infty$ abnehmendes $\xi_\nu$ werden die Werthe der $y$ zum Theil gegen Null convergiren.

Das Gebiet $S$ zerfällt demnach in $2^{2\nu-n}$ Theilgebiete $S_1,S_2,\ldots$, deren jedes durch eine bestimmte Vorzeichen-Combination der reellen $y_1,y_2,\ldots$ charakterisirt ist, und demnach zerfällt das Integral $U$ in ebenso viele Theilintegrale $U_1,U_2,\ldots$
\[
U=U_1+U_2+\cdots.
\]
Bei der Ausführung der Integration in einem dieser Bestandtheile, die sich alle als von gleichem Werthe ergeben, fangen wir mit einem imaginären Paare an, wenn ein solches vorhanden ist.

Betrachten wir $u_1,u_2$ als rechtwinkelige Coordinaten in einer Ebene, so können wir $e^{\frac12 z},\vartheta$ nach (5) als Polarcoordinaten auffassen, und erhalten
\[
\int\!\int du_1\,du_2=\frac12\int\!\int e^z\,dz\,d\vartheta,
\]
worin die Integration in Bezug auf $\vartheta$, die sich von 0 bis $2\pi$ erstreckt, ausgeführt werden kann. Man findet so
\[
\int\!\int du_1\,du_2=\pi\int e^z\,dz.
\]

\srcpage{693}
Nun ändert man die Reihenfolge der Integration, und stellt ein etwa noch vorhandenes zweites imaginäres Paar voran, das man wieder in derselben Weise behandelt.

Die verschiedenen Bestandtheile $U_1,U_2,\ldots$ unterscheiden sich durch die Vorzeichen der den reellen $y$ entsprechenden $u$. Nehmen wir das Zeichen so, dass $\pm u$ positiv wird, und setzen
\begin{equation}
z=\log(\pm u),
\qquad \pm du=e^z\,dz,
\tag{8}
\end{equation}
erhalten wir für die verschiedenen $U_1,U_2,\ldots$ denselben Ausdruck durch ein $\nu$-faches Integral
\begin{equation}
U_1=\pi^{n-\nu}\int\!\int\cdots\int e^{z_1+z_2+\cdots+z_\nu}\,dz_1\,dz_2\cdots dz_\nu,
\tag{9}
\end{equation}
wenn wir $dz_1,dz_2,\ldots,dz_\nu$ positiv annehmen.

Die Variablen $z_k$ sind aber nach (6) und (8) hier keine anderen, als in \S~189, (8), (9), (10), und danach ist
\[
z_1+z_2+\cdots+z_\nu=n\xi_\nu.
\]
Wenn wir nun die Transformationsformel (1) auf das $\nu$-fache Integral (9) anwenden, um die Variablen $\xi_1,\xi_2,\ldots,\xi_\nu$ einzuführen, so haben wir die Determinante
\[
\sum\pm
\frac{\partial z_1}{\partial\xi_1}
\frac{\partial z_2}{\partial\xi_2}\cdots
\frac{\partial z_\nu}{\partial\xi_\nu}
\]
zu bilden, die nach \S~188, (3) dem absoluten Werthe nach mit $nL$, d.~h. mit dem $n$fachen Werthe des Körperregulators $L$ übereinstimmt, und danach ist
\[
U_1=\pi^{n-\nu}nL\int_0^1d\xi_1\cdots\int_0^1d\xi_{\nu-1}\int_{-\infty}^{0}e^{n\xi_\nu}\,d\xi_\nu
=\pi^{n-\nu}L
\]
und
\[
U=2^{2\nu-n}\pi^{n-\nu}L.
\]
Daraus ergiebt sich also endlich nach (4) der Satz:

\begin{center}
\textbf{1. Bedeutet $T$ die Anzahl der durch $\mathfrak a$ theilbaren nicht associirten ganzen Zahlen, deren absolute Norm kleiner als $t$ ist, so ist}
\end{center}
\begin{equation}
\lim_{t=\infty}\frac{T}{t}
=
\frac{2^\nu\pi^{n-\nu}L}{wN(\mathfrak a)\sqrt{\pm\Delta}}
=
\frac{g}{N(\mathfrak a)},
\tag{10}
\end{equation}
\begin{center}
\textbf{worin $g$ eine durch die Natur des Körpers $\Omega$ völlig bestimmte positive Zahl ist, nämlich:}
\end{center}
\begin{equation}
g=\frac{2^\nu\pi^{n-\nu}L}{w\sqrt{\pm\Delta}}.
\tag{11}
\end{equation}

\section*{\S\,191. Sätze aus der Reihenlehre.}

\srcpage{694}
Bei den weiteren Anwendungen der bisherigen Resultate sind einige Sätze aus der Lehre von den unendlichen Reihen erforderlich, die zunächst hier abgeleitet werden sollen.

\begin{enumerate}[label=\textbf{\arabic*.},leftmargin=2.2em]
\item Ist $a_1,a_2,a_3,\ldots,a_n,\ldots$ ein unbegrenztes System reeller (positiver, negativer oder auch verschwindender) Grössen, von dem wir voraussetzen, dass die Summe
\begin{equation}
\sigma_n=a_1+a_2+\cdots+a_n
\tag{1}
\end{equation}
für jedes beliebige $n$ dem absoluten Werthe nach unter einer endlichen Grenze $C$ bleibt, so ist die Reihe
\begin{equation}
S=\frac{a_1}{1^s}+\frac{a_2}{2^s}+\frac{a_3}{3^s}+\cdots
\tag{2}
\end{equation}
für jedes positive $s$ nicht nur convergent, sondern auch eine stetige Function von $s$.
\end{enumerate}

Um diesen Satz zu beweisen, zerlegen wir $S$ in der Weise:
\[
S=S_m+R_m,
\]
worin
\[
S_m=\frac{a_1}{1^s}+\frac{a_2}{2^s}+\cdots+\frac{a_{m-1}}{(m-1)^s},
\]
\[
R_m=\frac{a_m}{m^s}+\frac{a_{m+1}}{(m+1)^s}+\frac{a_{m+2}}{(m+2)^s}+\cdots .
\]
Nun ist nach (1):
\[
a_m=\sigma_m-\sigma_{m-1},\quad a_{m+1}=\sigma_{m+1}-\sigma_m,\quad
a_{m+2}=\sigma_{m+2}-\sigma_{m+1},\ldots
\]
und daher
\[
\begin{aligned}
R_m+\frac{\sigma_{m-1}}{m^s}
&=\sigma_m\left(\frac1{m^s}-\frac1{(m+1)^s}\right)\\
&\quad+\sigma_{m+1}\left(\frac1{(m+1)^s}-\frac1{(m+2)^s}\right)+\cdots .
\end{aligned}
\]
Da nun nach der Voraussetzung $\sigma_{m-1},\sigma_m,\sigma_{m+1},\ldots$ zwischen endlichen Grenzen $\pm C$ eingeschlossen sind, so ist hiernach $R_m+\frac{\sigma_{m-1}}{m^s}$ zwischen den beiden Grenzen
\[
\pm C\left(\frac1{m^s}-\frac1{(m+1)^s}+\frac1{(m+1)^s}-\frac1{(m+2)^s}+\cdots\right)
=\pm\frac{C}{m^s}
\]

\srcpage{695}
eingeschlossen, und es ist dem absoluten Werthe nach
\[
R_m<\frac{2C}{m^s},
\]
oder, wenn $s>c$ ist,
\[
R_m<\frac{2C}{m^c}.
\]

Diese obere Grenze für $R_m$, die von $s$ unabhängig ist, kann aber, wenn $c$ positiv ist, dadurch, dass man $m$ hinlänglich gross annimmt, unter jeden noch so kleinen Werth herabgedrückt werden.

Daraus folgt aber nicht nur die Convergenz, sondern auch die Stetigkeit von $S$. Denn nimmt man $m$ hinlänglich gross, so wird nicht nur $R_m$, sondern es werden auch die Schwankung von $R_m$ bei veränderlichem $s$ unendlich klein, und $S_m$ ist für ein feststehendes $m$ eine stetige Function von $s$. Also ist auch $S$ stetig.

Es ist dabei noch zu bemerken, dass die Voraussetzung über $\sigma_n$ keineswegs die Convergenz der unendlichen Reihe $\sum a_k$ voraussetzt. Es ist nicht einmal erforderlich, dass die $a_k$ reell seien; denn wenn sie imaginär sind, so braucht man nur den Satz auf den reellen und den imaginären Bestandtheil anzuwenden. Endlich ist es auch nicht nothwendig, die $a_1,a_2,a_3,\ldots$ als Constanten vorauszusetzen. Alles bleibt gültig, wenn es stetige Functionen von $s$ sind.

\begin{enumerate}[label=\textbf{\arabic*.},leftmargin=2.2em,start=2]
\item Bedeutet $\mu_1,\mu_2,\ldots,\mu_n,\ldots$ eine unendliche Menge positiver Zahlen von der Beschaffenheit, dass zwei endliche Zahlen $\alpha,\beta$ so angegeben werden können, dass für jedes noch so grosse $n$
\begin{equation}
\alpha<\frac{n}{\mu_n}<\beta,
\tag{3}
\end{equation}
so ist die unendliche Reihe
\[
S=\frac1{\mu_1^s}+\frac1{\mu_2^s}+\frac1{\mu_3^s}+\cdots
\]
convergent, so lange der Exponent $s$ grösser als $1$ ist, und das Product $(s-1)S$ bleibt bei hinlänglicher Annäherung von $s$ an den Werth $1$ gleichfalls zwischen den Grenzen $\alpha$ und $\beta$, oder nähert sich wenigstens einem dieser Werthe unbegrenzt an.
\end{enumerate}

\srcpage{696}
Der Beweis ist durch Zurückführung auf ein Integral sehr einfach zu führen. Offenbar ist nämlich
\[
\frac1{(n+1)^s}<\int_n^{n+1}\frac{dx}{x^s}<\frac1{n^s},
\]
weil sich der erste oder der zweite Werth, $(n+1)^{-s}$ oder $n^{-s}$, ergiebt, wenn unter dem Integralzeichen für $x^{-s}$ der zu kleine oder der zu grosse Werth $(n+1)^{-s}$ oder $n^{-s}$ gesetzt wird.

Demnach ergiebt sich aus (3) für jedes positive $s$
\[
\alpha^s\int_n^{n+1}\frac{dx}{x^s}<\frac1{\mu_n^s}<\beta^s\int_{n-1}^{n}\frac{dx}{x^s}.
\]
Setzen wir
\[
R_m=\frac1{\mu_{m+1}^s}+\frac1{\mu_{m+2}^s}+\frac1{\mu_{m+3}^s}+\cdots,
\]
so ergiebt sich hieraus
\[
\alpha^s\int_{m+1}^{\infty}\frac{dx}{x^s}<R_m<\beta^s\int_m^{\infty}\frac{dx}{x^s}
\]
oder
\begin{equation}
\frac{\alpha^s}{(s-1)(m+1)^{s-1}}<R_m<\frac{\beta^s}{(s-1)m^{s-1}} .
\tag{4}
\end{equation}

Es bleibt also $R_m$, wie viele Glieder auch summirt werden mögen, wenn $s>1$ ist, endlich, und dies ist bei Reihen mit positiven Gliedern eine ausreichende Bedingung für die Convergenz.

Setzt man ferner (4) in die Form
\[
\frac{\alpha^s}{(m+1)^{s-1}}<(s-1)R_m<\frac{\beta^s}{m^{s-1}},
\]
so sieht man, dass die beiden Grenzen beliebig nahe an $\alpha,\beta$ gebracht werden können, wenn man $s-1$ klein genug annimmt, und daraus ergiebt sich, dass $(s-1)R_m$ bei hinlänglich kleinem $s-1$ nicht mehr ausserhalb der Grenzen $\alpha,\beta$ liegen kann, wenn es sich nicht mit abnehmendem $s-1$ einem dieser Werthe unbegrenzt annähert. Da sich nun $R_m$ von $S$ nur durch einen Bestandtheil unterscheidet, der für jedes positive $s$ endlich ist, so ist der Satz hiermit bewiesen.

Es gilt der Satz aber auch dann noch, wenn die Ungleichheit (3) nicht für alle $n$ gilt, sondern wenn sie nur besteht, sobald $n$ eine beliebig gegebene Grenze $m$ überschritten hat. Auf

\srcpage{697}
Grund dieser Bemerkung kann man, wenn sich $n:\mu_n$ einem endlichen Grenzwerthe $\gamma$ nähert, $\alpha$ und $\beta$ diesem Grenzwerthe beliebig nahe annehmen, und erhält so die etwas präcisere Fassung des Theorems:

\begin{enumerate}[label=\textbf{\arabic*.},leftmargin=2.2em,start=3]
\item Sind die positiven Zahlen $\mu_1,\mu_2,\mu_3,\ldots$ so beschaffen, dass
\[
\lim_{n=\infty}\frac{n}{\mu_n}=\gamma
\]
ein endlicher Grenzwerth ist, so ist die unendliche Reihe
\[
S=\frac1{\mu_1^s}+\frac1{\mu_2^s}+\frac1{\mu_3^s}+\cdots
\]
für jedes $s$, was grösser als $1$ ist, convergent, und es ist
\[
\lim_{s=1}(s-1)S=\gamma.
\]
\end{enumerate}

Es sei jetzt irgend ein System $M$ von unendlich vielen positiven Zahlen $\mu_n$ gegeben, die wir so ordnen
\begin{equation}
\mu_1\leqq\mu_2\leqq\mu_3\leqq\cdots\quad (M),
\tag{5}
\end{equation}
so dass in der Reihe (5) niemals ein kleineres Glied auf ein grösseres folgt, und es bedeute $T$ die Anzahl von Gliedern in $M$, die nicht grösser als eine beliebig gegebene positive Grösse $t$ sind. Es gilt dann folgender Satz:

\begin{enumerate}[label=\textbf{\arabic*.},leftmargin=2.2em,start=4]
\item Wenn einer der beiden Grenzwerthe
\begin{equation}
\lim_{n=\infty}\frac{n}{\mu_n},\qquad \lim_{t=\infty}\frac{T}{t}
\tag{6}
\end{equation}
endlich ist, so hat der andere denselben endlichen Werth.
\end{enumerate}

Es sei zunächst
\begin{equation}
\lim_{n=\infty}\frac{n}{\mu_n}=\gamma
\tag{7}
\end{equation}
ein endlicher Grenzwerth. Dann werden die $\mu_n$ mit unendlich wachsendem $n$ nothwendig ins Unendliche wachsen müssen, und es giebt für jedes positive $t$ einen Werth $m$, so dass
\begin{equation}
\mu_m\leqq t<\mu_{m+1};
\tag{8}
\end{equation}
$m$ wächst zugleich mit $t$ ins Unendliche und es ist $T=m$ [wegen (5)]. Daher nach (8)
\begin{equation}
\frac{m}{\mu_m}\geqq\frac{T}{t}>\frac{m}{\mu_{m+1}} .
\tag{9}
\end{equation}

\srcpage{698}
Wenn nun $\gamma=0$ ist, so folgt hieraus, indem man $t$ und $m$ ins Unendliche wachsen lässt, dass auch $T:t$ den Grenzwerth $0$ hat. Ist aber $\gamma$ von $0$ verschieden, so ist nach (7)
\[
\lim\frac{\mu_m}{\mu_{m+1}}=\lim\frac{m}{m+1}=1,
\]
und daher nähern sich die beiden Grenzwerthe in (9) dem Werthe $\gamma$ und es folgt also:
\begin{equation}
\lim_{t=\infty}\frac{T}{t}=\gamma.
\tag{10}
\end{equation}

Setzen wir zweitens umgekehrt die Grenzgleichung (10) voraus, so wird auch jetzt $\mu_n$ mit $n$ ins Unendliche wachsen müssen, denn sonst würde, da die Gesammtzahl aller $\mu$ unendlich ist, $T$ schon für ein endliches $t$ unendlich, und $\gamma$ könnte nicht endlich sein.

Nehmen wir nun einen Werth $\mu_n$, der in der Reihe der unter einander gleichen
\[
\mu_{m+1},\mu_{m+2},\ldots,\mu_{m+l}
\]
vorkommt, so dass
\[
m+1\leqq n\leqq m+l,
\]
so wird $T$, wenn $t$ durch den Werth $\mu_n$ geht, plötzlich um $l$ Einheiten wachsen, und das Verhältniss $T:t$ wächst um $l:\mu_n$. Da aber $T:t$ einen endlichen Grenzwerth haben soll, so muss
\begin{equation}
\lim\frac{l}{\mu_n}=0
\tag{11}
\end{equation}
sein. Wenn nun $t=\mu_n$ ist, so ist $T=m+l$, und folglich ist
\begin{equation}
\frac{T}{t}=\frac{m+l}{\mu_n},\qquad
\lim\frac{m+l}{\mu_n}=\gamma.
\tag{12}
\end{equation}
Ferner
\[
\frac{m}{\mu_n}<\frac{n}{\mu_n}\leqq\frac{m+l}{\mu_n},
\]
und folglich ist wegen (11) und (12)
\[
\lim\frac{n}{\mu_n}=\lim\frac{m}{\mu_n}=\gamma.
\]
Also ist aus der Gleichung (10) die Gleichung (7) gefolgert\footnote{Dieser Satz ist im Wesentlichen auf dieselbe Weise bewiesen bei Dirichlet-Dedekind, Supplement II. Vgl. auch Riemann's mathematische Werke, 2. Auflage, Nr. XXX.}.

Mit Benutzung des Satzes 4. können wir nun dem Satze 3. auch die Form geben:

\srcpage{699}
\begin{enumerate}[label=\textbf{\arabic*.},leftmargin=2.2em,start=5]
\item Ist $T$ die Anzahl der Grössen $\mu_n$, die nicht grösser als $t$ sind, so ist
\begin{equation}
\lim_{s=1}\sum_{n=1}^{\infty}\frac{s-1}{\mu_n^s}=\lim_{t=\infty}\frac{T}{t},
\tag{13}
\end{equation}
wenn der Grenzwerth von $T:t$ endlich ist.
\end{enumerate}

Nehmen wir an, dass die Grössen $\mu_1,\mu_2,\mu_3,\ldots$ nicht nur den Bedingungen des Satzes 3., sondern noch den weiteren genügen, dass, wenn
\begin{equation}
\mu_n=\frac{n+c_n}{\gamma}
\tag{14}
\end{equation}
gesetzt wird, die $c_n$ mit unendlich wachsendem $n$ nicht unendlich werden (d.\ h. für jedes noch so grosse $n$ zwischen endlichen Grenzen eingeschlossen sind), so können wir den Satz 3. durch den noch schärferen ersetzen:

\begin{enumerate}[label=\textbf{\arabic*.},leftmargin=2.2em,start=6]
\item Haben die Zahlen $\mu_n$ die Eigenschaft, dass sich ein endliches $\gamma$ so bestimmen lässt, dass
\begin{equation}
\gamma\mu_n-n=c_n
\tag{15}
\end{equation}
mit unendlich wachsendem $n$ nicht unendlich wird, so hat die für jedes $s>1$ definirte Function von $s$
\begin{equation}
S=\frac1{\mu_1^s}+\frac1{\mu_2^s}+\cdots
\tag{16}
\end{equation}
die Eigenschaft, dass die Differenz
\begin{equation}
S-\frac{\gamma}{s-1}=C,
\tag{17}
\end{equation}
sich mit unendlich abnehmendem $s-1$ einer endlichen Grenze nähert.
\end{enumerate}

Um dies zu beweisen, setzen wir nach (14) und (16):
\begin{equation}
\sum_{n=1}^{\infty}\frac{\gamma^s}{n^s}-S
=\gamma^s\sum_{n=1}^{\infty}\frac1{n^s}
\left[1-\left(1+\frac{c_n}{n}\right)^{-s}\right].
\tag{18}
\end{equation}
Wenn nun $\varepsilon$ ein positiver echter Bruch ist, so ist, wenn wir
\[
a_n=\frac1{n^\varepsilon}\left[1-\left(1+\frac{c_n}{n}\right)^{-s}\right]
\]
setzen, $a_n n^{1+\varepsilon}$ für ein unendlich wachsendes $n$ nicht unendlich, und folglich ist nach einem bekannten elementaren Satze der Reihenlehre $a_1+a_2+\cdots$ eine unbedingt convergente Reihe.

\srcpage{700}
Setzen wir nun
\[
s=s_1+\varepsilon,
\]
so ergiebt sich aus (18):
\[
\sum\frac{\gamma^s}{n^s}-S=\gamma^s\sum\frac{a_n}{n^{s_1}},
\]
und dies ist nach 1. eine für positive $s_1$ endliche und stetige Function von $s_1$. Für $s_1=1-\varepsilon$ wird aber $s=1$, und folglich ist
\begin{equation}
\sum\frac{\gamma^s}{n^s}-S=D_s
\tag{19}
\end{equation}
für $s=1$ endlich, nämlich gleich
\[
D_1=\sum\frac{\gamma c_n}{n(n+c_n)}.
\]

Mit Anwendung einfacher Sätze aus der Theorie der $\Gamma$-Functionen lässt sich aber die Summe $\sum\frac1{n^s}$ durch ein bestimmtes Integral ausdrücken. Es ist
\[
\frac1{n^s}=\frac1{\Gamma(s)}\int_0^\infty e^{-nx}x^{s-1}\,dx,
\]
und folglich
\[
\sum\frac1{n^s}=\frac1{\Gamma(s)}\int_0^\infty \frac{x^{s-1}e^{-x}\,dx}{1-e^{-x}}.
\]
Ferner ist [nach dem Satze $(s-1)\Gamma(s-1)=\Gamma(s)$]:
\[
\frac1{s-1}=\frac1{\Gamma(s)}\int_0^\infty e^{-x}x^{s-2}\,dx,
\]
und daraus
\[
\sum\frac1{n^s}-\frac1{s-1}
=\frac1{\Gamma(s)}\int_0^\infty x^{s-1}e^{-x}
\left(\frac1{1-e^{-x}}-\frac1x\right)dx.
\]
Hierin ist die rechte Seite eine für alle positiven Werthe von $s$ stetige Function, die für $s=1$ in die Euler'sche Constante
\[
\int_0^\infty e^{-x}\left(\frac1{1-e^{-x}}-\frac1x\right)dx
=-\Gamma'(1)=0{,}57721566\ldots
\]
übergeht. Nun ist nach (19):
\[
D_s=\gamma^s\left(\sum\frac1{n^s}-\frac1{s-1}\right)
-\left(S-\frac{\gamma^s}{s-1}\right),
\]

\srcpage{701}
und daher
\begin{equation}
\lim_{s=1}\left(S-\frac{\gamma^s}{s-1}\right)
=-\gamma\Gamma'(1)+D_1=C_1,
\tag{20}
\end{equation}
wie zu beweisen war, endlich\footnote{Die hier gebrauchten Sätze über $\Gamma$-Functionen finden sich in den ausführlicheren Lehrbüchern der Integralrechnung, z.\ B. Serret-Harnack, Lehrbuch der Differential- und Integralrechnung, Bd. II, S. 170 f. (Leipzig 1885).}.
\section*{\S\,192. Anwendung auf die Bestimmung der Classenzahl.}

Wir machen von diesen Sätzen jetzt die Anwendung auf die Theorie des Körpers $\Omega$. Wir verstehen unter den Zahlen $\mu_n$ des Theorems 5. die Normen der sämmtlichen Ideale des Körpers $\Omega$, also unter $T$ die Anzahl der Ideale in $\Omega$, deren Norm nicht grösser als eine gegebene positive Grösse $t$ ist, und erhalten
\begin{equation}
\lim_{s=1}\sum\frac{s-1}{N(\mathfrak a)^s}=\lim_{t=\infty}\frac{T}{t},
\tag{1}
\end{equation}
und darin erstreckt sich die Summe der linken Seite auf alle Ideale $\mathfrak a$ des Körpers. Die Ideale zerfallen nun nach \S~153 in eine endliche Anzahl von Classen
\[
A_1,A_2,\ldots,A_h,
\]
und das grosse Ziel ist die Bestimmung dieser Zahl $h$, der Classenzahl.

Die Ideale $\mathfrak a_1$ einer dieser Classen $A_1$ sind dadurch charakterisirt, dass ihre Producte mit einem und demselben ganzen Functionale $\varphi_1$, das der Classe $A_1^{-1}$ angehört, Hauptideale sind, dass also
\[
\varphi_1\mathfrak a_1=\alpha
\]
eine ganze Zahl des Körpers $\Omega$ ist.

Ist die Norm von $\mathfrak a_1$ nicht grösser als $t$, so ist
\[
N_a(\alpha)\leqq N_a(\varphi_1)t=t_1.
\]
Ist $T_1$ die Anzahl der Ideale der Classe $A_1$, deren Normen nicht grösser als $t$ sind, so ist $T_1$ zugleich die Anzahl der nicht

\srcpage{702}
associirten durch $\varphi_1$ theilbaren ganzen Zahlen, deren Normen nicht grösser als $t_1$ sind, und nach dem Satze \S~190, 1. ist
\begin{equation}
\lim_{t_1=\infty}\frac{T_1}{t_1}=\frac{g}{N_a(\varphi_1)},
\tag{2}
\end{equation}
wenn $g$ die an der erwähnten Stelle angegebene Bedeutung hat, also eine von Null verschiedene, durch die Natur des Körpers $\Omega$ bestimmte Zahl ist.

Wenn jetzt $T_2,t_2,\ldots,T_h,t_h$ die entsprechende Bedeutung für die Classen $A_2,\ldots,A_h$ haben, wie $T_1,t_1$ für $A_1$, so ist
\[
T=T_1+T_2+\cdots+T_h,
\]
\[
\frac{T}{t}=\frac{T_1}{t_1}N_a(\varphi_1)+\frac{T_2}{t_2}N_a(\varphi_2)+\cdots+
\frac{T_h}{t_h}N_a(\varphi_h),
\]
und aus (1) und (2) ergiebt sich die fundamentale Formel:
\begin{equation}
\lim_{s=1}\sum\frac{s-1}{N(\mathfrak a)^s}=gh.
\tag{3}
\end{equation}

Die Zahl $g$ können wir nach \S~185, (11) als bekannt betrachten, wenn auch ihre wirkliche Berechnung noch auf der Voraussetzung beruht, dass ein Fundamentalsystem von Einheiten bekannt sei, und wenn auch die Ermittelung eines solchen Systems in höheren Körpern immer als eine der grössten Schwierigkeiten betrachtet worden ist. Dann hängt die Berechnung der Classenzahl noch von der Bestimmung des Grenzwerthes auf der linken Seite von (3) ab, die natürlich auch nur in besonderen Fällen gelingt, doch aber oft wichtige Schlüsse über die Natur der Classenzahl gestattet. Wir wollen noch eine die Berechnung vorbereitende Umformung der Summe
\begin{equation}
\sum\frac1{N(\mathfrak a)^s}=\Phi(s)
\tag{4}
\end{equation}
in ein unendliches Product entwickeln.

Es sei $\mathfrak p$ irgend ein Primideal $f^{\mathrm{ten}}$ Grades im Körper $\Omega$, also
\[
N(\mathfrak p)=p^f.
\]
Dann ist die unendliche geometrische Reihe
\[
1+\frac1{N(\mathfrak p)^s}+\frac1{N(\mathfrak p)^{2s}}+\cdots
=\frac1{1-p^{-sf}};
\]
wenn man diese Reihen für alle verschiedenen Primideale $\mathfrak p$ mit einander multiplicirt, so erhält man nach bekannten Sätzen aus

\srcpage{703}
der Lehre von den unendlichen Reihen eine Summe von Gliedern der Form
\[
\left(\frac1{N(\mathfrak p)^x}\frac1{N(\mathfrak p')^{x'}}\frac1{N(\mathfrak p'')^{x''}}\cdots\right)^s
=\left(\frac1{N(\mathfrak p^x\mathfrak p'^{x'}\mathfrak p''^{x''}\cdots)}\right)^s,
\]
die alle in der Form $N(\mathfrak a)^{-s}$ enthalten sind, und jedes solche Glied ergiebt sich ein- und nur einmal. Danach ist also
\begin{equation}
\Phi(s)=\prod_{\mathfrak p}\frac1{1-N(\mathfrak p)^{-s}}.
\tag{5}
\end{equation}
Sind daher $\mathfrak p_1,\mathfrak p_2,\ldots,\mathfrak p_e$ die von einander verschiedenen Primfactoren von $p$, und $f_1,f_2,\ldots,f_e$ ihre Grade (\S~143), so ergiebt sich
\begin{equation}
\Phi(s)=\prod_p
\frac1{(1-p^{-sf_1})(1-p^{-sf_2})\cdots(1-p^{-sf_e})},
\tag{6}
\end{equation}
worin das unendliche Product $\prod$ über alle natürlichen Primzahlen $p$ auszudehnen ist, und nach (3), (4)
\begin{equation}
gh=\lim_{s=1}(s-1)\Phi(s).
\tag{7}
\end{equation}

Hieraus lässt sich ein für mannigfache Anwendungen wichtiger Schluss ziehen:

Wenn wir aus der Entwickelung
\[
1+\frac1{p^s}+\frac1{p^{2s}}+\cdots=\frac1{1-p^{-s}}
\]
das Product für alle Primzahlen $p$ bilden, so ergiebt sich:
\begin{equation}
\prod_p\frac1{1-p^{-s}}=\sum_{n=1}^{\infty}\frac1{n^s},
\tag{8}
\end{equation}
worin sich die Summe auf der rechten Seite auf alle natürlichen Zahlen $n$ erstreckt, und es hat also dies Product für alle Werthe von $s$, die grösser als $1$ sind, einen endlichen Werth. Der reciproke Werth, nämlich das Product
\begin{equation}
P=\prod_p(1-p^{-s}),
\tag{9}
\end{equation}
dessen Factoren sämmtlich kleiner als $1$ sind, hat aber, so lange $s>1$ ist, einen von Null verschiedenen Werth. Diese Eigenschaft bleibt nun erhalten, wenn wir bei der Bildung des Productes $P$ nur einen Theil aller Primzahlen $p$ berücksichtigen, weil das Product durch Weglassen beliebiger Factoren nur vergrössert wird, ohne die Einheit je zu übersteigen.

\srcpage{704}
Daraus ergiebt sich, dass in dem Producte (6) die Theilproducte
\[
\prod\frac1{1-p^{-sf}},
\]
die sich über alle Primzahlen $p$ erstrecken, für die $f>1$ ist, auch für $s=1$ noch endlich bleiben. Da andererseits $g,h$ bestimmte von Null verschiedene Werthe haben, so ergiebt sich aus (7) der wichtige Satz:

\begin{center}
\textbf{I. Durchläuft $\mathfrak p$ die Gesammtheit der Primideale ersten Grades irgend eines algebraischen Körpers $\Omega$, so hat das Product}
\end{center}
\begin{equation}
(s-1)\prod\frac1{1-N(\mathfrak p)^{-s}}
\tag{10}
\end{equation}
\begin{center}
\textbf{für $s=1$ einen endlichen von Null verschiedenen Grenzwerth.}
\end{center}

Diese Eigenschaft bleibt auch dann noch erhalten, wenn bei der Bildung des Productes eine beliebige endliche Anzahl von Primidealen ersten Grades ausgelassen wird.

Die Normen $N(\mathfrak p)$ sind in der Formel (10) gleich natürlichen Primzahlen $p$. Eine Primzahl $p$ kommt aber darin so oft vor, als sie Primfactoren ersten Grades in $\Omega$ enthält, also höchstens $n$ mal. Ist $\Omega$ ein Normalkörper, so kommt auch wirklich jede Primzahl $p$, die in Primfactoren ersten Grades zerlegbar ist, genau $n$ mal darin vor, und wir können für das Product (10) auch setzen:
\[
(s-1)\prod\frac1{(1-p^{-s})^n},
\]
worin sich das Product $\prod$ auf alle in Primfactoren ersten Grades zerlegbaren Primzahlen $p$ erstreckt (\S~160).

Eine unmittelbare Folgerung des Satzes I. ist die:
\begin{center}
\textbf{In jedem algebraischen Körper giebt es unendlich viele Primideale ersten Grades.}
\end{center}

\section*{Vierundzwanzigster Abschnitt.\\Classenzahl der Kreistheilungskörper.}
\section*{\S\,193. Classenzahldarstellung im Kreistheilungskörper $\Omega_m$.}

\srcpage{705}
Die allgemeinen Resultate über die Classenzahl $h$ sollen nun angewandt werden auf den vollen Kreistheilungskörper $\Omega_m$ unter der bisherigen Voraussetzung, dass $m$ eine Potenz einer Primzahl $q$ sei, und dass also
\begin{equation}
m=q^\alpha,\qquad \varphi(m)=q^{\alpha-1}(q-1)=\mu
\tag{1}
\end{equation}
gesetzt sei.

Im einundzwanzigsten Abschnitte (\S~169, 171) haben wir gesehen, dass in $q$ nur ein Primideal $1^{\mathrm{ten}}$ Grades aufgeht, und dass eine zum Exponenten $f$ gehörige, von $q$ verschiedene Primzahl $p$ in $e$ verschiedene Primideale $f^{\mathrm{ten}}$ Grades zerfällt. Darin bedeutet $f$ den kleinsten positiven Exponenten, für den
\begin{equation}
p^f\equiv1\pmod m
\tag{2}
\end{equation}
ist, und $e$ ist durch die Gleichung
\begin{equation}
\mu=ef
\tag{3}
\end{equation}
bestimmt.

Die am Ende des vorigen Paragraphen durch die Formel (6) bestimmte Function $\Phi(s)$ erhält daher hier den Ausdruck:
\begin{equation}
\Phi(s)=\frac1{1-q^{-s}}\prod\frac1{(1-p^{-sf})^e},
\tag{4}
\end{equation}
und darin erstreckt sich das Productzeichen $\prod$ auf alle Primzahlen $p$, die von $q$ verschieden sind. Darin ist $s$ eine Variable, die immer grösser als $1$ ist, die sich schliesslich der Grenze $1$ nähern soll.

\srcpage{706}
Um nun diesen Ausdruck für die Function $\Phi$ weiter umzuformen und zur Berechnung vorzubereiten, hat man die Sätze über die Abel'schen Gruppen und Gruppencharaktere zu benutzen, die wir in den Paragraphen 14 bis 16 dieses Bandes kennen gelernt haben.

Die Gesammtheit der durch $q$ nicht theilbaren Zahlen $n$ bildet, nach dem Modul $m$ genommen, eine Abel'sche Gruppe $\mathfrak R$ vom Grade $\mu$, deren Charaktere $\chi(n)$ im \S~16 bestimmt sind. Es war dabei ein kleiner Unterschied, je nachdem $q$ ungerade oder $q=2$ ist.

1) Ist $q$ ungerade, so giebt es eine primitive Wurzel $c$ von $m$, und für jede Zahl $n$ giebt es einen Index $\gamma$, so dass
\begin{equation}
n\equiv c^\gamma\pmod m.
\tag{5}
\end{equation}
Ist dann $\Theta$ eine $\mu^{\mathrm{te}}$ Einheitswurzel, so ist
\begin{equation}
\chi(n)=\Theta^\gamma,
\tag{6}
\end{equation}
und die sämmtlichen $\mu$ Charaktere erhält man, wenn man für $\Theta$ die verschiedenen $\mu^{\mathrm{ten}}$ Einheitswurzeln setzt.

Gehört $n$ zum Exponenten $f$, so ist $e$ der grösste gemeinschaftliche Theiler von $\gamma$ und $\mu$, und $\chi(n)$ ist $f^{\mathrm{te}}$ Einheitswurzel.

2) Für $q=2$, $m>4$ (den Fall $m=4$ berücksichtigen wir hier nicht) hat jede ungerade Zahl $n$ zwei Indices $\alpha,\beta$, die nach den Moduln $2,\frac12\mu$ bestimmt sind, für die
\begin{equation}
n\equiv(-1)^\alpha 5^\beta\pmod m
\tag{7}
\end{equation}
ist. Als Charaktere erhält man, wenn $\varepsilon=\pm1$ und $\Theta$ gleich einer $\frac12\mu^{\mathrm{ten}}$ Einheitswurzel gesetzt wird:
\begin{equation}
\chi(n)=\varepsilon^\alpha\Theta^\beta;
\tag{8}
\end{equation}
$f$ ist die kleinste positive Lösung der beiden Congruenzen:
\[
\alpha f\equiv0\pmod2,\qquad \beta f\equiv0\pmod{\tfrac12\mu},
\]
und $\chi(n)$ ist auch hier $f^{\mathrm{te}}$ Einheitswurzel.

Für beide Fälle gilt aber die Bemerkung:

3) Gehört $n$ zum Exponenten $f$, so erhält man, wenn man $\chi$ die gesammte Gruppe der $\mu$ Charaktere durchlaufen lässt, aus $\chi(n)$ jede $f^{\mathrm{te}}$ Einheitswurzel $e$ mal.

Denn da die $\chi$ eine Abel'sche Gruppe bilden, so erhält man zunächst jede $f^{\mathrm{te}}$ Einheitswurzel, die überhaupt vorkommt, gleich oft, nämlich so oft als den Werth $1$, weil, wenn $\chi_1(n)=\chi_2(n)$

\srcpage{707}
ist, $\chi_1\chi_2^{-1}(n)=\chi_0(n)=1$ ist. Andererseits erhält man aber jede $f^{\mathrm{te}}$ Einheitswurzel; denn wenn man im Falle 1) für $\Theta$ eine primitive $\mu^{\mathrm{te}}$, im Falle 2) eine primitive $\frac12\mu^{\mathrm{te}}$ Einheitswurzel, und im letzteren Falle zugleich $\varepsilon=-1$ setzt, so ist $\chi(n)$ eine primitive $f^{\mathrm{te}}$ Einheitswurzel und aus deren Potenzen kann man alle $f^{\mathrm{ten}}$ Einheitswurzeln herleiten.

Bedeutet daher jetzt $x$ eine Variable, so ist, wenn $n$ zum Exponenten $f$ gehört, das über sämmtliche Charaktere $\chi$ erstreckte Product
\begin{equation}
\prod_\chi[1-\chi(n)x]=(1-x^f)^e,
\tag{9}
\end{equation}
und wenn man daher $x=p^{-s}$ setzt, so ergiebt sich aus (4):
\begin{equation}
\Phi(s)=\frac1{1-q^{-s}}\prod_\chi\prod_p\frac1{1-\chi(p)p^{-s}},
\tag{10}
\end{equation}
wenn sich das erste Productzeichen auf alle Charaktere $\chi$, das zweite auf alle Primzahlen $p$ erstreckt. Es ist also $\Phi(s)$ ein Product aus einer endlichen Zahl, $\mu$, von Factoren, deren jeder ein unendliches Product ist.

Jetzt wollen wir jeden Factor eines dieser unendlichen Producte nach steigenden Potenzen von $p^{-s}$ entwickeln. Dadurch ergiebt sich
\[
\frac1{1-\chi(p)p^{-s}}
=1+\chi(p)p^{-s}+\chi(p^2)p^{-2s}+\chi(p^3)p^{-3s}+\cdots .
\]
Das Product aus allen Factoren, die man hieraus erhält, wenn $\chi$ festgehalten wird, während $p$ alle von $q$ verschiedenen Primzahlen durchläuft, ist ein Aggregat von Gliedern der Form
\[
\chi(p^k)\chi(p'^{\,k'})\chi(p''^{\,k''})\cdots p^{-sk}p'^{-sk'}p''^{-sk''}\cdots
=\chi(n)n^{-s},
\]
und jedes Glied dieser Form kommt in dem Producte ein- und nur einmal vor [vgl. \S~192, (8)]. Danach ergiebt sich die Umformung von $\Phi(s)$ in ein endliches Product von unendlichen Reihen:
\begin{equation}
\Phi(s)=\frac1{1-q^{-s}}\prod_\chi\sum_n\frac{\chi(n)}{n^s},
\tag{11}
\end{equation}
worin sich die Summe auf alle durch $q$ nicht theilbaren Zahlen $n$ erstreckt und das Product auf alle $\mu$ Charaktere $\chi$.

Dieser Ausdruck ist geeignet, um den Grenzwerth von $(s-1)\Phi(s)$ für $s=1$ zu bestimmen. Betrachten wir zunächst die dem Hauptcharakter $\chi=1$ entsprechende Summe
\[
\sum_n\frac1{n^s}.
\]

\srcpage{708}
Man erhält die Zahlen $n$, wenn man von der Gesammtheit aller natürlichen Zahlen $k$ die Vielfachen von $q$, d.\ h. die Gesammtheit der Zahlen $qk$ wegnimmt. Hiernach ist
\[
\sum_n\frac1{n^s}=\sum_k\frac1{k^s}-\sum_k\frac1{q^s k^s}
=(1-q^{-s})\sum_k\frac1{k^s}.
\]
Wenn man aber in dem Satze 3., \S~191, für die $\mu_n$ die Reihe der natürlichen Zahlen $k$ setzt, so folgt:
\[
\lim_{s=1}(s-1)\sum_k\frac1{k^s}=1,
\]
und wir erhalten daher
\begin{equation}
\lim_{s=1}\frac{s-1}{1-q^{-s}}\sum_n\frac1{n^s}=1.
\tag{12}
\end{equation}

Die anderen Factoren des Productes $\Phi(s)$ aber sind, wenn die unendlichen Reihen nach steigenden Werthen von $n$ geordnet werden, nach dem Satze 1., \S~191, stetige Functionen von $s$.

Denn die nach steigenden Werthen von $n$ geordnete Summe
\begin{equation}
\sum\chi(n)
\tag{13}
\end{equation}
ist zwar nicht convergent, kann aber doch dem absoluten Werthe nach nicht über eine endliche Grenze hinausgehen. Denn so oft $n$ ein volles Restsystem nach dem Modul $m$ durchläuft, kommt zu der Summe (13) ein Beitrag hinzu, der nach \S~11, 6. den Werth $0$ hat. Demnach ist, wenn $\chi$ nicht der Hauptcharakter ist,
\begin{equation}
\lim_{s=1}\sum_n\frac{\chi(n)}{n^s}=\sum_n\frac{\chi(n)}{n},
\tag{14}
\end{equation}
und demnach erhalten wir für die Classenzahl $h$ im Kreistheilungskörper $\Omega_m$ jetzt den Ausdruck
\begin{equation}
gh=\prod_\chi\sum_n\frac{\chi(n)}{n},
\tag{15}
\end{equation}
in dem sich das Productzeichen $\prod$ nur noch auf die vom Hauptcharakter verschiedenen Charaktere $\chi$ bezieht.

\section*{\S\,194. Bestimmung der Summen $X$.}

Die unendlichen Reihen
\begin{equation}
X=\sum_n\frac{\chi(n)}{n},
\tag{1}
\end{equation}

\srcpage{709}
von denen hiernach die Bestimmung der Classenzahl noch abhängt, lassen sich durch endliche Ausdrücke darstellen. Es ist nämlich
\[
\frac1n=\int_0^1x^{n-1}\,dx,
\]
und demnach
\[
X=\int_0^1\sum_n\chi(n)x^{n-1}\,dx.
\]
Nun bleibt $\chi(n)$ ungeändert, wenn $n$ um ein Vielfaches von $m$ wächst. Versteht man aber unter $t$ den kleinsten positiven Rest von $n$ und setzt
\[
n=t+ml,
\]
so ist $\chi(n)=\chi(t)$, und es folgt:
\[
X=\int_0^1\sum_t\chi(t)x^{t-1}\sum_{l=0}^{\infty}x^{ml}\,dx.
\]
Setzt man jetzt
\begin{equation}
f(x)=\sum_t\chi(t)x^t,
\tag{2}
\end{equation}
so ist $f(x)$ eine ganze Function von $x$ vom Grade $m-1$, und zugleich ist wegen der Relation $\sum\chi(t)=0$:
\begin{equation}
f(0)=0,\qquad f(1)=0,
\tag{3}
\end{equation}
also $f(x)$ durch $x(1-x)$ theilbar. Für $X$ ergiebt sich dann nach der Summenformel für die geometrische Reihe $\sum x^{ml}$:
\begin{equation}
X=\int_0^1\frac{f(x)\,dx}{x(1-x^m)}.
\tag{4}
\end{equation}

Um ein solches Integral zu finden, schreibt die Integralrechnung vor, den rationalen Bruch unter dem Integralzeichen in Partialbrüche zu zerlegen. Diese Partialbruch-Zerlegung ergiebt aber, wegen (3), wenn
\[
r=e^{\frac{2\pi i}{m}}
\]
gesetzt wird:
\[
\frac{f(x)}{x(1-x^m)}
=-\frac1m\sum_{s=1}^{m-1}\frac{f(r^s)}{x-r^s}
\]
[Bd. I, \S~29, (11)], und nun haben wir weiter nach bekannten

\srcpage{710}
elementaren Sätzen der Integralrechnung, wenn $\alpha$ irgend einen Winkel zwischen $0$ und $2\pi$ bedeutet,
\[
\int_0^1\frac{dx}{x-e^{i\alpha}}
=\frac12\log\left(4\sin^2\frac{\alpha}{2}\right)
+i\,\frac{\pi-\alpha}{2}.
\]
Hieraus ergiebt sich nach (4):
\begin{equation}
X=-\frac1m\sum_{s=1}^{m-1}f(r^s)
\left[\frac12\log\left(4\left(\sin\frac{s\pi}{m}\right)^2\right)-\frac{i\pi(m-2s)}{2m}\right].
\tag{5}
\end{equation}
Hierin ist nach (2)
\begin{equation}
f(r^s)=\sum_t\chi(t)r^{st},
\tag{6}
\end{equation}
und daraus ergiebt sich [wegen (3)]:
\begin{equation}
\sum_{s=1}^{m-1}f(r^s)=\sum_{s=0}^{m-1}f(r^s)=0,
\tag{7}
\end{equation}
da
\[
\sum_{s=0}^{m-1}r^{st}=0
\]
ist. Hiernach vereinfacht sich die Formel (5) noch etwas und ergiebt
\begin{equation}
X=-\frac1m\sum_{s=1}^{m-1}f(r^s)
\left[\frac12\log\left(\sin\frac{s\pi}{m}\right)^2+\frac{i\pi s}{m}\right],
\tag{8}
\end{equation}
wofür man auch, wenn man $s$ durch $m-s$ ersetzt und wieder (7) benutzt,
\begin{equation}
X=-\frac1m\sum_{s=1}^{m-1}f(r^{-s})
\left[\frac12\log\left(\sin\frac{s\pi}{m}\right)^2-\frac{i\pi s}{m}\right]
\tag{9}
\end{equation}
setzen kann.

Aus (6) folgt aber
\[
f(r^{-s})=\sum_t\chi(t)r^{-st},
\]
und die nach $t$ genommene Summe ändert sich nicht, wenn $t$ durch $m-t$ ersetzt wird. Daraus folgt:
\[
f(r^{-s})=\chi(-1)f(r^s).
\]
Der Factor $\chi(-1)$ hat den Werth $+1$ oder $-1$, da sein Quadrat $=+1$ ist. Wir unterscheiden daher jetzt zwei Arten von Charakteren $\chi_1(n),\chi_2(n)$, so dass
\begin{equation}
\chi_1(-1)=1,\qquad \chi_2(-1)=-1
\tag{10}
\end{equation}

\srcpage{711}
ist, und danach sind auch die Functionen $f$ in $f_1$ und $f_2$ zu unterscheiden, so dass
\begin{equation}
f_1(r^{-s})=f_1(r^s),\qquad f_2(r^{-s})=-f_2(r^s)
\tag{11}
\end{equation}
wird. Ebenso unterscheiden wir die Ausdrücke (8), (9) als $X_1$ und $X_2$, und es ergiebt sich, wenn man beide Ausdrücke addirt, nach (11):
\begin{equation}
\begin{aligned}
X_1&=-\frac1{2m}\sum_{s=1}^{m-1}f_1(r^s)\log\left(\sin\frac{s\pi}{m}\right)^2,\\
X_2&=-\frac{i\pi}{m^2}\sum_{s=1}^{m-1}s\,f_2(r^s).
\end{aligned}
\tag{12}
\end{equation}


\section*{\S\,195. Ueber die Classenzahl in dem in $\Omega_m$ enthaltenen reellen Körper.}

\srcpage{711}
Unsere allgemeinen Principien können auch angewandt werden, um die Classenzahlen in den Theilern des Körpers $\Omega_m$ zu bestimmen.

Insbesondere sind die Formeln (6), (7), \S~192 anwendbar, wenn wir die Grade der Primideale in einem solchen Theiler kennen. Die Frage nach dem Verhältniss der Classenzahlen in den Theilern eines Körpers $\Omega_m$ zu der Classenzahl in $\Omega_m$ selbst ist von grossem Interesse und soll hier für den Fall behandelt werden, dass es sich um den in $\Omega_m$ enthaltenen reellen Körper $H_m$ handelt, den wir ja schon im \S~176 f. untersucht haben\footnote{Vgl. Kummer, „Bestimmung der Anzahl u.s.f.“. Crelle's Journ., Bd. 40 (1849).}. Es hat sich dort ergeben, dass $q$ im Körper $H_m$ die $\frac12\mu$te Potenz einer Primzahl ersten Grades ist, dass die anderen Primzahlen $p$ in $e_1$ Primfactoren $f_1$ten Grades zerfallen, so dass
\begin{equation}
f_1e_1=\frac12\mu,
\tag{1}
\end{equation}
und dass zwei Arten von Primzahlen $p_1,p_2$ zu unterscheiden sind, so dass bei den Primzahlen erster Art $f_1=\frac12 f$, $e_1=e$, bei denen der zweiten Art $f_1=f$, $e_1=\frac12 e$ ist.

Für den Körper $H_m$ erhalten wir demnach aus \S~192, (6), (7) die Classenzahl $h_1$ nach der Formel:
\srcpage{712}
\begin{equation}
g_1h_1=\lim_{s=1}(s-1)\Phi_1(s),
\tag{2}
\end{equation}
\begin{equation}
\Phi_1(s)=\frac1{1-q^{-s}}\prod_{p_1}\frac1{(1-p_1^{-\frac12sf})^e}\prod_{p_2}\frac1{(1-p_2^{-sf})^{\frac12e}},
\tag{3}
\end{equation}
worin sich die beiden Productzeichen auf alle Primzahlen $p_1$ und $p_2$ beziehen. Mit Benutzung der oben erklärten Zeichen $f_1,e_1$ kann man dafür auch setzen:
\begin{equation}
\Phi_1(s)=\frac1{1-q^{-s}}\prod_p\frac1{(1-p^{-sf_1})^{e_1}},
\tag{4}
\end{equation}
und das Productzeichen auf alle Primzahlen $p$ erstrecken.

Die Bedeutung von $g_1$ ergiebt sich aus \S~190, (11).

Nun wollen wir aus der Gruppe $C$ der Charaktere $\chi$ einen Theiler $C_1$ aussondern, der aus allen den Charakteren $\chi_1$ bestehen soll, für die
\begin{equation}
\chi_1(-1)=+1
\tag{5}
\end{equation}
ist. Zunächst ist klar, dass diese Charaktere $\chi_1$ eine Gruppe $C_1$ bilden. Diese Gruppe ist aber nicht mit $C$ identisch, weil es gewiss einen Charakter $\chi_0$ giebt, für den $\chi_0(-1)=-1$ ist. Wir brauchen, um einen solchen zu erhalten, nur in \S~193, (6) für $\Theta$ eine primitive $\mu$te Einheitswurzel, oder in \S~193, (8) $\eps=-1$ zu setzen. Dann bildet die Gesammtheit der Charaktere $\chi_0\chi_1$, die alle nicht zu $C_1$ gehören, eine Nebengruppe $C_2$, deren Elemente mit $\chi_2$ bezeichnet werden mögen. Jeder Charakter $\chi$ ist also dann entweder in $C_1$ oder in $C_2$ enthalten, denn wenn $\chi$ nicht in $C_1$ vorkommt, so kommt $\chi_0^{-1}\chi$ darin vor, und folglich $\chi$ in $C_2$. Es ist daher $C_1$ ein Theiler von $C$ vom Index 2.

Die Charaktere $\chi_1$ kann man auf folgende Art bilden:

1) Ist $m$ ungerade, so setzen wir, wie in \S~193,
\[
n\equiv c^\gamma\pmod m,
\]
und lassen in
\begin{equation}
\chi_1(n)=\Theta^\gamma
\tag{6}
\end{equation}
$\Theta$ die Gesammtheit der $\frac12\mu$ten Einheitswurzeln durchlaufen.

2) Ist $m$ gerade, so setzen wir
\[
n\equiv(-1)^\alpha5^\beta,
\]
und erhalten
\begin{equation}
\chi_1(n)=\Theta^\beta,
\tag{7}
\end{equation}
worin $\Theta$ wieder die sämmtlichen $\frac12\mu$ten Einheitswurzeln durchläuft.

\srcpage{713}
Beides ergiebt sich leicht dadurch, dass für $n=-1$ im ersten Falle $\gamma=\frac12\mu$, im zweiten $\alpha=1$, $\beta=0$ ist, also beide Male $\chi_1(-1)=1$ ist, und dass jede der Formeln (6), (7) genau $\frac12\mu$ verschiedene Charaktere darstellt.

Worauf es nun wesentlich ankommt, ist Folgendes:
\begin{center}
\textbf{Ist $p$ irgend eine von $q$ verschiedene Primzahl, so erhält man in der Form $\chi_1(p)$, wenn $\chi_1$ die Gruppe $C_1$ durchläuft, jede $f_1$te Einheitswurzel, und jede gleich oft, also $e_1$ mal.}
\end{center}
Dass alle $\chi_1(p)$ Einheitswurzeln vom Grade $f_1$ sind, ist klar, denn es ist immer (vgl. \S~177) $p^{f_1}\equiv\pm1\pmod m$, und daher
\[
[\chi_1(p)]^{f_1}=\chi_1(\pm1)=1.
\]
Wenn aber unter den $\chi_1(p)$ eine primitive $f_1$te Einheitswurzel vorkommt, so kommen alle anderen $f_1$ten Einheitswurzeln auch darunter vor, und gleich oft, nämlich so oft wie die Einheitswurzel $1$.

Dies folgt unmittelbar aus der Gruppennatur der $\chi_1$. Es ist also nun noch zu zeigen, dass unter den $\chi_1(p)$ eine primitive $f_1$te Einheitswurzel vorkommt.

Diese Möglichkeit ergiebt sich aber aus der Bestimmung von $\chi_1$ sehr einfach. Ist zunächst $\Theta$ in (6) bei ungeradem $m$ eine primitive Einheitswurzel vom Grade $\frac12\mu$, und ist $\gamma$ der Index von $p$, so ist eine Potenz von $\Theta^\gamma$, $\Theta^{\gamma\lambda}$ dann und nur dann gleich $1$, wenn
\[
2\gamma\lambda\equiv0\pmod\mu.
\]
Nun ist $\mu=ef$ und $e$ der grösste gemeinschaftliche Theiler von $\gamma$ und $\mu$; und daher wird diese Bedingung $2\lambda\equiv0\pmod f$. Es ist aber nach \S~177, 3, $f_1=f$ oder $=\frac12f$, je nachdem $f$ ungerade oder gerade ist, und folglich muss $\lambda\equiv0\pmod {f_1}$ sein. Also ist $\Theta^\gamma$ primitive $f_1$te Einheitswurzel.

Ist aber $m$ gerade, so nehmen wir in (7) gleichfalls für $\Theta$ eine primitive $\frac12\mu$te Einheitswurzel und setzen
\[
p\equiv(-1)^\alpha5^\beta\pmod m,
\]
so dass $f$ die kleinste positive Zahl ist, die den Congruenzen
\[
\alpha f\equiv0\pmod2,
\qquad
\beta f\equiv0\pmod{\tfrac12\mu}
\]
genügt. Es ist nun $\Theta^{\beta\lambda}$ nur dann gleich $1$, wenn $\beta\lambda\equiv0\pmod{\frac12\mu}$ ist. Wenn nicht $\lambda=1$, also $\beta\not\equiv0$ ist, so ist der kleinste Werth, den $\lambda$ haben kann, gerade, und daher ist dieser kleinste Werth von $\lambda=f$.

\srcpage{714}
Wenn aber $\lambda=1$ ist, so muss $\beta=0$ sein, und es sind noch zwei Fälle möglich:
\[
\begin{array}{ll}
entweder & f=1,\quad \alpha=0,\quad p\equiv1\pmod m,\quad \lambda=f,\\[2mm]
oder & f=2,\quad \alpha=1,\quad p\equiv-1\pmod m,\quad \lambda=\frac12f.
\end{array}
\]
Nur in diesem letzten Falle gehört $p$ zur ersten Art, und es folgt also, dass auch hier allgemein $\Theta^\beta$ eine primitive $f_1$te Einheitswurzel ist. Hiermit ist dann der Satz bewiesen.

Dieser Satz gestattet uns die Anwendung der Formel \S~193, (9), die hier ergiebt:
\begin{equation}
(1-p^{-sf_1})^{e_1}=\prod_{\chi_1}[1-\chi_1(p)p^{-s}],
\tag{8}
\end{equation}
worin sich das Product auf alle Charaktere $\chi_1$ der Gruppe $C_1$ erstreckt, und nun lässt sich die Function $\Phi_1(s)$ ebenso wie $\Phi(s)$ im \S~193 weiter entwickeln. Man erhält auf demselben Wege, wie dort, die Formel (11),
\begin{equation}
\Phi_1(s)=\frac1{1-q^{-s}}\prod_{\chi_1}\sum_{n=1}^{\infty}\frac{\chi_1(n)}{n^s},
\tag{9}
\end{equation}
und wenn man zur Grenze $s=1$ übergeht,
\begin{equation}
g_1h_1=\prod_{\chi_1}\sum_{n=1}^{\infty}\frac{\chi_1(n)}{n},
\tag{10}
\end{equation}
worin aber jetzt das Product auf alle Charaktere der Gruppe $C_1$ mit Ausnahme des Hauptcharakters zu erstrecken ist. Dieses Product ist also ein Theil des Productes, durch welches im \S~193 die Zahl $gh$ ausgedrückt ist.

Die Zahlen $g,g_1$ ergeben sich aus dem im \S~190, (11) angegebenen allgemeinen Ausdrucke:
\begin{equation}
\frac{2^\nu\pi^{n-\nu}L}{w\sqrt{\pm\Delta}}.
\tag{11}
\end{equation}
Da, wie im \S~178 nachgewiesen ist, die Einheiten in $\Omega_m$ und $H_m$, von den Einheitswurzeln abgesehen, dieselben sind, so ist auch ein Fundamentalsystem von Einheiten in $H_m$ zugleich ein Fundamentalsystem in $\Omega_m$.

Bedeutet $\eps_1,\eps_2,\ldots,\eps_{\nu-1}$ daher ein Fundamentalsystem reeller Einheiten, und sind $\eps_{i,1},\eps_{i,2},\ldots,\eps_{i,\nu}$ die conjugirten Einheiten zu $\eps_i$, so sind die conjugirten Logarithmen im Körper $H_m$:
\[
\log|\eps_{i,1}|,\quad\log|\eps_{i,2}|,\quad\ldots,\quad\log|\eps_{i,\nu}|,
\]
dagegen im Körper $\Omega_m$, in dem nur conjugirt imaginäre Paare vorkommen (\S~185),
\[
2\log|\eps_{i,1}|,\quad2\log|\eps_{i,2}|,\quad\ldots,\quad2\log|\eps_{i,\nu}|.
\]

\srcpage{715}
Bezeichnen wir daher mit
\begin{equation}
E=\pm\begin{vmatrix}
\log|\eps_{1,1}|&\log|\eps_{1,2}|&\cdots&\log|\eps_{1,\nu-1}|\\
\log|\eps_{2,1}|&\log|\eps_{2,2}|&\cdots&\log|\eps_{2,\nu-1}|\\
\cdots&\cdots&\cdots&\cdots\\
\log|\eps_{\nu-1,1}|&\log|\eps_{\nu-1,2}|&\cdots&\log|\eps_{\nu-1,\nu-1}|
\end{vmatrix}
\tag{12}
\end{equation}
den Regulator des Körpers $H_m$ (\S~187), so ist bei der Anwendung des Ausdruckes (11) im Körper $\Omega_m$
\[
L=2^{\nu-1}E,
\]
und im Körper $H_m$
\[
L=E
\]
zu setzen. Es ist ferner
\[
\begin{array}{lll}
in\ \Omega_m: & n=\mu, & \nu=\frac12\mu,\\
in\ H_m: & n=\frac12\mu, & \nu=\frac12\mu.
\end{array}
\]
Im Körper $H_m$ sind nur die zwei Einheitswurzeln $\pm1$ enthalten. In $\Omega_m$ sind die mit positivem und negativem Zeichen genommenen Potenzen von $r$, sonst aber keine Einheitswurzeln enthalten (\S~175), und bei der Zählung ist noch zu beachten, dass bei geradem $m$ unter den Potenzen von $r$ auch $-1$ enthalten ist, bei ungeradem $m$ nicht. Daher haben wir
\[
\begin{array}{lll}
in\ \Omega_m: & w=2m, & m\text{ ungerade},\\
& w=m, & m\text{ gerade},\\
in\ H_m: & w=2.&
\end{array}
\]
Endlich haben wir noch, wenn $\Delta,\Delta_1$ die Grundzahlen der Körper $\Omega_m,H_m$ sind (\S~176):
\[
\pm\Delta=q\Delta_1^2\quad m\text{ ungerade},
\qquad
\pm\Delta=4\Delta_1^2\quad m\text{ gerade},
\]
und es folgt, wenn wir der Einfachheit wegen $\nu$ für $\frac12\mu$ setzen:
\begin{align}
g&=\frac{2^{2(\nu-1)}\pi^\nu E}{m\sqrt q\,\Delta_1}, &&m\text{ ungerade},\notag\\
g&=\frac{2^{2(\nu-1)}\pi^\nu E}{m\Delta_1}, &&m\text{ gerade},\notag\\
g_1&=\frac{2^{\nu-1}E}{\sqrt{\Delta_1}},\tag{13}
\end{align}
woraus
\begin{align}
g&=\frac{2^{\nu-1}\pi^\nu g_1}{m\sqrt{q\Delta_1}}, &&m\text{ ungerade},\notag\\
g&=\frac{2^{\nu-1}\pi^\nu g_1}{m\sqrt{\Delta_1}}, &&m\text{ gerade}.\tag{14}
\end{align}

\srcpage{716}
Wenn wir die Formel (10) mit der Formel \S~193, (15) verbinden und die im \S~194 eingeführte Bezeichnung
\[
X_1=\sum\frac{\chi_1(n)}{n},
\qquad
X_2=\sum\frac{\chi_2(n)}{n}
\]
gebrauchen, so ergiebt sich
\[
gh=g_1h_1\prod X_2.
\]
Ersetzt man $g$ durch seinen Werth (14), so folgt für die Classenzahl
\begin{equation}
h=kh_1,
\tag{15}
\end{equation}
wenn
\begin{equation}
\begin{aligned}
k&=\frac{m\sqrt{q\Delta_1}}{2^{\nu-1}\pi^\nu}\prod X_2 &&\text{bei ungeradem }m,\\
&=\frac{m\sqrt{\Delta_1}}{2^{\nu-1}\pi^\nu}\prod X_2 &&\text{bei geradem }m,
\end{aligned}
\tag{16}
\end{equation}
und worin [nach (10), (13)]:
\begin{equation}
h_1=\frac{\sqrt{\Delta_1}}{2^{\nu-1}E}\prod X_1
\tag{17}
\end{equation}
die Classenzahl des Körpers $H_m$ ist. Im letzteren Ausdrucke erstreckt sich das Product auf alle Charaktere $\chi_1$ der Gruppe $C_1$ mit Ausnahme des Hauptcharakters.

Dass $k$ eine ganze Zahl ist, und daher $h$ ein Vielfaches von $h_1$, wird sich bald zeigen. Die Zahlen $k$ und $h_1$ heissen der erste und zweite Factor der Classenzahl.

In dem Falle, dass
\[
m=2^\varkappa
\]
eine Potenz von 2 ist, haben wir nach \S~176, (7):
\[
\Delta_1=2^{(\varkappa-1)2^{\varkappa-2}-1},
\]
und wenn wir diesen Werth einsetzen, so können wir für die vorstehenden Formeln in diesem Falle schreiben:
\begin{equation}
\begin{aligned}
k&=\frac{\sqrt2\,2^{(\varkappa-3)2^{\varkappa-3}}2^\varkappa}{\pi^\nu}\prod X_2,\\
h_1&=\frac{\sqrt2\,2^{(\varkappa-3)2^{\varkappa-3}}}{E}\prod X_1.
\end{aligned}
\tag{18}
\end{equation}

\section*{\S\,196. Classenzahl im Körper der achten Einheitswurzeln.}

\srcpage{717}
Wir wollen, einerseits zur Veranschaulichung der bisherigen Resultate, andererseits um für die später anzuwendende vollständige Induction eine Basis zu gewinnen, den Fall $m=8$, $\mu=4$, $\nu=2$ zunächst behandeln.

Hier haben wir ausser dem Hauptcharakter nur drei Charaktere, nämlich, wenn
\[
n\equiv(-1)^\alpha5^\beta\pmod8
\]
ist,
\[
\chi_1(n)=(-1)^\beta,
\qquad
\chi_2(n)=(-1)^\alpha,
\qquad
\chi_3(n)=(-1)^{\alpha+\beta}.
\]
Für $n=-1$ ist $\alpha=1$, $\beta=0$, und folglich gehört $\chi_1$ zur ersten, $\chi_2,\chi_3$ zur zweiten Art [\S~195, (5)], und es ergeben sich für die vier Zahlen $n=1,3,5,7$ folgende Werthe dieser Charaktere:
\[
\begin{array}{c|rrrr}
&1&3&5&7\\\hline
\chi_1&+1&-1&-1&+1\\
\chi_2&+1&-1&+1&-1\\
\chi_3&+1&+1&-1&-1.
\end{array}
\]
Daraus erhält man die drei Functionen $f_1,f_2,f_3$ [\S~194, (2)]:
\[
\begin{aligned}
f_1(x)&=x-x^3-x^5+x^7=x(1-x^2)(1-x^4),\\
f_2(x)&=x-x^3+x^5-x^7=x(1-x^2)(1+x^4),\\
f_3(x)&=x+x^3-x^5-x^7=x(1+x^2)(1-x^4).
\end{aligned}
\]
Für $r$ können wir setzen:
\[
r=\frac{1+i}{\sqrt2},\quad r^2=i,
\quad r^3=-\frac{1-i}{\sqrt2},\quad r^4=-1,
\]
und danach ergiebt sich:
\[
\begin{array}{lll}
f_1(r)=2\sqrt2,& f_2(r)=0,& f_3(r)=2\sqrt2\,i,\\
f_1(r^2)=0,& f_2(r^2)=4i,& f_3(r^2)=0,\\
f_1(r^3)=-2\sqrt2,& f_2(r^3)=0,& f_3(r^3)=2\sqrt2\,i.
\end{array}
\]
Für $x=r^4$ verschwinden alle drei Functionen $f(x)$, und für $x=r^5,r^6,r^7$ erhält $f_1(x)$ dieselben, $f_2(x)$ und $f_3(x)$ die entgegengesetzten Werthe, wie für $x=r^3,r^2,r$.

\srcpage{718}
Demnach ergiebt sich aus \S~194, (12) mit Rücksicht auf die trigonometrische Formel
\[
\sin\frac{3\pi}{8}=\cos\frac\pi8
\]
\[
X_1=-\frac1{\sqrt2}\log\tan\frac\pi8,
\qquad
X_2=-\frac\pi4,
\qquad
X_3=-\frac\pi{2\sqrt2}.
\]
Nun ist
\[
\tau=\frac{r-1}{r^2(r+1)}=\tan\frac\pi8=\sqrt2-1
\]
eine Einheit des Körpers $H_8$, der hier nichts Anderes ist, als der aus $\sqrt2$ entspringende quadratische Körper, und die Einheiten darin erhält man also aus den Lösungen der Pell'schen Gleichung
\[
u^2-2v^2=\pm1
\]
in der Form $u+v\sqrt2$.

Die kleinste positive Lösung dieser Pell'schen Gleichung ist aber $u=1$, $v=1$, und folglich findet man nach Bd. I, \S~128 alle Einheiten in $H_8$ in der Form $\pm(1+\sqrt2)^a$, worin $a$ ein positiver oder negativer ganzzahliger Exponent ist; $\tau$ ergiebt sich hieraus für $a=-1$, und daher sind alle Einheiten auch in der Form
\[
\pm\tau^a
\]
enthalten. Es ist also $\tau$ eine fundamentale Einheit, und die in den Formeln des vorigen Paragraphen vorkommende Determinante $E$ [\S~195, (12)] wird hier (da $\tau$ ein echter Bruch ist):
\[
E=-\log\tau.
\]
Man erhält also aus \S~195, (15), (18):
\[
h_1=1,\quad k=1,\quad h=1.
\]
Der Körper $\Omega_8$ ist also ein einclassiger, d. h. ein solcher, in dem die Zerlegung der ganzen Zahlen in Primfactoren nach denselben Gesetzen, wie im Körper der rationalen Zahlen, möglich ist.

Es hat sich dabei zugleich ergeben, dass alle Einheiten im Körper $H_8$ in der Form $\pm\tau^a$ darstellbar sind. Die conjugirten Werthe der Einheit $\tau$ sind aber
\[
\tau_1=\tan\frac\pi8,
\qquad
\tau_3=-\tan\frac{3\pi}{8}=-\tau_1^{-1},
\]
und haben also verschiedene Zeichen. Wir schliessen daraus für diesen Fall auf den Satz:
\begin{center}
\textbf{Eine Einheit in $H_8$, die mit ihren Conjugirten von einerlei Zeichen ist, ist, vom Zeichen abgesehen, das Quadrat einer Einheit.}
\end{center}

\section*{\S\,197. Recurrente Berechnung der Classenzahl im Körper $\Omega_m$, wenn $m$ eine Potenz von 2 ist.}

\srcpage{719}
Zur allgemeinen Bestimmung der Classenzahl im Körper $\Omega_m$ unter der Voraussetzung, dass $m$ irgend eine Potenz von 2 und grösser als 8 ist, wenden wir ein recurrirendes Verfahren an. Es sei also
\begin{equation}
m=2^\varkappa,
\quad \mu=2^{\varkappa-1},
\quad \nu=2^{\varkappa-2},
\quad \nu'=2^{\varkappa-3}.
\tag{1}
\end{equation}
Neben dem Körper $\Omega_m$ betrachten wir den Körper $\Omega_\mu$, der ein Theiler von $\Omega_m$ ist, und den wir hier auch mit $\Omega'_m$ bezeichnen wollen. Es sei $h'$ die Classenzahl im Körper $\Omega'_m$, d. h. die Zahl, die sich aus $h$ ergiebt, wenn $\varkappa$ in $\varkappa-1$ verwandelt wird, und wir setzen
\begin{equation}
h=h'H.
\tag{2}
\end{equation}
Setzen wir $h'$ als schon bekannt voraus, so kommt es nur noch auf die Berechnung von $H$ an, was, wie sich zeigen wird, eine ganze Zahl ist. Indem wir $h$ und $h'$ wie im \S~195 zerlegen, setzen wir
\begin{equation}
h=kh_1,
\qquad
h'=k'h'_1,
\qquad
H=AB,
\tag{3}
\end{equation}
\begin{equation}
k=k'A,
\qquad
h_1=h'_1B,
\tag{4}
\end{equation}
worin $k'$ aus $k$ und $h'_1$ aus $h_1$ durch die Vertauschung von $\varkappa$ mit $\varkappa-1$ hervorgeht. Wir wollen überhaupt jetzt durch Accente an den Buchstaben andeuten, dass $\varkappa-1$ statt $\varkappa$ gesetzt sein soll.

Nach \S~176, (7) ist dann
\[
\Delta_1=2^{(\varkappa-1)2^{\varkappa-2}-1},
\qquad
\Delta'_1=2^{(\varkappa-2)2^{\varkappa-3}-1},
\]
und folglich
\begin{equation}
\Delta_1=2^{\varkappa2^{\varkappa-3}}\Delta'_1.
\tag{5}
\end{equation}
Hat $E$ die Bedeutung \S~195, (12), und geht $E'$ aus $E$ hervor, wenn $\varkappa$ durch $\varkappa-1$ ersetzt wird, so setzen wir noch
\begin{equation}
E=E'D,
\tag{6}
\end{equation}

\srcpage{720}
wodurch $D$ als eine von Null verschiedene positive Zahl definirt ist. Was die Summen
\[
X=\sum\frac{\chi(n)}{n}
\]
betrifft, so ist daran zu erinnern, dass $\chi(n)$ [nach \S~193, (8)] definirt ist durch
\begin{equation}
n\equiv(-1)^\alpha5^\beta\pmod m,
\qquad
\chi(n)=(\pm1)^\alpha\Theta^\beta,
\tag{7}
\end{equation}
worin $\Theta$ eine $\nu$te Einheitswurzel bedeutet.

Unter den $\nu$ten Einheitswurzeln sind aber auch die $\nu'$ten Einheitswurzeln enthalten, und unter den Charakteren $\chi$ auch die Charaktere für den Körper $\Omega'_m$. Wir unterscheiden demnach primitive und imprimitive Charaktere des Körpers $\Omega_m$, indem wir die dem Körper $\Omega_m$ eigenthümlichen Charaktere $\chi$, d. h. die, in denen $\Theta$ eine primitive $\nu$te Einheitswurzel ist, als primitiv bezeichnen.

Unter den Summen $X$ kommen auch alle zu dem Körper $\Omega'_m$ gehörigen Summen vor:
\[
X'=\sum\frac{\chi'(n)}{n}.
\]
Wenn man nun in (4) für $k,k'$ und $h_1,h'_1$ die im \S~195, (16), (17) gebildeten Ausdrücke einsetzt, und dann die Summen $X'$ weghebt, so ergiebt sich:
\begin{equation}
\begin{aligned}
\pi^{\nu'}A&=2\cdot2^{2^{\varkappa-4}(\varkappa-2)}\prod X_2,\\
DB&=2^{2^{\varkappa-4}(\varkappa-2)}\prod X_1,
\end{aligned}
\tag{8}
\end{equation}
worin sich jetzt aber die Producte $\prod$ nur noch auf die mit den primitiven Charakteren $\chi_1,\chi_2$ gebildeten Summen $X_1,X_2$ erstrecken.

In den im \S~194, (12) gegebenen Ausdrücken für die $X$:
\begin{equation}
\begin{aligned}
X_1&=-\frac1{2m}\sum_{s=1}^{m-1} f_1(r^s)\log\left(\sin\frac{s\pi}{m}\right)^2,\\
X_2&=-\frac{i\pi}{m^2}\sum_{s=1}^{m-1}s f_2(r^s)
\end{aligned}
\tag{9}
\end{equation}
treten nun, unter der Voraussetzung, dass $m$ eine Potenz von 2 und dass $\chi_1,\chi_2$ primitive Charaktere sind, bedeutende Vereinfachungen ein.

\srcpage{721}
Die Zahl $1+\mu$ hat die Indices $\alpha=0$, $\beta=\nu'$, und folglich ist nach (7)
\[
\chi(1+\mu)=\Theta^{\nu'}=-1,
\]
wenn $\chi$ ein primitiver Charakter ist. Ferner ist für ein ungerades $n$:
\[
n(1+\mu)\equiv n+\mu\pmod m,
\]
und folglich
\[
\chi(n+\mu)=-\chi(n).
\]
Wenn wir nun die Function $f(r^s)$ betrachten:
\begin{equation}
f(r^s)=\sum_n\chi(n)r^{sn},
\tag{10}
\end{equation}
worin $n$ die ungeraden Zahlen eines vollen Restsystems für den Modul $m$ durchläuft, so erhalten wir, da wir in (10) unter dem Summenzeichen $n$ durch $n+\mu$ ersetzen dürfen:
\[
f(r^s)=-r^{\mu s}\sum\chi(n)r^{sn},
\]
und da $r^\mu=-1$ ist:
\[
f(r^s)=-(-1)^sf(r^s).
\]
Es ist also
\begin{equation}
f(r^s)=0,\quad \text{wenn }s\text{ gerade ist,}
\tag{11}
\end{equation}
und demnach können wir in den Formeln (9) den Summationsbuchstaben $s$ auf die ungeraden Zahlen beschränken, die kleiner als $m$ sind.

Es ist aber nach (10):
\[
f(r^s)=\chi(s)^{-1}\sum_n\chi(ns)r^{sn},
\]
und wenn man $sn$ durch $n$ ersetzt, was bei ungeradem $s$ gestattet ist:
\begin{equation}
f(r^s)=\chi(s)^{-1}f(r),\quad \text{wenn }s\text{ ungerade ist.}
\tag{12}
\end{equation}
Die Function $f(r)$ hat, wenn $\alpha,\beta$ die Indices von $n$ sind, den Ausdruck
\begin{equation}
f(r)=\sum(\pm1)^\alpha\Theta^\beta r^n,
\tag{13}
\end{equation}
und ist also mit der im \S~17 dieses Bandes betrachteten Kreistheilungsresolvente
\[
(\pm1,\Theta,r)
\]
gleichbedeutend, für die im \S~17, (17) die Relation aufgestellt war:
\begin{equation}
(\pm1,\Theta,r)(\pm1,\Theta^{-1},r)=\pm m.
\tag{14}
\end{equation}
Hier gelten durchweg die oberen Zeichen für die Summen $X_1$, die unteren für die Summen $X_2$ [\S~195, (7)].

Wir theilen jetzt die Reihe der in (9) vorkommenden Zahlen $s$ in vier Theile. Es soll $t$ ein Zeichen sein, welches die
\srcpage{722}
Reihe der positiven ungeraden Zahlen von 1 bis $\nu-1$ durchläuft; dann durchläuft das ungerade $s$ die vier Zahlenreihen:
\[
t,
\quad t+\mu,
\quad m-t,
\quad m-t-\mu,
\]
und es ist:
\[
\begin{gathered}
\chi(t+\mu)=-\chi(t),\\
\chi_1(m-t)=\chi_1(t),\qquad \chi_1(m-\mu-t)=-\chi_1(t),\\
\chi_2(m-t)=-\chi_2(t),\qquad \chi_2(m-\mu-t)=\chi_2(t).
\end{gathered}
\]
Nach (12) ergiebt sich hieraus:
\[
\begin{gathered}
f(r^{t+\mu})=-f(r^t),\\
f_1(r^{m-t})=f_1(r^t),\qquad f_1(r^{m-\mu-t})=-f_1(r^t),\\
f_2(r^{m-t})=-f_2(r^t),\qquad f_2(r^{m-\mu-t})=f_2(r^t).
\end{gathered}
\]
Theilt man demnach die Summen (9) in je vier Theile und benutzt noch die trigonometrische Gleichung
\[
\sin\frac{(t+\mu)\pi}{m}=\cos\frac{t\pi}{m},
\]
so ergiebt sich
\[
X_1=-\frac1m\sum_t f_1(r^t)\log\left(\tan\frac{t\pi}{m}\right)^2,
\qquad
X_2=\frac{i\pi}{m}\sum_t f_2(r^t).
\]
In dieser Summe ist
\[
\frac{t\pi}{m}<\frac\pi4,
\qquad
\tan\frac{t\pi}{m}>0,
\]
und folglich erhalten wir mit Rücksicht auf (12), (13), wenn wir $\Theta^{-1}$ für $\Theta$ setzen, wodurch $\chi(t)^{-1}$ in $\chi(t)$ übergeht:
\begin{equation}
X_1=-\frac2m(+1,\Theta^{-1},r)
\sum_{1\le t\le\nu-1}\chi_1(t)\log\tan\frac{t\pi}{m},
\tag{15}
\end{equation}
\begin{equation}
X_2=\frac{i\pi}{m}(-1,\Theta^{-1},r)
\sum_{1\le t\le\nu-1}\chi_2(t).
\tag{16}
\end{equation}

\section*{\S\,198. Der Classenzahlfactor $A$.}

Um nun zunächst $A$ zu berechnen, haben wir den Ausdruck (16) für $X_2$ in die erste Formel (8) des vorigen Paragraphen einzusetzen. Diese Formel können wir auch so darstellen:
\begin{equation}
\pi^{\nu'}A=2m^{1/4}2^{-\nu'}\prod X_2.
\tag{1}
\end{equation}

\srcpage{723}
Bezeichnen wir mit $\alpha,\beta$ die Indices von $t$, so ist $\chi_2(t)=(-1)^\alpha\Theta^\beta$;
\begin{equation}
\sum_t\chi_2(t)=\sum(-1)^\alpha\Theta^\beta=\varphi(\Theta)
\tag{2}
\end{equation}
ist eine ganze Zahl des Körpers $\Omega_\nu$, und es wird:
\[
X_2=\frac{i\pi}{m}(-1,\Theta^{-1},r)\varphi(\Theta).
\]
Nun hat man für $\Theta$ alle Wurzeln der Gleichung
\begin{equation}
\Theta^{\nu'}+1=0
\tag{3}
\end{equation}
zu setzen, und erhält für $A$, mit Rücksicht auf \S~197, (14):
\begin{equation}
A=2^{1-\nu'}\prod_\Theta\varphi(\Theta).
\tag{4}
\end{equation}
Die Summe, durch die in (2) die Zahl $\varphi(\Theta)$ definirt ist, enthält $\nu'$ Glieder von der Form $(-1)^\alpha\Theta^\beta$, worin $\alpha,\beta$ so zu bestimmen sind, dass
\begin{equation}
(-1)^\alpha5^\beta\equiv t\pmod m
\tag{5}
\end{equation}
und $t$ zwischen 0 und $\nu$ liegt. Nun ist zunächst ersichtlich, dass unter den Exponenten $\beta$ nicht zwei nach dem Modul $\nu'$ congruente vorkommen können. Denn ist $\beta\equiv\beta'\pmod{\nu'}$, so ist $5^\beta\equiv5^{\beta'}\pmod\mu$, und aus
\[
(-1)^\alpha5^\beta\equiv t,
\qquad
(-1)^{\alpha'}5^{\beta'}\equiv t'\pmod m
\]
folgt:
\[
(-1)^\alpha t-(-1)^{\alpha'}t'\equiv0\pmod\mu.
\]
Da aber $t$ und $t'$ absolut kleiner als $\nu=\frac12\mu$ sind, so ist dies nur möglich, wenn $t=t'$, $\alpha=\alpha'$ ist. Demnach durchläuft $\beta$ in (2) ein volles Restsystem nach dem Modul $\nu'$.

Da aber $\beta$ nach dem Modul $\nu$ zu nehmen ist, so werden die nach (5) bestimmten $\beta$ theils kleiner, theils grösser als $\nu'$ ausfallen, und wenn wir also $\beta_1$ die Reihe der Zahlen $0,1,\ldots,\nu'-1$ durchlaufen lassen, so erhalten wir in (2) zweierlei Arten von Gliedern:
\begin{equation}
\begin{array}{ll}
1. & (-1)^\alpha\Theta^\beta=(-1)^\alpha\Theta^{\beta_1},\quad \beta=\beta_1<\nu',\\[1mm]
2. & (-1)^\alpha\Theta^\beta=-(-1)^\alpha\Theta^{\beta_1},\quad \beta=\beta_1+\nu'\geqq\nu'.
\end{array}
\tag{6}
\end{equation}
Nun ist nach (5) der absolut kleinste Rest von $(-1)^\alpha5^\beta$ nach dem Modul $m$ positiv, und wenn wir daher eine Zahl $c_\beta=\pm1$ so bestimmen, dass der absolut kleinste Rest von $c_\beta5^\beta$ positiv wird, so ist in dem Falle (6), 1,
\[
c_{\beta_1}=(-1)^\alpha.
\]

\srcpage{724}
Ferner ist $5^{\nu'}\equiv1\pmod\mu$, und da $5^{\nu'}$ nicht auch nach dem Modul $m$ mit 1 congruent ist, so folgt:
\begin{equation}
5^{\nu'}\equiv1+\mu\pmod m.
\tag{7}
\end{equation}
Demnach ist im Falle (6), 2,
\[
(-1)^\alpha5^\beta=(-1)^\alpha5^{\beta_1}5^{\nu'}\equiv(-1)^\alpha5^{\beta_1}+\mu\pmod m,
\]
und folglich ist hier der absolut kleinste Rest von $(-1)^\alpha5^{\beta_1}$ nach dem Modul $m$ negativ. Es ist also im Falle (6), 2,
\[
c_{\beta_1}=-(-1)^\alpha
\]
zu setzen.

Daraus ergiebt sich nach (2) (wenn wir wieder $\beta$ für $\beta_1$ schreiben):
\begin{equation}
\varphi(\Theta)=\sum_{\beta=0}^{\nu'-1}c_\beta\Theta^\beta
=c_0+c_1\Theta+c_2\Theta^2+\cdots+c_{\nu'-1}\Theta^{\nu'-1},
\tag{8}
\end{equation}
und hierin ist
\begin{equation}
c_\beta=+1\quad\text{oder}\quad=-1,
\tag{9}
\end{equation}
je nachdem der absolut kleinste Rest von $5^\beta$ positiv oder negativ ist. Hiernach ist also $\varphi(\Theta)$ eine ganze Zahl des Körpers $\Omega_\nu$.

Es kommt noch darauf an, das Verhalten dieser Zahl zu der Zahl 2, oder vielmehr zu dem in $\Omega_\nu$ enthaltenen Primfactor $(1-\Theta)$ von 2 zu ermitteln. Bilden wir zu diesem Zwecke
\begin{equation}
(1-\Theta)\varphi(\Theta)=2\psi(\Theta),
\tag{10}
\end{equation}
so ergiebt sich
\begin{equation}
\psi(\Theta)=\frac{c_0+c_{\nu'-1}}2+\frac{c_1-c_0}2\Theta+\frac{c_2-c_1}2\Theta^2+\cdots+
\frac{c_{\nu'-1}-c_{\nu'-2}}2\Theta^{\nu'-1}.
\tag{11}
\end{equation}
Die Coefficienten in diesem Ausdrucke sind alle $=0$ oder $=1$, und daher ist auch $\psi(\Theta)$ eine ganze Zahl des Körpers $\Omega_\nu$. Für diese Zahl ergiebt sich aber, wenn man $\Theta\equiv1$ setzt:
\begin{equation}
\psi(\Theta)\equiv c_{\nu'-1}\equiv\pm1\pmod{(1-\Theta)},
\tag{12}
\end{equation}
und folglich ist $\psi(\Theta)$ relativ prim zu 2.

Das in (4) vorkommende Product ist nichts Anderes, als die in Bezug auf den Körper $\Omega_\nu$ genommene Norm der Zahl $\varphi(\Theta)$, und wenn wir diese Norm mit $N$ bezeichnen, so ist nach \S~169
\[
N(1-\Theta)=2,
\]

\srcpage{725}
und folglich nach (10)
\[
N\varphi(\Theta)=2^{\nu'-1}N\psi(\Theta).
\]
Daraus ergiebt sich
\begin{equation}
A=N\psi(\Theta),
\tag{13}
\end{equation}
und hieraus folgt mit Rücksicht auf (12), dass $A$ eine ungerade ganze Zahl ist.

Bezeichnen wir die zu $m=2^\varkappa$ gehörige Zahl $A$ mit $A_\varkappa$, so ergiebt sich aus \S~197, (4) durch vollständige Induction der Ausdruck für den ersten Factor der Classenzahl:
\begin{equation}
k=A_4A_5\cdots A_\varkappa,
\tag{14}
\end{equation}
und daraus der Satz:
\begin{center}
\textbf{1. Der erste Factor der Classenzahl ist eine ungerade ganze Zahl.}
\end{center}
Die Zahl $A_\varkappa$ ist nach der Formel (10) für die ersten Werthe von $\varkappa$ nicht schwer zu berechnen.

Für $m=16$ findet man $\nu'=2$, $\beta=0,1$, $c_0=c_1=1$, folglich
\[
\psi(\Theta)=1,
\qquad
A_4=1.
\]
Für $m=32$, $\nu'=4$, $\beta=0,1,2,3$,
\[
c_0=1,
\quad c_1=1,
\quad c_2=-1,
\quad c_3=-1,
\]
\[
\psi(\Theta)=-\Theta^2,
\]
also auch $A_5=1$.

Für $m=64$, $\nu'=8$, $\beta=0,1,2,3,4,5,6,7$ sind die absolut kleinsten Reste von $5^\beta$
\[
1,\;5,\;25,\;-3,\;-15,\;-11,\;9,\;-19;
\]
also sind die $c_\beta$:
\[
+1,+1,+1,-1,-1,-1,+1,-1
\]
und
\[
\psi(\Theta)=-\Theta^3+\Theta^6-\Theta^7,
\qquad
\psi(\Theta)\psi(-\Theta)=\Theta^6(\Theta^6-2i),
\]
woraus sich leicht ergiebt:
\[
A_6=17.
\]
Eine etwas längere Rechnung ergiebt noch
\[
A_7=21121.
\]

\section*{\S\,199. Der Classenzahlfactor $B$.}

\srcpage{726}
Auf ganz andere Weise muss der Classenzahlfactor $B$ berechnet werden. Hier ist, wenn $\Theta$ wieder eine Primitivwurzel der Gleichung \S~198, (3) ist ($\Theta^{\nu'}+1=0$), und
\begin{equation}
t\equiv(-1)^\alpha5^\beta\pmod m,
\tag{1}
\end{equation}
\begin{equation}
\chi_1(t)=\Theta^\beta
\tag{2}
\end{equation}
zu setzen, und die Formel \S~197, (15) ergiebt:
\begin{equation}
X_1=-\frac2m(1,\Theta^{-1},r)\sum_t\Theta^\beta\log\tan\frac{t\pi}{m}.
\tag{3}
\end{equation}
Nun ist für jedes ungerade $t$, wenn [mit Rücksicht auf (1)]
\[
r=e^{2\pi i/m},
\qquad
r^\nu=i,
\qquad
r^{t\nu}=(-1)^\alpha i
\]
gesetzt wird,
\begin{equation}
\tan\frac{t\pi}{m}=(-1)^\alpha r^t\frac{1-r^t}{1+r^t},
\tag{4}
\end{equation}
und dies ist, wie wir schon im \S~178 gesehen haben, eine Einheit des reellen Körpers $H_m$. Man erhält die conjugirten Werthe dieser Einheit in Bezug auf diesen Körper, wenn man $t\equiv5^\beta\pmod m$ setzt, und $\beta$ ein volles Restsystem nach dem Modul $\mu$ durchlaufen lässt. Wir setzen demnach:
\begin{equation}
\tau_\beta=\tan\frac{5^\beta\pi}{m}.
\tag{5}
\end{equation}
Nach \S~198, (7) ist $5^{\beta+\nu'}\equiv5^\beta+\mu\pmod m$, und daraus ergiebt sich nach einer bekannten trigonometrischen Formel:
\begin{equation}
\tau_{\beta+\nu'}=-\frac1{\tau_\beta}.
\tag{6}
\end{equation}
Setzen wir also nun
\begin{equation}
\lambda_\beta=\frac12\log\tau_\beta^2,
\tag{7}
\end{equation}
so ist nach (6)
\begin{equation}
\lambda_{\beta+\nu'}=-\lambda_\beta,
\tag{8}
\end{equation}
und demnach
\begin{equation}
\Theta^{\beta+\nu'}\lambda_{\beta+\nu'}=\Theta^\beta\lambda_\beta;
\tag{9}
\end{equation}
unter den in der Summe (3) vorkommenden Werthen von $\beta$ sind, wie wir schon im vorigen Paragraphen gesehen haben, nicht zwei nach dem Modul $\nu'$ congruent, und wenn wir daher mittelst
\srcpage{727}
der Formel (9) alle Werthe von $\beta$ auf das Intervall von 0 bis $\nu'-1$ reduciren, so bekommen wir alle Zahlen dieses Intervalles.

Demnach lässt sich die Formel (3) so darstellen:
\begin{equation}
X_1=-\frac2m(1,\Theta^{-1},r)\sum_{\beta=0}^{\nu'-1}\Theta^\beta\lambda_\beta.
\tag{10}
\end{equation}
Dies haben wir in dem im \S~197, (8) angegebenen Ausdrucke
\[
DB=2^{2^{\varkappa-4}(\varkappa-2)}\prod X_1=m^{\nu'}2^{-\nu'}\prod X_1
\]
einzusetzen, und dann ergiebt sich, wenn wir
\begin{equation}
U_\Theta=\sum_{\beta=0}^{\nu'-1}\Theta^\beta\lambda_\beta
\tag{11}
\end{equation}
setzen, mit Rücksicht auf \S~197, (14):
\begin{equation}
DB=\prod_\Theta U_\Theta.
\tag{12}
\end{equation}
Dies Product lässt sich nun auf folgende Weise durch eine Determinante darstellen, die nur die conjugirten Logarithmen der Einheiten $\tau$ enthält.

Lassen wir den Index $\Theta$ bei $U$ weg, so können wir das Gleichungssystem ansetzen:
\[
\begin{aligned}
U&=\lambda_0+\lambda_1\Theta+\lambda_2\Theta^2+\cdots+\lambda_{\nu'-1}\Theta^{\nu'-1},\\
\Theta U&=-\lambda_{\nu'-1}+\lambda_0\Theta+\lambda_1\Theta^2+\cdots+\lambda_{\nu'-2}\Theta^{\nu'-1},\\
&\hspace{2.5cm}\cdots\\
\Theta^{\nu'-1}U&=-\lambda_1-\lambda_2\Theta-\lambda_3\Theta^2+\cdots+\lambda_0\Theta^{\nu'-1},
\end{aligned}
\]
woraus durch Elimination der Potenzen von $\Theta$ folgt:
\begin{equation}
\begin{vmatrix}
\lambda_0-U&\lambda_1&\lambda_2&\cdots&\lambda_{\nu'-1}\\
-\lambda_{\nu'-1}&\lambda_0-U&\lambda_1&\cdots&\lambda_{\nu'-2}\\
\cdots&\cdots&\cdots&\cdots&\cdots\\
-\lambda_1&-\lambda_2&-\lambda_3&\cdots&\lambda_0-U
\end{vmatrix}=0.
\tag{13}
\end{equation}
Dies ist für $U$ eine Gleichung $\nu'$ten Grades, deren Wurzeln die Grössen $U_\Theta$ sind. Ordnet man nach Potenzen von $U$, so erhält, da $\nu'$ gerade ist, die höchste Potenz $U^{\nu'}$ von $U$ den Coefficienten 1, und man findet also das in (12) vorkommende Product der Wurzeln, wenn man $U=0$ setzt, aus der Determinante (13).

Die so gewonnene Determinante wollen wir etwas anders anordnen, indem wir die erste Zeile an ihrer Stelle lassen, die
\srcpage{728}
übrigen aber in umgekehrter Ordnung und mit veränderten Vorzeichen schreiben. Dabei sind $\frac12\nu'$ Vorzeichenänderungen nöthig, und wir erhalten daher für das Product (12) die durch folgende Gleichung definirte Determinante, deren absoluten Werth wir mit $T$ bezeichnen:
\begin{equation}
(-1)^{\frac12\nu'}
\begin{vmatrix}
\lambda_0&\lambda_1&\cdots&\lambda_{\nu'-2}&\lambda_{\nu'-1}\\
\lambda_1&\lambda_2&\cdots&\lambda_{\nu'-1}&-\lambda_0\\
\lambda_2&\lambda_3&\cdots&-\lambda_0&-\lambda_1\\
\cdots&\cdots&\cdots&\cdots&\cdots\\
\lambda_{\nu'-1}&-\lambda_0&\cdots&-\lambda_{\nu'-3}&-\lambda_{\nu'-2}
\end{vmatrix}=\pm T,
\tag{14}
\end{equation}
worin rechts von der von links unten nach rechts oben laufenden Diagonalreihe die negativen Zeichen stehen.

Die Determinante $T$ ist hier so geordnet, dass jede Zeile aus der vorangehenden durch die Substitution $(r,r^{5^\beta})$ entsteht. Jede Colonne enthält also ein System conjugirter Logarithmen.

Nach (12) ist jetzt
\begin{equation}
B=\frac{T}{D}.
\tag{15}
\end{equation}
Um über die Natur dieses Resultates Klarheit zu bekommen, ist es nöthig, die Zahl $D$ etwas genauer zu betrachten, die durch \S~195, (12), \S~197, (6) definirt war; und dies erfordert ein Eingehen auf die Theorie der Einheiten des Körpers $H_m$.




\section*{\S\,200. Normaleinheiten in $H_m$.}

\srcpage{728}
Wir führen jetzt eine vereinfachende Bezeichnung ein, indem wir
\begin{equation}
r^{5^j}=r_j
\tag{1}
\end{equation}
setzen, und $\beta$ ein volles Restsystem nach dem Modul $\nu$ durchlaufen lassen. Es folgt hieraus
\begin{equation}
r_{\beta+\nu'}=-r_\beta,
\tag{2}
\end{equation}
und in der Form $(r,r_j)$ sind dann zwar nicht alle Substitutionen der Gruppe des Körpers $\Omega_m$ enthalten, sondern es müssen noch die Substitutionen $(r,r_j^{-1})$ hinzukommen; wohl aber liefert uns $(r,r_j)$ alle Substitutionen der Gruppe von $H_m$.

Wir bezeichnen jetzt mit $\mathfrak E(r)$ eine Einheit des Körpers $H_m$. Unter diesen sind auch die Einheiten des Körpers $H_\mu$ (als
\srcpage{729}
imprimitive Zahlen) enthalten und durch die Gleichung $\mathfrak E(r)=\mathfrak E(-r)$ charakterisirt.

Wir wollen nun unter einer \textbf{Normaleinheit} des Körpers $H_m$ eine Einheit verstehen, die nicht $=\pm1$ ist, aber der Gleichung
\begin{equation}
\mathfrak E(r)\mathfrak E(-r)=\pm1
\tag{3}
\end{equation}
genügt.

Es ist klar, dass eine Einheit des Körpers $H_\mu$ dieser Bedingung jedenfalls nicht genügen kann. Aber nicht jede Einheit in $H_m$, die dem Körper $H_\mu$ nicht angehört, wird Normaleinheit sein\footnote{Der Verfasser hat in der oben erwähnten Arbeit (Acta mathematica, Bd. 8) die Normaleinheiten als „primitive Einheiten“ bezeichnet. Dieser Name ist hier verworfen, weil er zu dem Irrthum Veranlassung geben könnte, als ob jede Einheit des Körpers $H_m$, die nicht zugleich dem Körper $H_\mu$ angehört, dazu zu rechnen sei.}.

Die Einheiten $\tau_j$ [\S~199, (5)] gehören aber zu den Normaleinheiten. Es ist nämlich
\begin{equation}
\tau_j=\tg\frac{5^j\pi}{m}=\frac{1-r_j}{1+r_j}r_j,
\tag{4}
\end{equation}
und wenn darin die Substitution $(r,-r)=(r_\beta,r_{\beta+\nu'})$ gemacht wird, so geht $\tau_j$ in
\begin{equation}
\tau_{\beta+\nu'}=-\frac1{\tau_\beta}
\tag{5}
\end{equation}
über, was der Relation (3) entspricht.

Ein System von $\nu'$ Normaleinheiten
\[
\mathfrak E_0(r),\ \mathfrak E_1(r),\ldots,\mathfrak E_{\nu'-1}(r)
\]
soll ein unabhängiges System heissen, wenn die Determinante
\begin{equation}
\pm\sum\pm\log|\mathfrak E_0(r_0)|\log|\mathfrak E_1(r_1)|\cdots\log|\mathfrak E_{\nu'-1}(r_{\nu'-1})|
=L(\mathfrak E_0,\mathfrak E_1,\ldots,\mathfrak E_{\nu'-1})
\tag{6}
\end{equation}
einen von Null verschiedenen Werth hat.

Den absoluten Werth dieser Determinante, $L$, nennen wir auch hier den \textbf{Regulator} des Systems $\mathfrak E_0,\mathfrak E_1,\ldots,\mathfrak E_{\nu'-1}$.

Die Einheiten $\tau_0,\tau_1,\ldots,\tau_{\nu'-1}$ bilden ein unabhängiges System normaler Einheiten, denn ihr Regulator
\[
L(\tau_0,\tau_1,\ldots,\tau_{\nu'-1})
\]
ist eben die im vorigen Paragraphen definirte Determinante $T$ [\S~199, (14), (15)]. Und dass diese nicht Null sein kann, ergiebt sich unmittelbar daraus, dass sie in dem Ausdrucke für die
\srcpage{730}
Classenzahl als Factor vorkommt. Die Classenzahl ist aber, da gewiss immer wenigstens eine Classe, die Hauptclasse, existirt, von Null verschieden, und also kann auch $T$ nicht gleich Null sein. Wir haben daher den Satz:

\begin{center}\textbf{1. Die Zahlen $\tau_0,\tau_1,\ldots,\tau_{\nu'-1}$ sind ein unabhängiges System von Normaleinheiten des Körpers $H_m$.}\end{center}

Ist $\mathfrak E(r)$ eine Normaleinheit, so ist, wenn wir
\begin{equation}
l_\beta=\log|\mathfrak E(r_\beta)|
\tag{7}
\end{equation}
setzen, wegen (3)
\begin{equation}
l_{\beta+\nu'}=-l_\beta.
\tag{8}
\end{equation}

Ist nun $\mathfrak E_0,\mathfrak E_1,\ldots,\mathfrak E_{\nu'-1}$ irgend ein System unabhängiger Normaleinheiten, so setzen wir
\begin{equation}
l_{i\beta}=\log|\mathfrak E_i(r_\beta)|,\qquad l_{i,\beta+\nu'}=-l_{i\beta},
\tag{9}
\end{equation}
so dass der Regulator des Systems der $\mathfrak E_i$ den Ausdruck erhält:
\begin{equation}
\pm\sum\pm l_{00}l_{11}\cdots l_{\nu'-1,\nu'-1}=L(\mathfrak E_0,\mathfrak E_1,\ldots,\mathfrak E_{\nu'-1})
\tag{10}
\end{equation}
und eine von Null verschiedene Zahl ist, die wir mit $T_\mathfrak E$ bezeichnen.

Bedeutet $\mathfrak E(r)$ irgend eine andere Normaleinheit, so können wir die Unbekannten $\xi_0,\xi_1,\ldots,\xi_{\nu'-1}$ aus den $\nu'$ linearen Gleichungen bestimmen, die wir erhalten, wenn wir in
\begin{equation}
l_\beta=\xi_0l_{0\beta}+\xi_1l_{1\beta}+\cdots+\xi_{\nu'-1}l_{\nu'-1,\beta}
\tag{11}
\end{equation}
den Index $\beta$ von $0$ bis $\nu'-1$ gehen lassen, und wegen (8) und (9) ist diese Gleichung auch für die übrigen Werthe von $\beta$ richtig, so dass man daraus die für alle $r$ gültige Formel ableiten kann:
\begin{equation}
\mathfrak E(r)=\pm\mathfrak E_0(r)^{\xi_0}\mathfrak E_1(r)^{\xi_1}\cdots\mathfrak E_{\nu'-1}(r)^{\xi_{\nu'-1}}.
\tag{12}
\end{equation}

Bedeuten andererseits $m_0,m_1,\\ldots,m_{\nu'-1}$ irgend welche ganze rationale Zahlen, so sind alle in der Form
\[
\pm\mathfrak E_0(r)^{m_0}\mathfrak E_1(r)^{m_1}\cdots\mathfrak E_{\nu'-1}(r)^{m_{\nu'-1}}
\]
enthaltenen Zahlen Normaleinheiten des Körpers $H_m$. Demnach lassen sich auf das Gleichungssystem (10) die Betrachtungen anwenden, die wir im \S~186 ausführlich durchgeführt haben, woraus sich ergiebt:

\begin{center}\textbf{2. Die aus (10) oder (11) bestimmten Exponenten $\xi_0,\xi_1,\ldots,\xi_{\nu'-1}$ sind rationale Zahlen.}\end{center}

\srcpage{731}
Hieraus kann, wie im \S~187, gefolgert werden, dass es Systeme von $\nu'$ unabhängigen Normaleinheiten giebt, deren Regulator $T_0$ so klein als möglich wird, und solche Systeme heissen \textbf{Fundamentalsysteme von Normaleinheiten}.

Ist $\mathfrak E_0,\mathfrak E_1,\ldots,\mathfrak E_{\nu'-1}$ ein solches Fundamentalsystem, so ergeben sich die Exponenten $\xi_0,\xi_1,\ldots,\xi_{\nu'-1}$ in (10) und (11) als \textbf{ganze rationale Zahlen}. Auch dies kann ebenso bewiesen werden, wie im \S~187.

Stellen wir unter dieser Voraussetzung über die $\mathfrak E_i$ die Gleichungen (11) für irgend ein anderes System unabhängiger Einheiten $\mathfrak E_0^{(1)},\mathfrak E_1^{(1)},\ldots,\mathfrak E_{\nu'-1}^{(1)}$ auf, so erhalten wir Ausdrücke von der Form:
\[
L_{i,k}^{(1)}=\xi_{i,0}L_{0,k}+\xi_{i,1}L_{1,k}+\cdots+\xi_{i,\nu'-1}L_{\nu'-1,k}.
\]
Ist also $T_1$ der Regulator des Systems der $\mathfrak E_i^{(1)}$, so ergiebt sich nach dem Multiplicationssatze der Determinanten, wenn mit $C$ der absolute Werth der Determinante der ganzzahligen Exponenten $\xi_{i,k}$, also eine ganze rationale positive Zahl bezeichnet wird:
\begin{equation}
T_1=CT_0.
\tag{13}
\end{equation}

Es kann nicht bewiesen werden und trifft für die höheren Werthe von $m$ wahrscheinlich auch nicht zu, dass $\tau_0,\tau_1,\ldots,\tau_{\nu'-1}$ ein Fundamentalsystem bilden\footnote{Es ist hier, wie bei allen Fragen, die die Einheiten betreffen, sehr schwierig, zu bestimmten numerischen Resultaten zu gelangen. Nach den Resultaten, die in dem reichhaltigen Werke von Reuschle, „Tafeln complexer Primzahlen, welche aus Wurzeln der Einheiten gebildet sind“ (Berlin 1875), enthalten sind, bilden für $m=16$ die $\tau$ noch ein Fundamentalsystem. Denn hier ist die Grundzahl des Körpers $H_m$ gleich 21, der Grad des Körpers gleich 4, und daher giebt es nach der Anmerkung zu \S~156 in jeder Idealclasse ein Ideal, dessen Norm kleiner als 6 ist. Nach den Tafeln von Reuschle kann aber jede natürliche Primzahl unter 100 im Körper $H_m$ in wirklich existirende Primfactoren zerlegt werden. Folglich gehen gewiss in allen Zahlen unter 6 nur Hauptideale auf, und folglich giebt es hier nur die eine Classe, die Hauptclasse. Da, wie wir gesehen haben, auch der erste Classenzahlfactor $=1$ ist, so ist der Kreistheilungskörper $\Omega_{16}$ einclassig. Es gelten in ihm dieselben Gesetze der Primzahlen, wie im Körper der rationalen Zahlen. Zweifelhaft ist dies für den Körper $\Omega_{32}$ und sicher nicht mehr zutreffend im Körper $\Omega_{64}$, in dem die Classenzahl ein Vielfaches von 17 ist.}.

Für unseren Zweck ist es aber von Wichtigkeit, dass sich diese Einheiten in Bezug auf den Nenner 2 in den Exponenten
\srcpage{732}
gewissermaassen wie ein Fundamentalsystem verhalten, was durch folgenden Satz bestimmter ausgedrückt wird:

\begin{center}\textbf{2. Wenn eine Normaleinheit des Körpers $H_m$ von der Form
\[
\mathfrak E(r)=\tau_0^{e_0}\tau_1^{e_1}\cdots\tau_{\nu'-1}^{e_{\nu'-1}},
\]
in der die Exponenten $e_0,e_1,\ldots,e_{\nu'-1}$ ganze rationale Zahlen sind, mit allen ihren conjugirten Werthen einerlei Vorzeichen hat, so müssen die sämmtlichen $e_0,e_1,\ldots,e_{\nu'-1}$ gerade Zahlen sein, und $\mathfrak E(r)$ ist das Quadrat einer Normaleinheit.}\end{center}

Den Beweis können wir wieder durch vollständige Induction führen, wollen aber zunächst den Ausdruck für $\tau_j$ etwas umformen.

Die conjugirten Werthe zu der Einheit $\mathfrak E(r)$ erhalten wir in der Form
\[
\mathfrak E(r_\beta)=\tau_\beta^{e_0}\tau_{\beta+1}^{e_1}\cdots\tau_{\beta+\nu'-1}^{e_{\nu'-1}},
\]
und diese Ausdrücke sollen also für alle Werthe von $\beta$ dasselbe Vorzeichen haben.

Es ist nach \S~199, (5)
\[
\tau_j=\tg\frac{5^j\pi}{m}=-\frac12\frac{\sin\left(\frac{2\pi}{m}5^j\right)}{\left[\cos\left(\frac{5^j\pi}{m}\right)\right]^2}.
\]
Hierin setzen wir zur Abkürzung
\begin{equation}
\sigma_j=\sin\left(\frac{2\pi}{m}5^j\right),
\tag{14}
\end{equation}
so dass $\sigma_j$ immer dasselbe Vorzeichen hat, wie $\tau_j$.

Daraus bilden wir das Product
\begin{equation}
S_\beta=\sigma_\beta^{e_0}\sigma_{\beta+1}^{e_1}\cdots\sigma_{\beta+\nu'-1}^{e_{\nu'-1}},
\tag{15}
\end{equation}
und unsere Voraussetzung ist also die, dass $S_\beta$ für alle Werthe von $\beta$ dasselbe Vorzeichen hat.

Es soll nachgewiesen werden, dass die ganzen Zahlen $e_0,e_1,\ldots,e_{\nu'-1}$ alle gerade sein müssen.

Der Satz ist richtig für $m=8$; denn in diesem Falle ist
\[
\sigma_0=\sin\frac{2\pi}{8},\qquad \sigma_1=\sin\frac{10\pi}{8}=-\sin\frac{6\pi}{8},
\]
\srcpage{733}
also $\sigma_0$ positiv, $\sigma_1$ negativ. Hier ist nun $S_\beta=\sigma_\beta^{e_0}$, und es kann also $S_0$ und $S_1$ nur dann gleiches Vorzeichen haben, wenn $e_0$ gerade ist.

Hiernach wenden wir die vollständige Induction an. Wir setzen zur Vereinfachung $\nu'=2\nu''$ und erinnern an die Bezeichnung
\begin{equation}
m=2^\kappa,\qquad \mu=\frac12m,\qquad \nu=\frac12\mu,
\qquad \nu'=\frac12\nu,
\qquad \nu''=\frac12\nu',
\tag{16}
\end{equation}
und setzen voraus, dass unser Satz als richtig erwiesen sei, wenn $\mu$ an die Stelle von $m$ tritt.

Es ist dann
\[
5^\nu\equiv1\pmod m,\\ 5^{\nu'}\equiv1+\mu\pmod m,
\quad 5^{\nu''}\equiv1+\nu\pmod\mu,
\]
\[
5^{\nu''}\equiv1+\nu+\mu\pmod m
\]
und daraus ergiebt sich
\[
\sigma_{\beta+\nu}=-\sigma_\beta,
\]
\[
\sigma_{\beta+\nu''}=-\cos\left(\frac{2\pi}{m}5^\beta\right),
\]
\begin{equation}
\sigma_\beta\sigma_{\beta+\nu''}=-\frac12\sin\left(\frac{2\pi}{\mu}5^\beta\right)=-\frac12\sigma'_\beta,
\tag{17}
\end{equation}
\begin{equation}
\sigma'_{\beta+\nu''}=-\sigma'_\beta,
\tag{18}
\end{equation}
wenn die $\sigma'_\beta$ aus den $\sigma_\beta$ durch den Uebergang von $\kappa$ zu $\kappa-1$ hervorgehen. Nun bilden wir nach (17), (18) das Product $S_\beta S_{\beta+\nu''}$, und erhalten, wenn noch zur Abkürzung
\[
e'=e_0+e_1+\cdots+e_{\nu'-1},
\qquad e''=e_0+e_1+\cdots+e_{\nu''-1}
\]
gesetzt wird,
\[
S_\beta S_{\beta+\nu''}=(-1)^{\nu''}\left(\frac12\right)^{e'}
(\sigma'_\beta)^{e_0+e_{\nu''}}(\sigma'_{\beta+1})^{e_1+e_{\nu''+1}}\cdots
(\sigma'_{\beta+\nu''-1})^{e_{\nu''-1}+e_{\nu'-1}}.
\]

Dies ist aber ein Product von derselben Form wie $S_\beta$, in dem $\mu$ an Stelle von $m$ getreten ist, was gleichfalls mit allen seinen Conjugirten einerlei Vorzeichen hat. Es ist daher nach der gemachten Voraussetzung
\begin{equation}
e_0\equiv e_{\nu''},\ e_1\equiv e_{\nu''+1},\ldots,e_{\nu''-1}\equiv e_{\nu'-1}\pmod2.
\tag{19}
\end{equation}
Setzen wir also
\[
S'_\beta=(\sigma'_\beta)^{e_0}(\sigma'_{\beta+1})^{e_1}\cdots(\sigma'_{\beta+\nu''-1})^{e_{\nu''-1}},
\]
\[
R_\beta=\sigma_{\beta+\nu''}^{\frac12(e_0-e_{\nu''})}
\sigma_{\beta+\nu''+1}^{\frac12(e_1-e_{\nu''+1})}\cdots
\sigma_{\beta+\nu'-1}^{\frac12(e_{\nu''-1}-e_{\nu'-1})},
\]
so folgt aus (15) und (17)
\[
S_\beta R_\beta=\left(-\frac12\right)^{e''}S'_\beta,
\]
\srcpage{734}
und hieraus ergiebt sich, dass $S'_\beta$ ebenso wie $S_\beta$ mit allen seinen conjugirten Werthen dasselbe Zeichen hat. $S'_\beta$ entsteht aber aus $S_\beta$ dadurch, dass $\mu$ an Stelle von $m$ gesetzt wird, und nach unserer Voraussetzung ist daher
\[
e_0\equiv0,\ e_1\equiv0,\ldots,e_{\nu''-1}\equiv0\pmod2,
\]
und mit Hülfe von (19)
\[
e_{\nu''}\equiv0,\ e_{\nu''+1}\equiv0,\ldots,e_{\nu'-1}\equiv0\pmod2,
\]
wodurch der Beweis des Satzes 2. geführt ist.

Wenn wir nun in der Formel (12) für $\mathfrak E_0,\mathfrak E_1,\ldots,\mathfrak E_{\nu'-1}$ das System $\tau_0,\tau_1,\ldots,\tau_{\nu'-1}$ wählen und die Exponenten $\xi$ in die Form setzen:
\[
\xi_0=\frac{e_0}{e},\quad \xi_1=\frac{e_1}{e},\ldots,\quad \xi_{\nu'-1}=\frac{e_{\nu'-1}}{e},
\]
worin $e,e_0,e_1,\ldots,e_{\nu'-1}$ ganze rationale Zahlen ohne gemeinsamen Theiler sind, so erhalten wir
\[
\mathfrak E(r)^e=\pm\tau_0^{e_0}\tau_1^{e_1}\cdots\tau_{\nu'-1}^{e_{\nu'-1}},
\]
und es folgt aus unserem Satze, dass $e$ ungerade sein muss; denn wäre es gerade, so wäre $\mathfrak E(r)^e$ ein Quadrat einer reellen Grösse, und folglich hätte $\pm\mathfrak E(r)^e$ mit seinen Conjugirten dasselbe Vorzeichen. Dann aber müssten nach unserem Satze alle $e_0,e_1,\ldots,e_{\nu'-1}$ durch 2 theilbar sein, was wieder gegen die Voraussetzung ist, dass die Exponenten $e,e_0,\ldots,e_{\nu'-1}$ keinen gemeinsamen Theiler haben sollen. Wir haben daher den Satz:

\begin{center}\textbf{3. Ist $\mathfrak E(r)$ irgend eine Normaleinheit des Körpers $H_m$, so lassen sich die ungerade Zahl $e$ und die ganzen Zahlen $e_0,e_1,\ldots,e_{\nu'-1}$ so bestimmen, dass}
\[
\mathfrak E(r_\beta)^e=\pm\tau_\beta^{e_0}\tau_{\beta+1}^{e_1}\cdots\tau_{\beta+\nu'-1}^{e_{\nu'-1}}
\tag{20}
\]
\textbf{wird.}\end{center}

Dies Ergebniss führt zu einigen wichtigen Consequenzen:

\begin{center}\textbf{4. Ist $T$ der Regulator des Systems der $\tau_j$, und $T_0$ der Regulator eines Fundamentalsystems normaler Einheiten in $H_m$, d. h. der Minimalwerth, den der Regulator irgend eines Systems unabhängiger Normaleinheiten annehmen kann, so ist}
\[
T=CT_0,\\ C\equiv1\pmod2,
\tag{21}
\]
\textbf{worin $C$ eine ungerade natürliche Zahl bedeutet.}\end{center}

\srcpage{735}
Dass nämlich $T:T_0=C$ eine ganze Zahl ist, folgt aus der Formel (13).

Wenn wir aber den Satz 3. auf alle Elemente des Fundamentalsystems $\mathfrak E_i(r)$ anwenden, so ergiebt sich ein System ganzer Zahlen $e_{i,k}$ und eine ungerade Zahl $e$, so dass
\[
\mathfrak E_i(r)^e=\pm\tau_0^{e_{i,0}}\tau_1^{e_{i,1}}\cdots\tau_{\nu'-1}^{e_{i,\nu'-1}}.
\]
Nehmen wir hiervon die Logarithmen, und bilden dann den Regulator $T_0$, so ergiebt sich
\[
e^{\nu'}T_0=\pm\sum\pm e_{0,0}e_{1,1}\cdots e_{\nu'-1,\nu'-1}T,
\]
folglich
\[
e^{\nu'}=\pm C\sum\pm e_{0,0}e_{1,1}\cdots e_{\nu'-1,\nu'-1},
\]
und da hierin die linke Seite ungerade ist, so kann $C$ nicht gerade sein, w. z. b. w.

Wenn alle conjugirten Werthe der Normaleinheit $\mathfrak E(r_\beta)$ einerlei Zeichen haben, so gilt dasselbe auch von $\mathfrak E(r_\beta)^e$, und folglich müssen in diesem Falle in der Formel (20) die Exponenten $e_0,e_1,\ldots,e_{\nu'-1}$ nach 2. gerade Zahlen sein. Schreibt man dann die Formel (20) so:
\[
\mathfrak E(r)=\pm\tau_0^{e_0}\tau_1^{e_1}\cdots\tau_{\nu'-1}^{e_{\nu'-1}}\mathfrak E(r)^{1-e},
\]
so erhält man daraus den Satz:

\begin{center}\textbf{5. Eine Normaleinheit des Körpers $H_m$, die mit ihren Conjugirten gleiches Vorzeichen hat, ist, vom Vorzeichen abgesehen, ein Quadrat einer Normaleinheit.}\end{center}

\section*{\S\,201. Fundamentalsystem von Einheiten des Körpers $H_m$.}

Es bleibt noch übrig, den Nenner $D$ in der Formel für $B$ [\S~197, (6)] zu bestimmen, der durch die Gleichung
\[
E=E'D
\]
definirt war, worin $E$ und $E'$ die Regulatoren der Körper $H_m,H_\mu$ bedeuteten.

Hier handelt es sich also nicht mehr um die Normaleinheiten, sondern um die Einheiten des Körpers $H_m$ überhaupt.

\srcpage{736}
Wir nehmen ein Fundamentalsystem von Einheiten im Körper $H_m$ an:
\begin{equation}
\varepsilon_1(r),\ \varepsilon_2(r),\ldots,\varepsilon_{\nu-1}(r),
\tag{1}
\end{equation}
und die conjugirten Logarithmen dieses Systems seien:
\begin{equation}
l_{1,k},\ l_{2,k},\ldots,l_{\nu-1,k},\qquad k=0,1,2,\ldots,\nu-1.
\tag{2}
\end{equation}
Daneben betrachten wir ein Fundamentalsystem im Körper $H_\mu$
\begin{equation}
\varepsilon'_1,\ \varepsilon'_2,\ldots,\varepsilon'_{\nu'-1}
\tag{3}
\end{equation}
mit den conjugirten Logarithmen
\begin{equation}
l'_{1,k},\ l'_{2,k},\ldots,l'_{\nu'-1,k},\qquad k=0,1,2,\ldots,\nu'-1.
\tag{4}
\end{equation}

Bei der Bildung der Regulatoren wird einer der conjugirten Werthe, etwa $k=0$, weggelassen (vgl. \S~186), und wir erhalten daher
\begin{equation}
E=\pm\sum\pm l_{1,1}l_{2,2}\cdots l_{\nu-1,\nu-1},
\qquad
E'=\pm\sum\pm l'_{1,1}l'_{2,2}\cdots l'_{\nu'-1,\nu'-1}.
\tag{5}
\end{equation}

Nun fügen wir zu dem Systeme (3) ein Fundamentalsystem von Normaleinheiten des Körpers $H_m$
\begin{equation}
\mathfrak E_0,\mathfrak E_1,\ldots,\mathfrak E_{\nu'-1},
\tag{6}
\end{equation}
und wir behaupten, dass das System
\begin{equation}
\varepsilon'_1,\varepsilon'_2,
\ldots,\varepsilon'_{\nu'-1},\mathfrak E_0,\mathfrak E_1,
\ldots,\mathfrak E_{\nu'-1}
\tag{7}
\end{equation}
ein unabhängiges System von Einheiten des Körpers $H_m$ ist (wenn auch nicht ein Fundamentalsystem).

Bezeichnen wir mit
\begin{equation}
L_{0,k},L_{1,k},\ldots,L_{\nu'-1,k},\qquad k=0,1,\ldots,\nu'-1
\tag{8}
\end{equation}
die conjugirten Logarithmen des Systems (6), so ist nach dem vorigen Paragraphen
\begin{equation}
\pm\sum\pm L_{0,0}L_{1,1}\cdots L_{\nu'-1,\nu'-1}=T_0
\tag{9}
\end{equation}
der Regulator dieses Systems, und wir haben mit Rücksicht auf die Definition der Normaleinheiten [\S~200, (9)]:
\begin{equation}
L_{i,k+\nu'}=-L_{i,k},
\tag{10}
\end{equation}
und weil die $\varepsilon'_i$ als Zahlen des Körpers $H_\mu$ durch die Substitution $(r,-r)$ ungeändert bleiben:
\begin{equation}
l'_{i,k+\nu'}=l'_{i,k}.
\tag{11}
\end{equation}

\srcpage{737}
Hiernach ergiebt sich der Regulator $R$ des Systems (7) aus der Determinante:
\[
\small
\left|\begin{array}{cccc|cccc}
l'_{1,1}&l'_{2,1}&\cdots&l'_{\nu'-1,1}&L_{0,1}&L_{1,1}&\cdots&L_{\nu'-1,1}\\
\cdots&\cdots&&\cdots&\cdots&\cdots&&\cdots\\
l'_{1,\nu'-1}&l'_{2,\nu'-1}&\cdots&l'_{\nu'-1,\nu'-1}&L_{0,\nu'-1}&L_{1,\nu'-1}&\cdots&L_{\nu'-1,\nu'-1}\\
l'_{1,0}&l'_{2,0}&\cdots&l'_{\nu'-1,0}&-L_{0,0}&-L_{1,0}&\cdots&-L_{\nu'-1,0}\\
l'_{1,1}&l'_{2,1}&\cdots&l'_{\nu'-1,1}&-L_{0,1}&-L_{1,1}&\cdots&-L_{\nu'-1,1}\\
\cdots&\cdots&&\cdots&\cdots&\cdots&&\cdots\\
l'_{1,\nu'-1}&l'_{2,\nu'-1}&\cdots&l'_{\nu'-1,\nu'-1}&-L_{0,\nu'-1}&-L_{1,\nu'-1}&\cdots&-L_{\nu'-1,\nu'-1}
\end{array}\right|.
\]
Diese Determinante hat $\nu-1$ Zeilen; wir addiren die erste zur $(\nu'+1)$ten, die zweite zur $(\nu'+2)$ten u.s.f., die $(\nu'-1)$te zur letzten. Dadurch heben sich in den letzten $(\nu'-1)$ Zeilen die $L$ heraus und die $l'$ bekommen den Factor 2 [nach (10), (11)], und hieraus ergiebt sich nach einem Determinantensatze (Bd. I, \S~23, V.):
\begin{equation}
R=2^{\nu'-1}E'T_0.
\tag{12}
\end{equation}

Da hiernach $R$ nicht verschwindet und daher in (7) ein System von $\nu-1$ unabhängigen Einheiten des Körpers $H_m$ vorliegt, so können wir [nach \S~186, (5)] die rationalen (ganzen oder gebrochenen) Zahlen $m_{i,k},M_{i,k}$ aus den Gleichungen
\begin{equation}
2l_{i,k}=m_{1,i}l'_{1,k}+\cdots+m_{\nu'-1,i}l'_{\nu'-1,k}
+M_{0,i}L_{0,k}+M_{1,i}L_{1,k}+\cdots+M_{\nu'-1,i}L_{\nu'-1,k}
\tag{13}
\end{equation}
bestimmen. Es lässt sich aber leicht durch die Relationen (10), (11) nachweisen, dass die Zahlen $m_{i,k},M_{i,k}$ hier \textbf{ganze} sein müssen. Denn es ist
\[
l_{i,k}+l_{i,k+\nu'}=\log|\varepsilon_i(r_k)\varepsilon_i(-r_k)|
=m_{1,i}l'_{1,k}+\cdots+m_{\nu'-1,i}l'_{\nu'-1,k},
\]
\[
l_{i,k}-l_{i,k+\nu'}=\log|\varepsilon_i(r_k)\varepsilon_i(-r_k)^{-1}|
=M_{0,i}L_{0,k}+M_{1,i}L_{1,k}+\cdots+M_{\nu'-1,i}L_{\nu'-1,k}.
\]
Die Einheiten $\varepsilon_i(r)\varepsilon_i(-r)$ gehören aber dem Körper $H_\mu$ an, und da die $\varepsilon'_i$ nach Voraussetzung ein Fundamentalsystem dieses Körpers sind, so müssen die $m_{i,k}$ nach \S~187 ganze rationale Zahlen sein.

Die Einheiten $\varepsilon_i(r)\varepsilon_i(-r)^{-1}$ sind Normaleinheiten des Körpers $H_m$, und weil die $\mathfrak E_i$ ein Fundamentalsystem von Normaleinheiten sein sollten, so sind nach \S~200 die $M_{i,k}$ ganze Zahlen.

\srcpage{738}
Zur besseren Uebersicht setzen wir die Gleichungen (13) noch einmal ohne den Index $k$, der die conjugirten Werthe von einander unterscheidet, hierher:
\begin{equation}
2l_i=m_{1,i}l'_1+m_{2,i}l'_2+\cdots+m_{\nu'-1,i}l'_{\nu'-1}
+M_{0,i}L_0+M_{1,i}L_1+\cdots+M_{\nu'-1,i}L_{\nu'-1}.
\tag{14}
\end{equation}
Bilden wir nun nach (13) den Regulator $E$ und bezeichnen mit $M$ den absoluten Werth der Determinante der $\nu-1$ Reihen der Zahlen $m_{k,i},M_{k,i}$:
\[
M=\pm\sum\pm m_{1,1}\cdots m_{\nu'-1,\nu'-1}M_{0,0}M_{1,1}\cdots M_{\nu'-1,\nu'-1},
\]
so folgt:
\begin{equation}
2^{\nu-1}E=MR,
\tag{15}
\end{equation}
und daraus mit Rücksicht auf (12):
\begin{equation}
E=2^{-\nu'}ME'T_0.
\tag{16}
\end{equation}

Es lässt sich nun weiter nachweisen, dass $M$ eine Potenz von 2 und durch $2^{\nu'}$ theilbar ist.

Da nämlich das System (1) als Fundamentalsystem in $H_m$ vorausgesetzt war, so giebt es ganze rationale Zahlen $n_{k,i},N_{k,i}$, die den Gleichungen genügen:
\begin{equation}
\begin{aligned}
l'_i&=n_{1,i}l_1+n_{2,i}l_2+\cdots+n_{\nu-1,i}l_{\nu-1},&& i=1,2,\ldots,\nu'-1,\\
L_i&=N_{1,i}l_1+N_{2,i}l_2+\cdots+N_{\nu-1,i}l_{\nu-1},&& i=0,1,2,\ldots,\nu'-1.
\end{aligned}
\tag{17}
\end{equation}
Bilden wir daraus die Determinante $R$ und bezeichnen mit $N$ den absoluten Werth der Determinante der $n_{k,i},N_{k,i}$, so folgt:
\[
R=NE,
\]
und daraus nach (15):
\[
MN=2^{\nu-1};
\]
daraus folgt, dass jede der beiden ganzen Zahlen $M,N$ eine Potenz von 2 sein muss.

Um zu entscheiden, durch welche Potenz von 2 die Determinante $M$ theilbar ist, bestimmen wir zunächst ein System ganzer rationaler Zahlen $a_i$ ohne gemeinschaftlichen Theiler, das den linearen Gleichungen
\begin{equation}
\sum_{i=1}^{\nu-1}a_im_{k,i}=0\qquad(k=1,2,\ldots,\nu'-1)
\tag{18}
\end{equation}
genügt. Setzen wir dann
\[
\sum_{i=1}^{\nu-1}a_iM_{k,i}=\xi_k,
\]
\srcpage{739}
so erhalten wir aus den Gleichungen (13)
\[
2\sum_{i=1}^{\nu-1}a_il_{i,k}=\xi_0L_{0,k}+\xi_1L_{1,k}+\cdots+\xi_{\nu'-1}L_{\nu'-1,k}.
\]
Daraus ergiebt sich nach (10)
\[
\sum_{i=1}^{\nu-1}a_il_{i,k+\nu'}=-\sum_{i=1}^{\nu-1}a_il_{i,k},
\]
und dies bedeutet, dass
\[
\varepsilon_1^{a_1}\varepsilon_2^{a_2}\cdots\varepsilon_{\nu-1}^{a_{\nu-1}}
\]
eine Normaleinheit des Körpers $H_m$ ist. Daraus folgt dann weiter, weil $\mathfrak E_0,\mathfrak E_1,\ldots,\mathfrak E_{\nu'-1}$ ein Fundamentalsystem von Normaleinheiten ist, dass die Zahlen $\xi_0,\xi_1,\ldots,\xi_{\nu'-1}$ gerade sein müssen. Daraus ergiebt sich:

\begin{center}\textbf{1. Das System der Gleichungen (18) hat die Congruenzen
\[
\sum_{i=1}^{\nu-1}a_iM_{k,i}\equiv0\pmod2\qquad(k=0,1,\ldots,\nu'-1)
\tag{19}
\]
zur Folge.}\end{center}

In den Gleichungen (18), deren Anzahl nur $\nu'-1$ beträgt, sind aber $\nu-1$ Unbekannte $a_i$ enthalten. Daher giebt es mehrere Systeme von einander unabhängiger Lösungen, was wir etwas genauer dahin präcisiren können:

\begin{center}\textbf{2. Es giebt $\nu'$ ganzzahlige Lösungen des Systems (18):}
\[
\begin{array}{cccc}
a_{1,1}&a_{2,1}&\cdots&a_{\nu-1,1}\\
a_{1,2}&a_{2,2}&\cdots&a_{\nu-1,2}\\
\cdots&\cdots&&\cdots\\
a_{1,\nu'}&a_{2,\nu'}&\cdots&a_{\nu-1,\nu'}
\end{array}
\tag{20}
\]
\textbf{von der Art, dass aus der Matrix (20) eine Determinante von $\nu'$ Reihen gebildet werden kann, die eine ungerade Zahl ist.}\end{center}

Denn es giebt zunächst eine Lösung von (18):
\[
a_{1,1},a_{2,1},\ldots,a_{\nu-1,1},
\]
in der eine ungerade Zahl vorkommt. Nehmen wir an, was gestattet ist, da hier auf die Anordnung der Indices nichts ankommt,
\srcpage{740}
es sei $a_{1,1}$ diese ungerade Zahl, dann bestimmen wir eine zweite Lösung:
\[
0,a_{2,2},a_{3,2},\ldots,a_{\nu-1,2},
\]
in der wieder eine ungerade Zahl, etwa $a_{2,2}$, vorkommt. Dann bestimmen wir eine dritte Lösung:
\[
0,0,a_{3,3},\ldots,
\]
und fahren mit Nullsetzen so lange fort, als die Anzahl der übrig bleibenden Unbekannten noch grösser ist, als die Anzahl der Gleichungen.

Im letzten Systeme werden also $\nu'-1$ Nullen vorkommen, damit $\nu'$ Unbekannte übrigbleiben.

Diese Zahlenreihen $a_{i,k}$ können wir für das System (20) setzen, denn darin reducirt sich die Determinante
\begin{equation}
\pm\sum\pm a_{1,1}a_{2,2}\cdots a_{\nu',\nu'}
\tag{21}
\end{equation}
auf das Diagonalglied $a_{1,1}a_{2,2}\cdots a_{\nu',\nu'}$, dessen Factoren alle ungerade sind.

Ist nun die Forderung des Satzes 2. erfüllt, so lässt sich die Matrix (20) durch Hinzufügung von $\nu'-1$ Zeilen zu einer Determinante von $\nu-1$ Reihen ergänzen, die gleichfalls einen ungeraden ganzzahligen Werth hat:
\begin{equation}
\alpha=\pm\sum\pm a_{1,1}a_{2,2}\cdots a_{\nu',\nu'}a_{\nu'+1,\nu'+1}\cdots a_{\nu-1,\nu-1}.
\tag{22}
\end{equation}
Man braucht nur, wenn die Determinante (21) als ungerade vorausgesetzt wird, in den hinzugefügten Zeilen für die Elemente lauter Nullen zu setzen, ausgenommen die in der Diagonalreihe stehenden, $a_{\nu'+1,\nu'+1},\ldots,a_{\nu-1,\nu-1}$, die gleich 1 genommen werden können. Dann erhält $\alpha$ denselben Werth wie die Determinante (21).

Wenn wir aber jetzt nach dem Multiplicationssatze der Determinanten das Product $\alpha M$ bilden, so ergiebt sich eine Determinante, in der die $\nu-1$ Summen
\[
\sum_{i=1}^{\nu-1}a_i m_{k,i},
\qquad
\sum_{i=1}^{\nu-1}a_i M_{k,i},
\]
\[
k=1,2,\ldots,\nu'-1,
\qquad
k=0,1,2,\ldots,\nu'-1
\]
für jeden Werth $s=1,2,\ldots,\nu'$ eine Reihe bilden, und in der also $\nu'$ Reihen vorkommen, deren Elemente alle durch 2 theilbar sind.

\srcpage{741}
Demnach ist $\alpha M$ und folglich auch $M$ durch $2^{\nu'}$ theilbar, und wenn wir
\begin{equation}
M=2^{\nu'+\theta},\qquad N=2^{\nu'-\theta-1}
\tag{23}
\end{equation}
setzen, so ist $\theta$ eine nicht negative ganze rationale Zahl.

Setzen wir dies in (16) ein, so folgt:
\[
E=2^\theta E'T_0.
\]
Nun war in \S~197, (6)
\[
E=E'D
\]
gesetzt, und es ergiebt sich also:
\[
D=2^\theta T_0.
\]
Hieraus folgt für den Classenzahlfactor $B$ nach \S~199, (15) und \S~200, (21):
\begin{equation}
B=2^{-\theta}\frac{T}{T_0}=2^{-\theta}C,
\tag{24}
\end{equation}
worin $C$ eine ungerade ganze Zahl ist.

Demnach ergiebt sich nach \S~197, (4) für den zweiten Factor der Classenzahl $h_1$, der, wie wir gesehen haben, eine ganze Zahl ist:
\begin{equation}
h_1=2^{-\theta}h'_1C.
\tag{25}
\end{equation}
Nehmen wir nun als bewiesen an, dass $h'_1$ ungerade ist, so folgt, da auch $C$ ungerade ist, dass $\theta=0$ und $h_1$ ungerade ist, und beides ist also damit durch vollständige Induction bewiesen.

Damit werden die abgeleiteten Resultate folgendermaassen vereinfacht:

\begin{center}\textbf{3. Es ist}
\[
M=2^{\nu'},\qquad D=T_0,
\qquad B=\frac{T}{T_0};
\tag{26}
\]
\textbf{der Classenzahlfactor $B$ ist gleich dem Verhältniss des Regulators $T$ der Einheiten $\tau$ zum Regulator $T_0$ eines Fundamentalsystems von Normaleinheiten, und ist eine ungerade ganze Zahl.}\end{center}

Bezeichnen wir die zu $m=2^\kappa$ gehörige Zahl $B$ mit $B_\kappa$, so ist die Classenzahl $h_1$, wie durch vollständige Induction folgt,
\begin{equation}
h_1=B_4B_5\cdots B_\kappa,
\tag{27}
\end{equation}
und ist also auch eine ungerade Zahl. Damit ist dann, mit Rücksicht auf \S~198, 1., das Haupttheorem bewiesen:

\begin{center}\textbf{A. Die Classenzahl im Körper $\Omega_m$ ist, wenn $m$ eine Potenz von 2 ist, ungerade.}\end{center}

\section*{\S\,202. Positive Einheiten.}

\srcpage{742}
Die zuletzt gefundenen Resultate zeigen, dass das System unabhängiger Einheiten \S~201, (7)
\[
\varepsilon'_1,\varepsilon'_2,\ldots,\varepsilon'_{\nu'-1},\mathfrak E_0,\mathfrak E_1,\ldots,\mathfrak E_{\nu'-1}
\]
im Körper $H_m$ kein Fundamentalsystem bildet. Denn sonst müssten in den Formeln \S~201, (13) die Zahlen $m_{k,i},M_{k,i}$ sämmtlich durch 2 theilbar sein, und es wäre die Determinante $M$, deren Werth wir $=2^{\nu'}$ gefunden haben, durch $2^{\nu-1}$ theilbar.

Wir betrachten nun an Stelle des Systems der Gleichungen (18), \S~201 die folgenden Congruenzen:
\begin{equation}
\begin{aligned}
\sum_{i=1}^{\nu-1}b_i m_{1,i}&\equiv1\pmod2,\\
\sum_{i=1}^{\nu-1}b_i m_{2,i}&\equiv0\pmod2,\\
&\ \vdots\\
\sum_{i=1}^{\nu-1}b_i m_{\nu'-1,i}&\equiv0\pmod2,
\end{aligned}
\tag{1}
\end{equation}
und beweisen, dass es möglich ist, ihnen durch ganzzahlige Werthe der $b_i$ zu genügen.

Wäre es nicht möglich, so müsste sich aus den $\nu'-2$ Congruenzen
\[
\sum_{i=1}^{\nu-1}b_i m_{k,i}\equiv0\pmod2\qquad k=2,3,\ldots,\nu'-1
\]
die letzte
\[
\sum_{i=1}^{\nu-1}b_i m_{1,i}\equiv0\pmod2
\]
als Folgerung ergeben, und dann würde, wie im vorigen Paragraphen, weiter folgen:
\[
\sum_{i=1}^{\nu-1}b_i M_{k,i}\equiv0\pmod2\qquad k=0,1,\ldots,\nu'-1.
\]

\srcpage{743}
Man könnte dann an Stelle der Matrix \S~201, (20) eine Matrix der $b_{k,i}$ von $\nu'+1$ Reihen setzen, und es würde sich auf dem im vorigen Paragraphen eingeschlagenen Wege schliessen lassen, dass $M$ durch $2^{\nu'+1}$ theilbar sein müsste, was nicht der Fall ist.

Wenden wir aber das System der Multiplicatoren $b_i$, das den Bedingungen (1) genügt, auf die Gleichungen (14) des vorigen Paragraphen an, indem wir mit $b_i$ multipliciren und dann summiren, so ergiebt sich, wenn wir Vielfache von 2 weglassen und dies auch hier (wo irrationale Zahlen, nämlich die Logarithmen, auftreten) durch das Zeichen der Congruenz ausdrücken:
\begin{equation}
l'_1\equiv L_0\sum M_{0,i}b_i+L_1\sum M_{1,i}b_i+\cdots+L_{\nu'-1}\sum M_{\nu'-1,i}b_i\pmod2.
\tag{2}
\end{equation}
Wollen wir die Congruenz in eine Gleichung verwandeln, so kommt eine Summe von Grössen
\[
l_1,l_2,\ldots,l_{\nu-1},\quad l'_1,l'_2,\ldots,l'_{\nu'-1}
\]
mit geraden ganzen Zahlen als Coefficienten hinzu.

Dieselbe Betrachtung lässt sich aber auf $l'_2,l'_3,\ldots,l'_{\nu'-1}$ anwenden, und danach kann man, wenn man von den Logarithmen wieder zu den Zahlen übergeht, die Formel (2) so in Worte fassen:

\begin{center}\textbf{1. Jede Einheit des Körpers $H_\mu$ ist als Product einer Normaleinheit des Körpers $H_m$ und eines Quadrates einer Einheit in $H_m$ darstellbar.}\end{center}

Wenn wir darauf den Satz 2., \S~200 anwenden, so ergiebt sich, dass jede Einheit $\mathfrak E(r)$ des Körpers $H_\mu$, die mit allen ihren Conjugirten dasselbe Zeichen hat, vom Vorzeichen abgesehen, das Quadrat einer Einheit in $H_m$ sein muss. Ist aber
\[
\pm\mathfrak E(r)=\Delta(r)^2,
\]
worin $\mathfrak E(r)$ eine Einheit in $H_\mu$ und $\Delta(r)$ eine Einheit in $H_m$ ist, so muss auch $\Delta(r)$ in $H_\mu$ enthalten sein.

Denn setzen wir
\[
\Delta(r)=\Delta_1(r^2)+r\Delta_2(r^2),
\]
so kann das Quadrat davon nur dann in $H_\mu$ enthalten sein, also durch die Vorzeichenänderung von $r$ ungeändert bleiben, wenn entweder $\Delta_1$ oder $\Delta_2=0$ ist. Es kann aber nicht $\Delta(r)=r\Delta_2(r^2)$ sein, denn sonst wäre $\Delta_2(r^2)$ eine Einheit in $\Omega_\mu$ und müsste
\srcpage{744}
nach \S~178, 1. das Product einer reellen Einheit und einer Einheitswurzel $r^{2\lambda}$ sein, und da auch $\Delta(r)$ reell ist, so müsste $r^{2\lambda+1}$ reell sein, was nicht möglich ist. Demnach ist $\Delta(r)$ in $H_\mu$ enthalten.

Da nun $\mu$ ebenso wie $m$ jede beliebige Potenz von 2 sein kann, so ergiebt sich hieraus das zweite Haupttheorem:

\begin{center}\textbf{B. Eine Einheit des Körpers $H_m$, die mit allen ihren Conjugirten positiv ist, ist das Quadrat einer Einheit desselben Körpers.}\end{center}

Von den beiden Theoremen A. und B. haben wir aber im \S~184 den Beweis des Hauptsatzes von den Abel'schen Zahlkörpern (\S~179, I.) abhängig gemacht.

\end{document}
